\parte{Medidas y conjuntos medibles}{Medir exige decidir primero qué conjuntos pueden medirse. Construiremos las familias medibles, las medidas exteriores y la medida de Lebesgue. El lector necesita operaciones elementales de conjuntos, sucesiones y las nociones de continuidad y compacidad del volumen I. La distinción entre álgebra, $\sigma$-álgebra y anillo explica por qué las extensiones requieren cubiertas numerables.}
\chapter{Álgebras de conjuntos y medidas}\label{cap:v2-01}
Medir exige decidir primero qué conjuntos se admitirán y después asignarles tamaños compatibles con las operaciones numerables. Una familia de conjuntos no debe confundirse con el espacio que la contiene: usaremos $X$ para el espacio y $\Acal$ para sus conjuntos medibles. Este capítulo desarrolla las lecturas 1 y 2 de Valencia Bucio \cite{valencia2}.

\section{Familias estables bajo operaciones}\label{lectura:v2-01}
\begin{definicion}\index{algebra de conjuntos@álgebra de conjuntos}\index{sigmaálgebra@\(\sigma\)-álgebra}
Un \emph{álgebra} sobre $X$ es una familia $\Acal\subseteq\mathcal P(X)$ que contiene $X$ y es estable bajo complementos relativos a $X$ e intersecciones finitas. Una \emph{$\sigma$-álgebra} es un álgebra estable, además, bajo uniones numerables.
\end{definicion}
Toda álgebra contiene $\varnothing=X\setminus X$. Por De Morgan contiene $A\cup B=(A^c\cap B^c)^c$; por inducción contiene uniones e intersecciones finitas. También contiene $A\setminus B=A\cap B^c$ y $A\mathbin\triangle B=(A\setminus B)\cup(B\setminus A)$. En una $\sigma$-álgebra, De Morgan da igualmente intersecciones numerables. Estas identidades explican cada propiedad: no son nuevas condiciones independientes.

\begin{ejemplo}
Las familias $\{\varnothing,X\}$ y $\mathcal P(X)$ son $\sigma$-álgebras. Los conjuntos finitos o cofinitos forman un álgebra: el complemento intercambia ambas clases; la unión de dos finitos es finita y una unión que contiene un cofinito es cofinita. Si $X$ es infinito, esta familia puede no ser una $\sigma$-álgebra. En $\N$, la unión de los singletons pares es infinita y tiene complemento infinito.

Más generalmente, para un cardinal infinito $\kappa$, los conjuntos de cardinal menor que $\kappa$ o con complemento de cardinal menor que $\kappa$ forman un álgebra. La unión de dos conjuntos de cardinal menor que $\kappa$ sigue teniendo cardinal menor que $\kappa$: el máximo de sus dos cardinales, multiplicado por $2$, permanece menor que $\kappa$. Los conjuntos de cardinal a lo sumo $\kappa$ o con complemento de cardinal a lo sumo $\kappa$ forman una $\sigma$-álgebra, pues una unión numerable de conjuntos de cardinal a lo sumo $\kappa$ tiene cardinal a lo sumo $\kappa\aleph_0=\kappa$. Se usa elección en esta aritmética cardinal.
\end{ejemplo}
\begin{observacion}[Una familia que no es un álgebra]
En la lectura 1, ejemplo 3(iii), se afirma que los intervalos y sus complementos forman un álgebra en $\R$. La afirmación es falsa: $(0,1)\cup(2,3)$ no es un intervalo y su complemento tampoco lo es. Una versión correcta es la familia de \emph{uniones finitas de intervalos}, permitiendo intervalos no acotados, singletons y cualquier combinación de extremos abiertos o cerrados. Una partición finita de la recta por todos los extremos involucrados verifica su estabilidad bajo las operaciones del álgebra.
\end{observacion}

\begin{proposicion}\label{prop:v2-generada}
La intersección de una familia no vacía de álgebras sobre $X$ es un álgebra; lo mismo vale para $\sigma$-álgebras. Por tanto, para cualquier $\mathcal G\subseteq\mathcal P(X)$ existen el álgebra $A(\mathcal G)$ y la $\sigma$-álgebra $\sigma(\mathcal G)$ generadas por $\mathcal G$: son las menores familias de la clase correspondiente que contienen $\mathcal G$.
\end{proposicion}
\begin{proof}
El conjunto $X$ pertenece a cada familia y, por ello, a la intersección. Si $A$ pertenece a todas, también $A^c$ pertenece a todas. Una intersección finita o unión numerable de elementos comunes pertenece a cada familia por su estabilidad. Para construir la familia generada intersectamos todas las familias admisibles que contienen $\mathcal G$. Hay al menos una, $\mathcal P(X)$; la intersección contiene $\mathcal G$ y está contenida en cualquier otra candidata.
\end{proof}
\begin{definicion}\index{Borel, sigmaalgebra de@Borel, sigmaálgebra de}
En un espacio topológico $X$, $\Bcal(X)=\sigma(\tau)$ es la \emph{$\sigma$-álgebra de Borel}. Para espacios medibles $(X,\Acal)$ y $(Y,\Bcal)$ se define
\[\Acal\otimes\Bcal=\sigma\{A\times B:A\in\Acal,\ B\in\Bcal\}.\]
La familia de rectángulos generadores generalmente no es un álgebra.
\end{definicion}

\section{Ejercicios sobre las familias medibles}
\begin{ejercicio}[Lectura 1, ejercicio 1]\label{ej:v2-l01-1}
Sea $\mathcal G$ una familia no vacía estable bajo complementos y uniones binarias. Pruebe que es un álgebra.
\end{ejercicio}

\begin{ejercicio}[Lectura 1, ejercicio 2]\label{ej:v2-l01-2}
Dada una sucesión $(A_n)$ en un álgebra, construya $B_n\subseteq A_n$, disjuntos dos a dos, con $\bigcup_nB_n=\bigcup_nA_n$.
\end{ejercicio}

\begin{ejercicio}[Lectura 1, ejercicio 3]\label{ej:v2-l01-3}
Pruebe que $\Acal|_Y=\{A\cap Y:A\in\Acal\}$ es un álgebra, o una $\sigma$-álgebra, sobre $Y\subseteq X$ si $\Acal$ lo es sobre $X$.
\end{ejercicio}

\begin{ejercicio}[Lectura 1, ejercicio 4]\label{ej:v2-l01-4}
Enliste todas las álgebras sobre un conjunto de cuatro elementos.
\end{ejercicio}

\begin{ejercicio}[Lectura 1, ejercicio 5]\label{ej:v2-l01-5}
Caracterice $A(\mathcal G)$ mediante expresiones finitas de intersecciones y complementos de elementos de una familia no vacía $\mathcal G$.
\end{ejercicio}

\begin{ejercicio}[Lectura 1, ejercicio 6]\label{ej:v2-l01-6}
En un espacio métrico separable, pruebe que las bolas abiertas generan $\Bcal(X)$.
\end{ejercicio}

\begin{ejercicio}[Lectura 1, ejercicio 7]\label{ej:v2-l01-7}
Pruebe que los rayos $(a,\infty)$ generan $\Bcal(\R)$.
\end{ejercicio}

\begin{ejercicio}[Lectura 1, ejercicio 8]\label{ej:v2-l01-8}
Si $X$ tiene al menos dos puntos, pruebe que los rectángulos $A\times B$ no forman un álgebra en $X\times X$. Determine el álgebra que generan.
\end{ejercicio}

\begin{ejercicio}[Lectura 1, ejercicio 9]\label{ej:v2-l01-9}
Si $f:X\to Y$ y $\Bcal$ es una $\sigma$-álgebra sobre $Y$, determine si $\{f^{-1}(B):B\in\Bcal\}$ es una $\sigma$-álgebra.
\end{ejercicio}


\section{Aditividad numerable y continuidad}\label{lectura:v2-02}
En $[0,\infty]$ adoptamos $0\cdot\infty=0$ y $a+\infty=\infty$. Nunca efectuaremos la resta $\infty-\infty$. Una serie de términos no negativos es el supremo de sus sumas parciales, incluso cuando vale infinito.
\begin{definicion}\index{medida}\index{sigmafinitud@\(\sigma\)-finitud}
Una \emph{medida} en $(X,\Acal)$ es una función $\mu:\Acal\to[0,\infty]$ tal que $\mu(\varnothing)=0$ y
\[\mu\left(\bigcup_nA_n\right)=\sum_n\mu(A_n)\quad\text{si los }A_n\text{ son disjuntos dos a dos}.\]
La terna $(X,\Acal,\mu)$ es un espacio de medida. Es finito si $\mu(X)<\infty$; es $\sigma$-finito si $X$ admite una partición numerable en conjuntos medibles de medida finita.
\end{definicion}
\begin{ejemplo}
La medida nula vale cero en todos los conjuntos. La medida de conteo asigna a un conjunto su número de elementos, o infinito si es infinito: en una unión disjunta los elementos se cuentan una sola vez. Para $x_0\in X$ y $\alpha\in[0,\infty]$, la medida puntual $\alpha\delta_{x_0}$ vale $\alpha$ cuando el conjunto contiene $x_0$ y cero en caso contrario. En una familia disjunta, $x_0$ pertenece a lo sumo a un término, lo que prueba la aditividad. Una medida de probabilidad es una medida con $\mu(X)=1$.\index{medida!de conteo}\index{medida!de Dirac}
\end{ejemplo}
\begin{teorema}\label{teo:v2-contmedida}\index{continuidad!de una medida}
Una medida es monótona y numerablemente subaditiva. Si $A_n\uparrow A$, entonces $\mu(A_n)\uparrow\mu(A)$. Si $A\subseteq B$ y $\mu(B)<\infty$, entonces $\mu(B\setminus A)=\mu(B)-\mu(A)$.
\end{teorema}
\begin{proof}
De $B=A\mathbin{\dot\cup}(B\setminus A)$ resulta $\mu(B)=\mu(A)+\mu(B\setminus A)\ge\mu(A)$. Si $\mu(B)$ es finita, ambos sumandos son finitos y se puede restar. Si $A_n\uparrow A$, ponga $D_1=A_1$, $D_n=A_n\setminus A_{n-1}$: son disjuntos y $A_n=\bigcup_{j\le n}D_j$. Por ello $\mu(A)=\sum_j\mu(D_j)=\lim_n\mu(A_n)$, con igualdad también si la serie diverge a infinito. Para una sucesión arbitraria use, en cambio, $D_n=A_n\setminus\bigcup_{j<n}A_j$. Entonces
\[\mu\left(\bigcup_nA_n\right)=\sum_n\mu(D_n)\le\sum_n\mu(A_n).\]
La diferencia con todos los términos anteriores es indispensable: quitar solamente $A_{n-1}$ no produce disjunción en general. Esta es la corrección necesaria en la prueba de subaditividad de la lectura 2.
\end{proof}

\section{Ejercicios sobre medidas y límites de conjuntos}
\begin{ejercicio}[Lectura 2, ejercicio 1]\label{ej:v2-l02-1}
Para $A\in\Acal$ fijo, pruebe que $\lambda(B)=\mu(A\cap B)$ es una medida.
\end{ejercicio}

\begin{ejercicio}[Lectura 2, ejercicio 2]\label{ej:v2-l02-2}
Pruebe que las combinaciones finitas no negativas de medidas son medidas. Pruebe también que $\lambda=\sum_n\alpha_n\mu_n$ es una medida si $\alpha_n\ge0$ y las $\mu_n$ son medidas finitas.
\end{ejercicio}

\begin{ejercicio}[Lectura 2, ejercicio 3]\label{ej:v2-l02-3}
Si $A_n\downarrow A$ y algún $A_N$ tiene medida finita, pruebe que $\mu(A_n)\to\mu(A)$. Muestre la necesidad de la hipótesis de finitud.
\end{ejercicio}

\begin{ejercicio}[Lectura 2, ejercicio 4]\label{ej:v2-l02-4}
Pruebe $\mu(\liminf A_n)\le\liminf\mu(A_n)$, donde $\liminf A_n=\bigcup_m\bigcap_{n\ge m}A_n$.
\end{ejercicio}

\begin{ejercicio}[Lectura 2, ejercicio 5]\label{ej:v2-l02-5}
Si $\mu(\bigcup_nA_n)<\infty$, pruebe $\limsup\mu(A_n)\le\mu(\limsup A_n)$.
\end{ejercicio}

\begin{ejercicio}[Lectura 2, ejercicio 6]\label{ej:v2-l02-6}
Defina $\mu_*(A)=\sup\{\mu(B):B\in\Acal,\ B\subseteq A,\ \mu(B)<\infty\}$. Pruebe que es una medida y que coincide con $\mu$ si esta es $\sigma$-finita.
\end{ejercicio}

\begin{ejercicio}[Lectura 2, ejercicio 7: Borel--Cantelli]\label{ej:v2-l02-7}\index{Borel--Cantelli, lema de}
Si $\sum_n\mu(A_n)<\infty$, pruebe que $\mu(\limsup A_n)=0$.
\end{ejercicio}

\begin{ejercicio}[Lectura 2, ejercicio 8]\label{ej:v2-l02-8}
Pruebe que la $\sigma$-finitud equivale a tener $E_n\uparrow X$ con $\mu(E_n)<\infty$.
\end{ejercicio}

\chapter{Funciones medibles y aproximación topológica}\label{cap:v2-02}
La medibilidad describe qué información de una función puede expresarse mediante conjuntos medibles. Su definición utiliza preimágenes porque estas respetan complementos y uniones. Este capítulo desarrolla la lectura 3 de Valencia Bucio \cite{valencia2}; la aproximación por funciones simples se retomará al construir la integral.

\section{Conjuntos de nivel y operaciones}\label{lectura:v2-03}
\begin{definicion}\index{funcion@función!medible}
Una función $f:X\to\overline\R=[-\infty,\infty]$ es medible si $\{f>a\}\in\Acal$ para cada $a\in\R$. Para funciones reales se utiliza la misma definición, restringiendo el contradominio a $\R$.
\end{definicion}
\begin{proposicion}\label{prop:v2-niveles}
Son equivalentes las condiciones de medibilidad de todas las familias $\{f>a\}$, $\{f\ge a\}$, $\{f<a\}$ y $\{f\le a\}$, al variar $a\in\R$.
\end{proposicion}
\begin{proof}
Las identidades
\[\{f\ge a\}=\bigcap_{n\ge1}\{f>a-1/n\},\qquad
\{f>a\}=\bigcup_{n\ge1}\{f\ge a+1/n\}\]
prueban la equivalencia de las dos primeras condiciones. Por ejemplo, si $f(x)>a$, existe $n$ con $1/n<f(x)-a$, lo que explica la inclusión no inmediata en la segunda identidad. Las otras condiciones se obtienen de $\{f<a\}=X\setminus\{f\ge a\}$ y $\{f\le a\}=X\setminus\{f>a\}$. Se usan uniones e intersecciones numerables, nunca una unión arbitraria de conjuntos medibles.
\end{proof}
Las funciones constantes son medibles porque sus conjuntos de nivel son $X$ o $\varnothing$. La indicadora $\ind A$ es medible exactamente cuando $A\in\Acal$: el nivel $\{\ind A>1/2\}$ recupera $A$, y todos los demás niveles son $\varnothing,A$ o $X$. Toda función continua en un espacio topológico es Borel-medible, pues sus preimágenes de abiertos son abiertas.

\begin{teorema}\label{teo:v2-algebra-medibles}
Si $f,g:X\to\R$ son medibles, también lo son $af+bg$, $fg$ y $|f|$. Los supremos, ínfimos, límites superiores e inferiores de sucesiones de funciones medibles son funciones medibles con valores extendidos. En particular, todo límite puntual real de funciones medibles es medible.
\end{teorema}
\begin{proof}
Para $a>0$, $\{af>t\}=\{f>t/a\}$; para $a<0$ es $\{f<t/a\}$, y para $a=0$ la función es constante. Para la suma usamos
\[\{f+g>t\}=\bigcup_{r\in\Q}\bigl(\{f>r\}\cap\{g>t-r\}\bigr).\]
Si $f(x)+g(x)>t$, el intervalo $(t-g(x),f(x))$ contiene un racional $r$, que sitúa a $x$ en el lado derecho. La inclusión inversa se obtiene sumando las dos desigualdades estrictas.

Si $t\ge0$, $\{f^2>t\}=\{f>\sqrt t\}\cup\{f<-\sqrt t\}$; si $t<0$, el conjunto es $X$. Por tanto $f^2$ es medible y la identidad $fg=((f+g)^2-(f-g)^2)/4$ prueba el caso del producto. Para $|f|$ usamos los niveles $\{f>t\}\cup\{f<-t\}$ cuando $t\ge0$, y $X$ cuando $t<0$.

Para una sucesión $(f_n)$,
\[\{\sup_nf_n>t\}=\bigcup_n\{f_n>t\},\qquad
\{\inf_nf_n\ge t\}=\bigcap_n\{f_n\ge t\}.\]
Esto cubre también supremos e ínfimos finitos. Finalmente,
\[\limsup_nf_n=\inf_m\sup_{n\ge m}f_n,\qquad
\liminf_nf_n=\sup_m\inf_{n\ge m}f_n,\]
de modo que basta aplicar dos veces lo anterior. Cuando existe el límite, coincide con el límite inferior.
\end{proof}

\begin{proposicion}[Valores infinitos]\label{prop:v2-infinitos}
Una función $f:X\to\overline\R$ es medible si y solo si $A=\{f=\infty\}$ y $B=\{f=-\infty\}$ son medibles y la función real $g$, igual a $f$ en $X\setminus(A\cup B)$ y a cero en $A\cup B$, es medible.
\end{proposicion}
\begin{proof}
Si $f$ es medible, $A=\bigcap_n\{f>n\}$ y $B=\bigcap_n\{f<-n\}$ son medibles. Si $t\ge0$, $\{g>t\}=\{f>t\}\setminus A$; si $t<0$, $\{g>t\}=\{f>t\}\cup B$. Esto prueba que $g$ es medible. Recíprocamente, para $t\ge0$ se tiene $\{f>t\}=\{g>t\}\cup A$ y, para $t<0$, $\{f>t\}=\{g>t\}\setminus B$. Todas estas operaciones preservan medibilidad.
\end{proof}

\section{Semicontinuidad y aproximaciones continuas}
\begin{definicion}\index{semicontinuidad}
Una función real $f$ en un espacio métrico es semicontinua inferior en $x$ si, dado $a<f(x)$, existe una vecindad $U$ de $x$ en la que $f>a$. Es semicontinua superior si, dado $a>f(x)$, existe una vecindad donde $f<a$.
\end{definicion}
La semicontinuidad inferior permite saltos hacia abajo en el valor del punto; prohíbe que los valores cercanos estén uniformemente por debajo de ese valor. Su caracterización secuencial es $f(x)\le\liminf_nf(x_n)$ cuando $x_n\to x$: una vecindad para cada $a<f(x)$ prueba una dirección; si la definición falla, elegimos $x_n\in B(x,1/n)$ con $f(x_n)\le a<f(x)$ y contradicimos la otra.

\begin{proposicion}\label{prop:v2-aprox-semicont}
Toda función real semicontinua inferior en un espacio métrico es límite creciente de funciones continuas reales. Toda función semicontinua superior es límite decreciente de funciones continuas. En un espacio compacto, una función real semicontinua inferior alcanza su mínimo y una semicontinua superior alcanza su máximo.
\end{proposicion}
\begin{proof}
Primero suponga $0<g<1$ y $g$ semicontinua inferior. Defina
\[g_n(x)=\inf_{y\in X}\bigl(g(y)+n d(x,y)\bigr).\]
Como la expresión con $y=x$ vale $g(x)$, tenemos $g_n\le g<1$. La desigualdad triangular implica $g_n(x)\le g_n(z)+n d(x,z)$, primero comparando con cada $y$ y después tomando ínfimo; intercambiando $x,z$ vemos que $g_n$ es $n$-Lipschitz. La sucesión es creciente porque aumenta cada expresión dentro del ínfimo. Para probar $g_n(x)>0$, tome $\delta>0$ tal que $g(y)>g(x)/2$ si $d(x,y)<\delta$; fuera de esa bola la expresión es al menos $n\delta$. Así $g_n(x)\ge\min\{g(x)/2,n\delta\}>0$.

Si $0<a<g(x)$, la semicontinuidad da una bola de radio $\delta$ donde $g>a$; para $n\delta\ge a$, también fuera de la bola se tiene $g(y)+nd(x,y)\ge a$. Por ello $g_n(x)\ge a$ y $g_n(x)\uparrow g(x)$. Para una función real general $f$, aplique este argumento a $g=\frac12+\pi^{-1}\arctan f$ y componga $g_n$ con la inversa creciente $t\mapsto\tan(\pi(t-1/2))$. Las funciones resultantes son continuas reales y crecen a $f$. Para semicontinuidad superior se aplica el resultado a $-f$.

En el caso compacto, primero $f$ está acotada inferiormente: las vecindades donde $f>f(x)-1$ cubren $X$, y una subcubierta finita da una cota inferior común. Tome puntos $x_n$ con $f(x_n)\to\inf_Xf$. Una subsucesión converge a $x$, y la caracterización secuencial da $f(x)\le\liminf f(x_n)=\inf_Xf$. La desigualdad inversa es la definición de ínfimo. El máximo se deduce aplicando esto a $-f$.
\end{proof}

\section{Ejercicios sobre medibilidad y semicontinuidad}
\begin{ejercicio}[Lectura 3, ejercicio 1]\label{ej:v2-l03-1}
Pruebe que $f:X\to\R$ es medible si y solo si $f^{-1}(U)$ es medible para cada abierto $U\subseteq\R$.
\end{ejercicio}
\begin{solucion}
La condición para abiertos incluye $U=(a,\infty)$, así que implica la definición. Recíprocamente, la preimagen de $(a,b)$ es $\{f>a\}\cap\{f<b\}$. Todo abierto real es unión de los intervalos de extremos racionales contenidos en él. Hay solamente una cantidad numerable de esos intervalos, por lo que su preimagen es medible. Además, $\{B\subseteq\R:f^{-1}(B)\in\Acal\}$ es una $\sigma$-álgebra que contiene los abiertos; contiene por ello todos los borelianos.
\end{solucion}
\begin{ejercicio}[Lectura 3, ejercicio 2]\label{ej:v2-l03-2}
Extienda las propiedades de operaciones, supremos, ínfimos y límites a funciones medibles con valores en $\overline\R$. Precise cuándo las operaciones están definidas.
\end{ejercicio}
\begin{solucion}
Los supremos e ínfimos numerables siempre existen en $\overline\R$; las identidades de niveles del teorema \ref{teo:v2-algebra-medibles} siguen siendo exactas. Por tanto también son medibles los límites superiores e inferiores y cualquier límite puntual extendido.

La suma se permite cuando no aparece $\infty+(-\infty)$; una combinación lineal requiere la misma precaución después de multiplicar por los escalares, con $0f=0$. Para comprobar medibilidad, parta $X$ en los conjuntos donde cada función es real, $+\infty$ o $-\infty$. Son nueve piezas medibles por la proposición \ref{prop:v2-infinitos}. En la pieza donde ambas son finitas se usa el álgebra de funciones reales; en cada una de las demás, una suma definida vale un infinito de signo determinado o una función medible multiplicada por un escalar finito. Las preimágenes de niveles son uniones finitas de las correspondientes preimágenes en las piezas.

Para el producto adoptamos $0\cdot(\pm\infty)=0$. Subdividimos las piezas anteriores según los signos y los conjuntos de ceros. En cada pieza el producto es real medible, cero o un infinito de signo fijo. Por la misma razón es medible. El valor absoluto transforma ambas infinitudes en $+\infty$ y conserva la medibilidad. Estas afirmaciones no autorizan a escribir una suma indeterminada: si $f=\infty$ y $g=-\infty$ en un punto, $f+g$ no es una función allí sin una convención adicional explícita.
\end{solucion}
\begin{ejercicio}[Lectura 3, ejercicio 3]\label{ej:v2-l03-3}
Construya $f$ no medible tal que $|f|$ y $f^2$ sean medibles.
\end{ejercicio}
\begin{solucion}
Tome $X=\{a,b\}$ con $\Acal=\{\varnothing,X\}$ y defina $f(a)=1$, $f(b)=-1$. Entonces $\{f>0\}=\{a\}\notin\Acal$, mientras $|f|=f^2=1$ son constantes medibles. Así conocer el tamaño de una función no basta para conocer la medibilidad de su signo.
\end{solucion}
\begin{ejercicio}[Lectura 3, ejercicio 4]\label{ej:v2-l03-4}
Pruebe: (i) continuidad equivale a semicontinuidad superior e inferior; (ii) la función piso es semicontinua superior; (iii) semicontinuidad superior equivale a que todos los conjuntos $\{f<a\}$ sean abiertos, y la inferior a que todos los $\{f>a\}$ sean abiertos; (iv) $A$ es abierto, respectivamente cerrado, si y solo si $\ind A$ es semicontinua inferior, respectivamente superior.
\end{ejercicio}
\begin{solucion}
(i) La continuidad produce simultáneamente vecindades donde $f(x)-\eps<f(y)<f(x)+\eps$. Recíprocamente, interseque las dos vecindades obtenidas de las semicontinuidades con esos umbrales; en la intersección se cumple $|f(y)-f(x)|<\eps$.

(ii) Si $k=\lfloor x\rfloor$, elija una vecindad de $x$ contenida en $(-\infty,k+1)$. En ella $\lfloor y\rfloor\le k<a$ para cualquier $a>\lfloor x\rfloor$. En los enteros hay un salto que impide semicontinuidad inferior, pero no superior.

(iii) Si $f(x)<a$, la definición superior proporciona una vecindad de $x$ contenida en $\{f<a\}$; por tanto ese conjunto es abierto. Recíprocamente, para $a>f(x)$, ese mismo abierto es la vecindad requerida. Para la versión inferior intercambie las desigualdades o aplique el resultado a $-f$.

(iv) Para $0\le t<1$, $\{\ind A>t\}=A$, y los restantes niveles estrictos superiores son $X$ o vacío. Esto prueba la equivalencia con apertura. Para $0<t\le1$, $\{\ind A<t\}=A^c$, y los demás niveles inferiores son $X$ o vacío. Esto prueba la equivalencia con cerradura.
\end{solucion}
\begin{ejercicio}[Lectura 3, ejercicio 5]\label{ej:v2-l03-5}
Si $f:X\to\R$ es medible y $g:\R\to\R$ es semicontinua superior o inferior, pruebe que $g\circ f$ es medible.
\end{ejercicio}
\begin{solucion}
Si $g$ es semicontinua inferior, $\{g>a\}$ es abierto. Entonces $\{g\circ f>a\}=f^{-1}(\{g>a\})$ es medible por el ejercicio \ref{ej:v2-l03-1}. Si es semicontinua superior, use $\{g<a\}$ y la caracterización por niveles inferiores. La continuidad no es necesaria; la medibilidad de $g$ respecto de los borelianos sí es la propiedad pertinente.
\end{solucion}
\begin{ejercicio}[Lectura 3, ejercicio 6]\label{ej:v2-l03-6}
Pruebe que el conjunto de puntos donde una sucesión de funciones reales medibles tiene límite real es medible.
\end{ejercicio}
\begin{solucion}
La completitud de $\R$ identifica convergencia real con la condición de Cauchy. Por ello el conjunto buscado es
\[E=\bigcap_{r\ge1}\bigcup_{N\ge1}\bigcap_{m,n\ge N}
\{|f_m-f_n|<1/r\}.\]
Cada nivel es medible y todas las operaciones son numerables. Si $x\in E$, dado $\eps>0$ tome $r$ con $1/r<\eps$; el correspondiente $N$ prueba Cauchy. La implicación inversa resulta de aplicar Cauchy con $\eps=1/r$. Si se admiten límites extendidos, añada los dos conjuntos medibles $\bigcap_{r\ge1}\bigcup_N\bigcap_{n\ge N}\{f_n>r\}$ y su análogo $\{f_n<-r\}$. Esto distingue explícitamente límite finito de límite infinito.
\end{solucion}

\chapter{Construcción y extensión de medidas}\label{cap:v2-03}
La longitud está definida inicialmente en intervalos, pero queremos medir uniones, intersecciones y límites de conjuntos. La construcción siguiente separa dos problemas: asignar un tamaño exterior a \emph{todo} subconjunto y determinar en qué conjuntos ese tamaño es aditivo. Desarrollamos las lecturas 4--6 de Valencia Bucio \cite{valencia2}. La organización de Carathéodory por lemas y conjuntos de prueba también se inspira en Maximenko \cite{maximenkoCaratheodory}; aquí se justifican los pasos que allí se proponen como ejercicios.

\section{Medidas exteriores y conjuntos que separan la medida}\label{lectura:v2-04}
\begin{definicion}\index{medida!exterior}
Una medida exterior es una función $\mu^*:\mathcal P(X)\to[0,\infty]$ nula en vacío, monótona y numerablemente subaditiva. Un conjunto $E\subseteq X$ es \emph{Carathéodory-medible} si
\[\mu^*(T)=\mu^*(T\cap E)+\mu^*(T\setminus E)\qquad(T\subseteq X).\]
Denotamos por $\Mcal_{\mu^*}$ a la familia de esos conjuntos.
\end{definicion}
Siempre vale la desigualdad $\le$ por subaditividad. Por tanto basta probar la desigualdad inversa. El conjunto $T$ es arbitrario: no se presupone medible, pues precisamente estamos construyendo la familia medible.

\begin{ejemplo}
La función que vale $0$ en vacío y $1$ en conjuntos no vacíos es una medida exterior. Si una unión es no vacía, al menos un término lo es y la suma de tamaños es al menos $1$. Otra medida exterior vale $0$ en los conjuntos numerables y $1$ en los no numerables: una unión numerable de numerables es numerable, y si la unión no lo es, al menos un término no lo es. Ninguna es aditiva sobre conjuntos disjuntos en general.
\end{ejemplo}
\begin{lema}
Si $\mu^*(E)=0$, entonces $E\in\Mcal_{\mu^*}$. La familia $\Mcal_{\mu^*}$ es estable bajo complementos.
\end{lema}
\begin{proof}
Para todo $T$, $\mu^*(T\cap E)+\mu^*(T\setminus E)\le0+\mu^*(T)$. Junto con la subaditividad esto da igualdad. Cambiar $E$ por $E^c$ solamente intercambia los dos sumandos de la condición, por lo que preserva medibilidad. En particular, también es medible cualquier conjunto de complemento exterior nulo.
\end{proof}
\begin{teorema}[Carathéodory]\label{teo:v2-carath}\index{Caratheodory, teorema de@Carathéodory, teorema de}
La familia $\Mcal_{\mu^*}$ es una $\sigma$-álgebra y $\mu^*|_{\Mcal_{\mu^*}}$ es una medida completa.
\end{teorema}
\begin{proof}
El vacío es medible por el lema; también lo es $X$. Si $E,F$ son medibles, parta primero $T$ mediante $E$ y después $T\setminus E$ mediante $F$:
\[\mu^*(T)=\mu^*(T\cap E)+\mu^*((T\setminus E)\cap F)+\mu^*(T\setminus(E\cup F)).\]
Los dos primeros conjuntos cubren $T\cap(E\cup F)$, por lo que su suma de tamaños es al menos el tamaño de esa unión. Así obtenemos la desigualdad que prueba medibilidad de $E\cup F$. Con los complementos, esto muestra que $\Mcal_{\mu^*}$ es un álgebra.

Sean ahora $E_j$ medibles disjuntos y $E=\bigcup_jE_j$. Por inducción,
\[\mu^*(T)=\sum_{j=1}^m\mu^*(T\cap E_j)+\mu^*\left(T\setminus\bigcup_{j=1}^mE_j\right).\]
El paso inductivo divide el último resto mediante $E_{m+1}$; la disjunción hace que la parte que se extrae sea exactamente $T\cap E_{m+1}$. Como $T\setminus E$ está contenido en ese resto,
\[\mu^*(T)\ge\sum_{j=1}^m\mu^*(T\cap E_j)+\mu^*(T\setminus E).\]
Al tomar el supremo sobre $m$ y usar subaditividad,
\[\mu^*(T)\ge\sum_{j=1}^\infty\mu^*(T\cap E_j)+\mu^*(T\setminus E)
\ge\mu^*(T\cap E)+\mu^*(T\setminus E).\]
Por tanto $E$ es medible. Una unión numerable arbitraria se convierte en disjunta mediante el ejercicio \ref{ej:v2-l01-2}; así la familia es una $\sigma$-álgebra. Tomando $T=E$ en la desigualdad anterior resulta $\mu^*(E)\ge\sum_j\mu^*(E_j)$; la desigualdad opuesta ya es subaditividad. Esto prueba aditividad numerable.

Finalmente, si $N$ es medible con medida cero y $A\subseteq N$, la monotonía da $\mu^*(A)=0$ y el lema prueba que $A$ es medible. Esta es la completitud. La prueba nunca resta cantidades infinitas.
\end{proof}

\section{Ejercicios sobre medidas exteriores}
\begin{ejercicio}[Lectura 4, ejercicio 1]\label{ej:v2-l04-1}
Pruebe que la medida obtenida por Carathéodory es completa.
\end{ejercicio}

\begin{ejercicio}[Lectura 4, ejercicio 2]\label{ej:v2-l04-2}
Pruebe que $\rho(A)=\inf_{B\subseteq A}[\mu^*(B)+\nu^*(A\setminus B)]$ es una medida exterior si $\mu^*,\nu^*$ lo son.
\end{ejercicio}

\begin{ejercicio}[Lectura 4, ejercicio 3]\label{ej:v2-l04-3}
Si $Y\subseteq X$, pruebe que (i) $\mu^*|_{\mathcal P(Y)}$ es una medida exterior sobre $Y$; (ii) $E\cap Y$ es medible para esta medida exterior si $E$ es $\mu^*$-medible.
\end{ejercicio}

\begin{ejercicio}[Lectura 4, ejercicio 4]\label{ej:v2-l04-4}
Si $R\subseteq X\times Y$ es una relación, pruebe que $\rho(A)=\mu^*(R^{-1}(A))$ es una medida exterior en $Y$.
\end{ejercicio}

\begin{ejercicio}[Lectura 4, ejercicio 5]\label{ej:v2-l04-5}
Si $E$ es $\mu^*$-medible, pruebe $\mu^*(E\cap A)+\mu^*(E\cup A)=\mu^*(E)+\mu^*(A)$ para todo $A\subseteq X$.
\end{ejercicio}


\section{De un anillo a una medida}\label{lectura:v2-05}
\begin{definicion}\index{anillo de conjuntos}\index{premedida}\index{volumen}
Un anillo de conjuntos es una familia no vacía estable bajo uniones finitas y diferencias; contiene vacío e intersecciones finitas. Es un $\sigma$-anillo si también contiene las uniones numerables. Un \emph{volumen} en un anillo es una función $v$ no negativa, nula en vacío y finitamente aditiva. Se llama \emph{premedida} si es numerablemente aditiva siempre que la unión disjunta considerada pertenezca al anillo.
\end{definicion}
Las álgebras son anillos que contienen $X$. Los subconjuntos acotados de un espacio métrico forman un anillo; si el espacio es no acotado, no forman un $\sigma$-anillo, pues una unión numerable de bolas puede ser no acotada. Si el espacio es acotado, en cambio, constituyen toda $\mathcal P(X)$: la salvedad corrige el ejemplo 3(ii) de la lectura 5.

Las uniones finitas de rectángulos de dos anillos forman un anillo: se usan $ (A\times B)\cap(C\times D)=(A\cap C)\times(B\cap D)$ y
\[(A\times B)\setminus(C\times D)=((A\setminus C)\times B)\mathbin{\dot\cup}((A\cap C)\times(B\setminus D)).\]
Una partición finita común de los factores convierte toda representación en rectángulos disjuntos. Si los factores llevan volúmenes, la suma de los productos de sus volúmenes es independiente de la partición: al refinar, la distributividad y la aditividad finita preservan la suma. Esto define el volumen producto.

\begin{observacion}[Otra corrección sobre anillos]
Las uniones numerables de intervalos $(a,b]$ no forman, en general, el $\sigma$-anillo afirmado en el ejemplo 3(iv) de la lectura 5. La estabilidad bajo diferencias falla, pues
\[(0,2]\setminus\bigcup_{n\ge2}(1,2-1/n]= (0,1]\cup\{2\}.
\]
El punto aislado $2$ no puede cubrirse por intervalos $(a,b]$ contenidos en ese conjunto. La familia segura para la construcción es el anillo de \emph{uniones finitas} de intervalos semiabiertos; su $\sigma$-álgebra generada es $\Bcal(\R)$.
\end{observacion}

\begin{proposicion}\label{prop:v2-volumen}
Un volumen es monótono y finitamente subaditivo. Si $A_1,\ldots,A_m\subseteq B$ son disjuntos, $\sum_jv(A_j)\le v(B)$. Para una premedida, $B\subseteq\bigcup_nA_n$ implica $v(B)\le\sum_nv(A_n)$.
\end{proposicion}
\begin{proof}
Parta $B$ en los $A_j$ y el resto para la primera desigualdad. Para conjuntos no disjuntos, disjunte una cubierta finita y aplique monotonía a sus piezas. En el caso numerable defina $D_n=(B\cap A_n)\setminus\bigcup_{j<n}A_j$. Son elementos del anillo, disjuntos, y su unión es $B$, que pertenece al anillo. Por la premedida,
$v(B)=\sum_nv(D_n)\le\sum_nv(A_n)$.
\end{proof}
\begin{teorema}[Extensión de una premedida]\label{teo:v2-extension}
Si $v$ es una premedida en un anillo $\mathcal R$, defina
\[v^*(E)=\inf\left\{\sum_nv(R_n):R_n\in\mathcal R, E\subseteq\bigcup_nR_n\right\},\]
con ínfimo del conjunto vacío igual a infinito. Entonces $v^*$ es una medida exterior, cada $R\in\mathcal R$ es $v^*$-medible y $v^*(R)=v(R)$. En consecuencia, la restricción a $\sigma(\mathcal R)$ extiende $v$.
\end{teorema}
\begin{proof}
Cubrir vacío mediante vacíos prueba $v^*(\varnothing)=0$. Toda cubierta de un conjunto también cubre sus subconjuntos, lo que prueba monotonía. Para $E=\bigcup_jE_j$, si $\sum_jv^*(E_j)=\infty$ no hay nada que probar. De otro modo, para cada $j$ elija una cubierta $(R_{j,k})_k$ con suma menor que $v^*(E_j)+\eps2^{-j}$. La familia doble es numerable y cubre $E$, de modo que $v^*(E)\le\sum_jv^*(E_j)+\eps$. Hacer $\eps\downarrow0$ prueba subaditividad.

Para $R\in\mathcal R$, la cubierta $R,\varnothing,\ldots$ da $v^*(R)\le v(R)$; toda cubierta numerable da la desigualdad inversa por la proposición \ref{prop:v2-volumen}. Para probar medibilidad de $R$, fije $E\subseteq X$ y una cubierta $(R_n)$ de $E$. Las familias $(R_n\cap R)$ y $(R_n\setminus R)$ cubren $E\cap R$ y $E\setminus R$ respectivamente, luego
\[v^*(E\cap R)+v^*(E\setminus R)\le\sum_n[v(R_n\cap R)+v(R_n\setminus R)]=\sum_nv(R_n).\]
Tomar ínfimo prueba la desigualdad de Carathéodory, aun si no existe una cubierta de suma finita. El teorema \ref{teo:v2-carath} completa la extensión.
\end{proof}
\begin{observacion}
La observación 7 de la lectura 5 utiliza el ínfimo sobre \emph{un solo conjunto} del anillo que contiene a $E$. Esa fórmula no siempre es una medida exterior: en el anillo de los subconjuntos finitos de $\N$, con el volumen idénticamente cero, asignaría cero a cada singleton e infinito a $\N$. Fallaría la subaditividad numerable. Las cubiertas numerables de la fórmula anterior reparan la construcción.
\end{observacion}
En un espacio métrico no acotado, los conjuntos acotados o de complemento acotado forman un álgebra. Asignarles respectivamente $0$ y una constante $\alpha>0$ define un volumen: dos conjuntos coacotados no pueden ser disjuntos, y añadir un acotado a otro acotado sigue dando uno acotado. No es premedida: las coronas obtenidas de una cubierta creciente por bolas tienen volumen cero, mientras su unión $X$ tiene volumen $\alpha$.

\section{Ejercicios sobre premedidas}
\begin{ejercicio}[Lectura 5, ejercicio 1]\label{ej:v2-l05-1}
Si $v(X)=1$ y $\sum_{k=1}^nv(A_k)>n-1$, pruebe que $\bigcap_{k=1}^nA_k\ne\varnothing$.
\end{ejercicio}

\begin{ejercicio}[Lectura 5, ejercicio 2]\label{ej:v2-l05-2}
Pruebe que un volumen es premedida si y solo si $v(B)\le\sum_nv(A_n)$ para toda cubierta $B\subseteq\bigcup_nA_n$ con $B,A_n$ en el anillo.
\end{ejercicio}

\begin{ejercicio}[Lectura 5, ejercicio 3]\label{ej:v2-l05-3}
Pruebe que un volumen en un álgebra es premedida si y solo si es continuo por abajo, siempre que la unión creciente pertenezca al álgebra.
\end{ejercicio}

\begin{ejercicio}[Lectura 5, ejercicio 4]\label{ej:v2-l05-4}
Si $v(X)<\infty$, pruebe que una partición de $X$ por elementos del álgebra de volumen estrictamente positivo es a lo sumo numerable.
\end{ejercicio}


\section{Longitud y medida de Lebesgue}\label{lectura:v2-06}\index{Lebesgue, medida de}
\begin{lema}[El volumen de cajas es una premedida]\label{lema:v2-cajas}
En $\R^d$, las uniones finitas disjuntas de cajas semiabiertas acotadas $\prod_{j=1}^d[a_j,b_j)$ forman un anillo. La regla que asigna a cada caja $\prod_j(b_j-a_j)$ y suma sobre una partición finita define una premedida.
\end{lema}
\begin{proof}
Dados finitos conjuntos de este tipo, corte cada eje por todos los extremos que aparecen. Los productos de los intervalos resultantes forman una partición finita común. Toda intersección y diferencia es unión de celdas de esa partición. La distributividad del producto muestra que la suma de los volúmenes de las celdas de una caja es su volumen; esto prueba independencia de representación y aditividad finita.

Para la subaditividad numerable, suponga que una caja $Q$ está cubierta por cajas $Q_n$ y que la suma de sus volúmenes es finita. Dado $\eps>0$, agrande ligeramente cada $Q_n$ a una caja abierta $U_n$ de volumen menor que $v(Q_n)+\eps2^{-n}$; esto es posible por continuidad del producto de longitudes. Dentro de $Q$ tome una caja compacta $K$ obtenida apartando una distancia $\delta$ de sus caras. Una subfamilia finita de los $U_n$ cubre $K$. La subaditividad finita, verificada con la misma partición por extremos y válida también para cajas cerradas o abiertas, da
\[v(K)\le\sum_nv(U_n)<\sum_nv(Q_n)+\eps.\]
Al hacer $\delta\downarrow0$ y después $\eps\downarrow0$, resulta $v(Q)\le\sum_nv(Q_n)$. Para una unión finita disjunta $R$ de cajas, aplique este argumento a cada caja y a las intersecciones de la cubierta con ella; la suma finita de los volúmenes de esas intersecciones no excede el volumen de cada conjunto de la cubierta. Obtenemos $v(R)\le\sum_nv(R_n)$ en el anillo. El ejercicio \ref{ej:v2-l05-2} prueba que $v$ es premedida.
\end{proof}
\begin{definicion}
La medida exterior de Lebesgue $\lambda_d^*$ es la medida exterior inducida por esta premedida. La $\sigma$-álgebra de Lebesgue $\Lcal_d$ consta de sus conjuntos Carathéodory-medibles; $\lambda_d=\lambda_d^*|_{\Lcal_d}$. En dimensión uno escribimos $\lambda$ y $\Lcal$.
\end{definicion}
La definición mediante cubiertas por intervalos compactos de las lecturas 4 y 6 produce la misma medida exterior: un semiabierto está contenido en su cierre de igual longitud, y cada intervalo compacto admite una ampliación semiabierta cuya longitud adicional es menor que cualquier tolerancia prefijada. Asignando a la ampliación $n$ un error $\eps2^{-n}$ y sumando se obtienen las dos desigualdades entre ínfimos. El mismo argumento permite cubiertas por abiertos.

\begin{teorema}\label{teo:v2-lebesgue-prop}
La medida de Lebesgue es completa, mide cada caja por su volumen, contiene todos los borelianos y es invariante bajo traslaciones. Para todo conjunto Lebesgue-medible $A\subseteq\R^d$,
\[\lambda_d(A)=\inf_{U\supseteq A, U\text{ abierto}}\lambda_d(U)
=\sup_{K\subseteq A, K\text{ compacto}}\lambda_d(K).\]
Su restricción a los borelianos es la única medida que asigna a cada caja semiabierta su volumen.
\end{teorema}
\begin{proof}
Completitud y volumen de las cajas se siguen de los teoremas \ref{teo:v2-carath} y \ref{teo:v2-extension}. Las cajas semiabiertas generan los borelianos: cualquier caja abierta es unión numerable de cajas semiabiertas interiores y las cajas abiertas de extremos racionales forman una base numerable. Toda caja del anillo es medible; así lo es su $\sigma$-álgebra generada.

Trasladar una cubierta por cajas preserva los volúmenes y cubre el conjunto trasladado, luego $\lambda_d^*(A+t)\le\lambda_d^*(A)$; trasladando por $-t$ obtenemos igualdad. Para probar medibilidad de $A+t$, aplique el criterio de $A$ al conjunto $T-t$ y traslade sus dos partes. Por tanto la traslación preserva también $\Lcal_d$.

Para regularidad exterior, si $\lambda_d(A)<\infty$, cubra $A$ por cajas de suma menor que $\lambda_d(A)+\eps/2$, y amplíe cada una a una caja abierta con aumento menor que $\eps2^{-n-1}$. Su unión abierta $U$ satisface $\lambda_d(U)<\lambda_d(A)+\eps$. Si la medida es infinita, todo superset tiene medida infinita y la igualdad es inmediata.

Para regularidad interior, si $A$ es acotado elija una caja compacta $C\supseteq A$. La regularidad exterior de $C\setminus A$ proporciona un abierto $U\supseteq C\setminus A$ con $\lambda_d(U)<\lambda_d(C\setminus A)+\eps$. Entonces $K=C\setminus U$ es compacto y contenido en $A$. Como $C\subseteq K\cup U$,
\[\lambda_d(K)\ge\lambda_d(C)-\lambda_d(U)>\lambda_d(A)-\eps.\]
Todas las cantidades restadas son finitas. Para $A$ no acotado, aplique este resultado a $A\cap[-n,n]^d$ y use continuidad por abajo; esto cubre también $\lambda_d(A)=\infty$.

Por último, si $\nu$ coincide con el volumen en cajas, toda cubierta de un boreliano $A$ por cajas da $\nu(A)\le\sum_nv(Q_n)$ y, al tomar ínfimo, $\nu(A)\le\lambda_d(A)$. Si $A\subseteq C$ para una caja de volumen finito, aplique la misma desigualdad a $C\setminus A$. De
\[\nu(A)+\nu(C\setminus A)=\nu(C)=\lambda_d(C)=\lambda_d(A)+\lambda_d(C\setminus A)\]
y las dos desigualdades resulta igualdad en ambas. Un conjunto boreliano arbitrario se agota mediante intersecciones con cajas crecientes, y continuidad por abajo concluye $\nu(A)=\lambda_d(A)$.
\end{proof}
\begin{teorema}[Steinhaus]\label{teo:v2-steinhaus}\index{Steinhaus, teorema de}
Si $A\subseteq\R$ es Lebesgue-medible y $\lambda(A)>0$, entonces $A-A$ contiene una vecindad de cero.
\end{teorema}
\begin{proof}
Por regularidad interior existe un compacto $K\subseteq A$ con $0<\lambda(K)<\infty$. Por regularidad exterior existe un abierto $U\supseteq K$ con $\lambda(U)<2\lambda(K)$. La función continua $x\mapsto\dist(x,U^c)$ es positiva en $K$ y alcanza un mínimo positivo $\delta$; $U^c$ es no vacío porque $U$ tiene medida finita. Si $|t|<\delta$, entonces $K+t\subseteq U$. Si $K$ y $K+t$ fueran disjuntos, la invariancia por traslación daría $2\lambda(K)=\lambda(K\cup(K+t))\le\lambda(U)$, contradicción. Por tanto existen $x,y\in K$ con $x=y+t$, y $t=x-y\in K-K\subseteq A-A$.
\end{proof}

\section{Ejercicios sobre conjuntos de Lebesgue}
\begin{ejercicio}[Lectura 6, ejercicio 1]\label{ej:v2-l06-1}
Pruebe que todo abierto de $\R$ es unión numerable disjunta de intervalos semiabiertos.
\end{ejercicio}

\begin{ejercicio}[Lectura 6, ejercicio 2]\label{ej:v2-l06-2}
Pruebe que el conjunto ternario de Cantor tiene medida cero.
\end{ejercicio}

\begin{ejercicio}[Lectura 6, ejercicio 3]\label{ej:v2-l06-3}
Pruebe la equivalencia entre la condición de medida cero mediante cubiertas de intervalos y medida de Lebesgue cero.
\end{ejercicio}

\begin{ejercicio}[Lectura 6, ejercicio 4]\label{ej:v2-l06-4}
Pruebe que $A\subseteq\R$ es Lebesgue-medible si y solo si $A=B\cup Z$, con $B$ de tipo $F_\sigma$ y $Z$ nulo.
\end{ejercicio}

\begin{ejercicio}[Lectura 6, ejercicio 5]\label{ej:v2-l06-5}
Pruebe $\lambda^*(A+t)=\lambda^*(A)$ para todo subconjunto $A\subseteq\R$.
\end{ejercicio}

\begin{ejercicio}[Lectura 6, ejercicio 6]\label{ej:v2-l06-6}
Sea $\nu$ una medida boreliana en $\R$, finita en acotados e invariante bajo traslaciones, con $\nu(\R)>0$. Pruebe $\nu=\alpha\lambda$ en borelianos para algún $\alpha>0$.
\end{ejercicio}

\begin{ejercicio}[Lectura 6, ejercicio 7: Smith--Volterra--Cantor]\label{ej:v2-l06-7}\index{Smith--Volterra--Cantor, conjunto de}
En la etapa $n\ge1$, retire de cada uno de los $2^{n-1}$ intervalos restantes su intervalo central abierto de longitud $4^{-n}$. Pruebe que la intersección $C$ es cerrada, no vacía, sin interior ni puntos aislados, y tiene medida $1/2$.
\end{ejercicio}

\begin{ejercicio}[Lectura 6, ejercicio 8]\label{ej:v2-l06-8}
Construya un boreliano $A$ tal que $0<\lambda(A\cap I)<\lambda(I)$ para todo intervalo abierto acotado no vacío $I$.
\end{ejercicio}

\begin{ejercicio}[Lectura 6, ejercicio 9]\label{ej:v2-l06-9}
Construya $A_n\subseteq[0,1]$ decrecientes, con $\lambda^*(A_n)=1$ y $\bigcap_nA_n=\varnothing$.
\end{ejercicio}

\begin{ejercicio}[Lectura 6, ejercicio 10]\label{ej:v2-l06-10}
Defina la medida de Lebesgue en $\R^d$; pruebe allí medibilidad boreliana, regularidad y unicidad por volúmenes de cajas, así como las versiones multidimensionales de los ejercicios 4, 5 y 6.
\end{ejercicio}

\begin{ejercicio}[Lectura 6, ejercicio 11]\label{ej:v2-l06-11}
Pruebe que $\lambda_2$ es invariante bajo isometrías del plano.
\end{ejercicio}

\parte{Construcción y propiedades de la integral}{La integral se construye desde funciones simples y aproximaciones monótonas. Después se justifican los pasos al límite y las integrales iteradas. Cada intercambio se apoya en una hipótesis concreta: positividad, dominación, integrabilidad o $\sigma$-finitud. Esta parte utiliza la teoría de medida de los tres primeros capítulos.}
\chapter{Funciones simples e integral no negativa}\label{cap:v2-04}
La integral comienza sumando valores multiplicados por medidas. Para pasar a funciones generales se aproxima desde abajo y se toma un supremo. El orden de esta construcción evita suponer de antemano un teorema de intercambio de límites e integrales. Desarrollamos aquí las lecturas 7 y 8 de Valencia Bucio \cite{valencia2}.

\section{Particiones y aproximación por funciones simples}\label{lectura:v2-07}
\begin{definicion}\index{funcion@función!simple}\index{representacion@representación!canonica de una funcion simple@canónica de una función simple}
Una función simple es una función medible de rango finito contenido en $\R$. Si sus valores distintos son $a_1,\ldots,a_m$, su representación canónica es
\[s=\sum_{j=1}^ma_j\ind{A_j},\qquad A_j=\{s=a_j\}.\]
Los $A_j$ son medibles, no vacíos, disjuntos y cubren $X$. También se permiten representaciones no canónicas, con valores repetidos, conjuntos vacíos o conjuntos que se superponen.
\end{definicion}
\begin{proposicion}
Las combinaciones lineales, productos, valores absolutos y máximos o mínimos finitos de funciones simples son simples.
\end{proposicion}
\begin{proof}
Si $s=\sum_ja_j\ind{A_j}$ y $t=\sum_kb_k\ind{B_k}$ son canónicas, los conjuntos $A_j\cap B_k$ forman una partición finita después de descartar los vacíos. En cada pieza, $\alpha s+\beta t$, $st$, $\max(s,t)$ y $\min(s,t)$ toman respectivamente los valores $\alpha a_j+\beta b_k$, $a_jb_k$, $\max(a_j,b_k)$ y $\min(a_j,b_k)$. Todos los niveles se obtienen uniendo finitas piezas medibles. Para $|s|$, se reemplaza cada valor $a_j$ por $|a_j|$. Iterar el procedimiento prueba el caso de cualquier número finito de funciones.
\end{proof}

\begin{teorema}[Aproximación simple]\label{teo:v2-aproxsimple}
Si $f:X\to[0,\infty]$ es medible, existe una sucesión creciente de funciones simples no negativas $s_n$ que converge puntualmente a $f$. Si $f$ está acotada, se pueden elegir con convergencia uniforme.
\end{teorema}
\begin{proof}
Para $f(x)<n$, defina $s_n(x)=2^{-n}\lfloor2^nf(x)\rfloor$; para $f(x)\ge n$, incluyendo $f(x)=\infty$, defina $s_n(x)=n$. Su rango es finito y sus fibras se describen mediante bandas medibles $\{k2^{-n}\le f<(k+1)2^{-n}\}$ y $\{f\ge n\}$. Por ello es simple. El refinamiento de la malla diádica y el aumento del nivel de truncamiento implican $s_{n+1}\ge s_n$: por debajo de $n$, el redondeo con malla mitad no disminuye; por encima de $n$, el nuevo valor es al menos $n$.

Siempre $0\le s_n\le f$. Si $f(x)<\infty$, para $n>f(x)$ se cumple $0\le f(x)-s_n(x)<2^{-n}$, de donde la convergencia. Si $f(x)=\infty$, $s_n(x)=n\to\infty$. Cuando $f\le M$, la misma cota vale simultáneamente para todos los puntos una vez que $n>M$, lo que prueba convergencia uniforme.
\end{proof}
\begin{observacion}[La construcción racional de la lectura original]
Enumere los racionales no negativos como $(r_j)$ y defina $u_n=\max_{j\le n}r_j\ind{\{f>r_j\}}$. Son simples, crecientes y menores o iguales que $f$. Si $0<a<f(x)$, elija un racional $r_j\in(a,f(x))$; para $n\ge j$ se tiene $u_n(x)\ge r_j>a$. Como esto vale para todo $a<f(x)$, $u_n(x)\to f(x)$, incluyendo $f(x)=\infty$. Si $f(x)=0$, todos los términos son cero. Este argumento completa también el caso cero, que no permite una sucesión creciente de racionales no negativos estrictamente menores que $f(x)$.
\end{observacion}

\section{La integral de una función simple}\label{lectura:v2-08}
\begin{definicion}\index{integral!de una funcion simple@de una función simple}
Para una función simple no negativa con representación canónica $s=\sum_ja_j\ind{A_j}$, se define
\[I(s)=\int_Xs\,d\mu=\sum_ja_j\mu(A_j)\in[0,\infty].\]
La convención $0\cdot\infty=0$ asegura que la zona donde $s=0$ no contribuya, aunque tenga medida infinita.
\end{definicion}
\begin{proposicion}\label{prop:v2-integralsimple}
Si $s=\sum_{j=1}^mb_j\ind{B_j}$ con $b_j\ge0$ y $B_j$ medibles, entonces $I(s)=\sum_jb_j\mu(B_j)$, aunque los conjuntos no sean disjuntos. Para funciones simples no negativas, $I(s+t)=I(s)+I(t)$, $I(as)=aI(s)$ si $a\ge0$, y $s\le t$ implica $I(s)\le I(t)$.
\end{proposicion}
\begin{proof}
Forme las celdas $C_\eta=\bigcap_{j=1}^mB_j^{\eta_j}$, donde $\eta\in\{0,1\}^m$, $B_j^1=B_j$ y $B_j^0=B_j^c$. Son disjuntas y cubren $X$. En $C_\eta$, el valor de $s$ es $\sum_jb_j\eta_j$. Agrupar celdas de igual valor recupera la representación canónica; por aditividad finita,
\[I(s)=\sum_\eta\left(\sum_jb_j\eta_j\right)\mu(C_\eta)
=\sum_jb_j\sum_{\eta:\eta_j=1}\mu(C_\eta)
=\sum_jb_j\mu(B_j).\]
Todas las sumas son finitas y no negativas, por lo que no hay cancelaciones indeterminadas. Aplicar esta fórmula a las representaciones concatenadas de $s+t$ prueba aditividad; reemplazar $b_j$ por $ab_j$ prueba homogeneidad. Si $s\le t$, $t-s$ es simple no negativa y $I(t)=I(s)+I(t-s)\ge I(s)$.
\end{proof}

\begin{lema}[Continuidad creciente para un límite simple]\label{lema:v2-limitesimple}
Si $s_n$ son simples no negativas, $s_n\uparrow s$ y $s$ es simple, entonces $I(s_n)\uparrow I(s)$.
\end{lema}
\begin{proof}
La monotonía da $L=\lim_nI(s_n)\le I(s)$. Fije $0<c<1$ y escriba $s=\sum_{j=1}^ma_j\ind{A_j}$, omitiendo el valor cero. Para cada $j$ ponga $A_{n,j}=A_j\cap\{s_n\ge ca_j\}$. Estos conjuntos crecen a $A_j$: si $s(x)=a_j>0$, la convergencia a $a_j$ obliga a superar el umbral estrictamente menor $ca_j$. Las funciones $t_n=\sum_jca_j\ind{A_{n,j}}$ satisfacen $t_n\le s_n$, de modo que
\[L\ge\lim_nI(t_n)=\sum_jca_j\lim_n\mu(A_{n,j})=c\sum_ja_j\mu(A_j)=cI(s).\]
Si $I(s)=\infty$, ya esta desigualdad fuerza $L=\infty$. Si es finita, haga $c\uparrow1$. Se ha usado solamente continuidad de la medida, no un teorema de convergencia de integrales que todavía no esté probado.
\end{proof}

\section{El supremo que define la integral}
\begin{definicion}\index{integral!de Lebesgue}
Para $f:X\to[0,\infty]$ medible, se define
\[\int_Xf\,d\mu=\sup\{I(s):s\text{ simple},\ 0\le s\le f\}.\]
Una función no negativa es integrable cuando este supremo es finito.
\end{definicion}
Si $f$ es simple, la definición coincide con $I(f)$: la monotonía da $I(s)\le I(f)$ y la función $s=f$ es admisible. Si $f\le g$, toda aproximación inferior de $f$ también lo es de $g$ y, por tanto, $\int f\le\int g$.

\begin{teorema}\label{teo:v2-integral-aprox}
Si $s_n$ son simples no negativas y $s_n\uparrow f$, entonces
\[\int_Xf\,d\mu=\lim_n\int_Xs_n\,d\mu.\]
\end{teorema}
\begin{proof}
El límite $L$ existe por monotonía y es menor o igual que $\int f$. Fije una simple $0\le t\le f$. Las funciones $t_n=\min(t,s_n)$ son simples, crecen a $t$ y satisfacen $t_n\le s_n$. El lema \ref{lema:v2-limitesimple} da $I(t)=\lim_nI(t_n)\le L$. Tomar supremo sobre todas las simples $t\le f$ produce $\int f\le L$. La conclusión no depende de una aproximación simple particular.
\end{proof}
\begin{proposicion}\label{prop:v2-integral-positiva}
Para funciones medibles no negativas $f,g$ y $a\ge0$,
\[\int(f+g)=\int f+\int g,\qquad \int af=a\int f.\]
Si dos funciones medibles no negativas coinciden casi en todas partes, tienen la misma integral, finita o infinita.
\end{proposicion}
\begin{proof}
Elija simples crecientes $s_n\uparrow f$ y $t_n\uparrow g$. Entonces $s_n+t_n\uparrow f+g$. El teorema anterior y la aditividad simple dan
\[\int(f+g)=\lim_n[I(s_n)+I(t_n)]=\int f+\int g.\]
La igualdad de límites vale en $[0,\infty]$: si algún sumando diverge, también la suma diverge; si ambos son finitos, es el álgebra usual de límites. La homogeneidad para $a>0$ se obtiene con $as_n\uparrow af$; para $a=0$ ambos lados son cero por convención.

Si $N$ es medible nulo, toda simple no negativa que se anula fuera de $N$ tiene integral cero. Por definición, cualquier función medible no negativa que se anula fuera de $N$ tiene integral cero. Dividir $f=f\ind{X\setminus N}+f\ind N$ muestra que su integral no cambia al anularla en $N$. Aplicando esto a dos funciones iguales fuera de $N$, obtenemos igualdad de sus integrales.
\end{proof}

\section{Ejercicios sobre la construcción de la integral}
La lectura 7 no contiene una lista final de ejercicios; sus comprobaciones sobre máximos, mínimos, operaciones y aproximación quedaron desarrolladas en las pruebas anteriores.
\begin{ejercicio}[Lectura 8, ejercicio 1]\label{ej:v2-l08-1}
Justifique la fórmula de la integral para una representación no canónica $s=\sum_{k=1}^m\alpha_k\ind{A_k}$ de una función simple no negativa.
\end{ejercicio}
\begin{solucion}
Si $\alpha_k\ge0$, las celdas determinadas por todas las pertenencias a los $A_k$ forman una partición finita. En cada celda $C$, $s$ es la suma de los coeficientes cuyos conjuntos contienen $C$. Al sumar ese valor multiplicado por $\mu(C)$ e intercambiar las dos sumas finitas, el coeficiente $\alpha_k$ queda multiplicado por la suma de las medidas de las celdas que componen $A_k$, esto es, $\mu(A_k)$. Por tanto $\int s=\sum_k\alpha_k\mu(A_k)$, como se detalló en la proposición \ref{prop:v2-integralsimple}.

Si se permiten coeficientes negativos, la fórmula exige que el lado derecho esté definido; por ejemplo, basta que todos los conjuntos con coeficiente no nulo tengan medida finita. Sin esa precaución, $0=\ind X-\ind X$ en un espacio de medida infinita produciría $\infty-\infty$. La no negatividad de $s$ por sí sola no elimina ese problema en una representación con coeficientes de signos distintos.
\end{solucion}
\begin{ejercicio}[Lectura 8, ejercicio 2]\label{ej:v2-l08-2}
Si $s,t$ son simples no negativas y $s\le t$ casi en todas partes, pruebe $\int s\le\int t$.
\end{ejercicio}
\begin{solucion}
El conjunto medible $N=\{s>t\}$ es nulo. Las funciones $s\ind{X\setminus N}$ y $t\ind{X\setminus N}$ satisfacen la desigualdad en todo punto. Sus integrales coinciden con las de $s,t$ porque las partes sobre $N$ tienen integral cero. Aplicar monotonía puntual concluye.
\end{solucion}
\begin{ejercicio}[Lectura 8, ejercicio 3]\label{ej:v2-l08-3}
Si $f\ge0$ es medible y $\int f<\infty$, pruebe: (i) $\{f>0\}$ es $\sigma$-finito; (ii) $\{f=\infty\}$ es nulo.
\end{ejercicio}
\begin{solucion}
(i) Para $E_n=\{f\ge1/n\}$, la simple $n^{-1}\ind{E_n}\le f$ da $\mu(E_n)/n\le\int f$, de donde $\mu(E_n)<\infty$. Los $E_n$ crecen y su unión es $\{f>0\}$. El ejercicio \ref{ej:v2-l02-8}, aplicado a ese espacio restringido, prueba $\sigma$-finitud.

(ii) Si $Z=\{f=\infty\}$, para cada entero $m>0$ se tiene $m\ind Z\le f$, luego $m\mu(Z)\le\int f$. Por tanto $\mu(Z)\le m^{-1}\int f$ para todo $m$, y hacer $m\to\infty$ da cero. La integrabilidad permite infinitos puntuales, pero solamente en un conjunto nulo.
\end{solucion}
\begin{ejercicio}[Lectura 8, ejercicio 4]\label{ej:v2-l08-4}
En $\N$ con la medida de conteo, pruebe que $\int_\N f\,d\mu_c=\sum_na_n$ para $f(n)=a_n\ge0$.
\end{ejercicio}
\begin{solucion}
Para cada $m$, la función $s_m=\sum_{n=1}^ma_n\ind{\{n\}}$ es simple y menor o igual que $f$ si los $a_n$ son finitos. Por ello $\int f\ge\sum_{n\le m}a_n$ y, tomando supremo, $\int f\ge\sum_na_n$. Para la otra desigualdad, si $s=\sum_{j=1}^rb_j\ind{B_j}\le f$ es simple no negativa, su integral es $\sum_jb_j\mu_c(B_j)=\sum_ns(n)\le\sum_na_n$: la igualdad intermedia se comprueba agrupando los términos según los niveles de $s$, con suma infinita cuando un nivel positivo tiene infinitos puntos. Tomar supremo sobre $s$ da la desigualdad inversa. Si algún $a_n=\infty$, las simples $M\ind{\{n\}}\le f$ para todo $M$ fuerzan integral infinita, de acuerdo con la serie.
\end{solucion}

\chapter{Integrabilidad y teoremas de convergencia}\label{cap:v2-05}
Para integrar funciones con signo necesitamos evitar la cancelación de dos infinitos. Para intercambiar límites e integrales necesitamos controlar de qué manera la masa puede concentrarse o escapar. Este capítulo desarrolla las lecturas 9, 10 y 12 de Valencia Bucio \cite{valencia2}; el diagnóstico de pruebas falsas de convergencia monótona sigue la propuesta didáctica de Maximenko \cite{maximenkoFalsas}.

\section{Funciones reales integrables}\label{lectura:v2-09}
\begin{definicion}\index{parte positiva}\index{parte negativa}\index{funcion@función!integrable}
Para una función real medible $f$, definimos $f^+=\max(f,0)$ y $f^-=\max(-f,0)$. Entonces $f=f^+-f^-$ y $|f|=f^++f^-$. Decimos que $f$ es integrable si $\int|f|<\infty$, y ponemos
\[\int f=\int f^+-\int f^-.\]
Esta resta está definida porque ambos términos son finitos. Si solamente uno de ellos es finito, puede definirse una integral extendida, pero $f$ no se llama integrable.
\end{definicion}
\begin{lema}\label{lema:v2-descomposicion-integral}
Si $u,v,r,s$ son no negativas integrables y $u-v=r-s$, entonces $\int u-\int v=\int r-\int s$.
\end{lema}
\begin{proof}
La igualdad puntual equivale a $u+s=r+v$. La aditividad para funciones no negativas da $\int u+\int s=\int r+\int v$. Como las cuatro integrales son finitas, podemos reordenar y restar.
\end{proof}
\begin{teorema}\label{teo:v2-integral-lineal}
Las funciones reales integrables forman un espacio vectorial. La integral es lineal, monótona y satisface $|\int f|\le\int|f|$. Funciones medibles iguales casi en todas partes son simultáneamente integrables y tienen la misma integral.
\end{teorema}
\begin{proof}
La desigualdad $|af+bg|\le|a||f|+|b||g|$ prueba integrabilidad. Para aditividad escriba $f+g=(f^++g^+)-(f^-+g^-)$ y use el lema anterior para comparar esta descomposición con las partes positiva y negativa de $f+g$. Luego aplique aditividad no negativa a cada sumando. Para $a>0$, $(af)^+=af^+$ y $(af)^-=af^-$; para $a<0$, $(af)^+=(-a)f^-$ y $(af)^-=(-a)f^+$; para $a=0$ la conclusión es inmediata. Esto prueba homogeneidad en todos los casos.

Si $f\le g$, $g-f\ge0$ y linealidad da $\int g-\int f=\int(g-f)\ge0$. Como $- |f|\le f\le|f|$, se obtiene la desigualdad del valor absoluto. Finalmente, la igualdad casi en todas partes de $f,g$ implica la de sus partes positivas, negativas y valores absolutos. La proposición \ref{prop:v2-integral-positiva} iguala las integrales correspondientes, lo que prueba simultáneamente la equivalencia de integrabilidad y la igualdad de integrales.
\end{proof}
Para $A$ medible escribimos $\int_Af\,d\mu=\int_Xf\ind A\,d\mu$. La integral en un conjunto se entiende, por tanto, como una integral en todo el espacio de una función anulada fuera de ese conjunto.

\begin{proposicion}[Markov y anulación de la integral]\label{prop:v2-markov}\index{Markov, desigualdad de}
Si $f\ge0$ es medible y $t>0$, entonces
\[t\mu\{f\ge t\}\le\int_{\{f\ge t\}}f\,d\mu\le\int_Xf\,d\mu.\]
En particular, $\int f=0$ si y solo si $f=0$ casi en todas partes.
\end{proposicion}
\begin{proof}
Integre la desigualdad $t\ind{\{f\ge t\}}\le f\ind{\{f\ge t\}}\le f$. Si $\int f=0$, cada conjunto $\{f\ge1/n\}$ es nulo. Su unión es $\{f>0\}$ y es nula por subaditividad. La recíproca es la invariancia de la integral al modificar una función medible en un conjunto nulo.
\end{proof}

\section{Convergencia monótona y el papel de los cuantificadores}\label{lectura:v2-10}
\begin{teorema}[Convergencia monótona]\label{teo:v2-mct}\index{convergencia!monotona@monótona}
Si $f_n\ge0$ son medibles y $f_n\uparrow f$ casi en todas partes, con $f$ medible, entonces $\int f_n\uparrow\int f$.
\end{teorema}
\begin{proof}
Primero suponga convergencia y monotonía en todo punto. La monotonía de la integral da $L=\lim_n\int f_n\le\int f$. Fije una simple no negativa $s\le f$ y un número $c\in(0,1)$. Los conjuntos
\[E_n=\{f_n\ge cs\}\]
son medibles, crecen y cubren $X$: donde $s=0$ todos los puntos pertenecen a todos los $E_n$; donde $s>0$, el límite $f\ge s>cs$ obliga a pertenecer desde algún índice. Como $cs\ind{E_n}\le f_n$, obtenemos $\int f_n\ge c\int s\ind{E_n}$. Escribir $s$ como suma finita de indicadores y usar continuidad por abajo de $\mu$ da $\int s\ind{E_n}\uparrow\int s$. Luego $L\ge c\int s$. Si $\int s=\infty$, esto ya fuerza $L=\infty$; si es finita, hacer $c\uparrow1$ da $L\ge\int s$. Finalmente tome el supremo sobre todas las simples $s\le f$.

En el caso casi en todas partes, reúna en un único conjunto medible nulo todos los puntos donde falla la monotonía de algún par consecutivo o donde falla la convergencia. Anule $f_n$ y $f$ en ese conjunto. Las funciones obtenidas satisfacen el caso puntual y conservan sus integrales.
\end{proof}
\begin{observacion}[Dos pasos falsos y su reparación]
De $f_n\uparrow f$ no se deduce que algún $f_n$ alcance a $f$: las constantes $f_n=1-1/n$ nunca alcanzan el límite $1$. Tampoco de $\forall x\ \exists n:\ f_n(x)\ge cs(x)$ se deduce $\exists n\ \forall x:\ f_n(x)\ge cs(x)$. Por ejemplo, en $\N$, $f_n=\ind{\{1,\ldots,n\}}$ crece a $s=1$, pero ningún índice sirve para todos los puntos. La prueba correcta no busca ese índice común: utiliza los conjuntos crecientes $E_n$ y la continuidad de la medida. Este diagnóstico se inspira expresamente en \cite{maximenkoFalsas}; las pruebas deliberadamente falsas de esa presentación no se utilizan como demostraciones.
\end{observacion}
\begin{observacion}[La diagonal del argumento original]
También puede demostrarse el teorema tomando, para cada $j$, simples $s_{j,k}\uparrow f_j$ y definiendo $u_n=\max_{j\le n}s_{j,n}$. Se tiene $u_n\le f_n$, y $u_{n+1}\ge u_n$ porque cada fila es creciente y se añade una fila. Para cada $j$ fijo y $n\ge j$, $u_n\ge s_{j,n}$, luego $\lim_nu_n\ge f_j$; tomar supremo en $j$ da $\lim_nu_n\ge f$. La desigualdad inversa resulta de $u_n\le f_n\le f$. Así el teorema \ref{teo:v2-integral-aprox} también proporciona la conclusión y quedan explicadas las dependencias de índices de la lectura 10.
\end{observacion}

\begin{teorema}[Fatou]\label{teo:v2-fatou}\index{Fatou, lema de}
Si $f_n\ge0$ son medibles, entonces $\int\liminf_nf_n\le\liminf_n\int f_n$.
\end{teorema}
\begin{proof}
Defina $g_m=\inf_{n\ge m}f_n$. Son medibles, no negativas y crecen a $\liminf_nf_n$. Para todo $n\ge m$ se tiene $g_m\le f_n$, luego $\int g_m\le\inf_{n\ge m}\int f_n$. Convergencia monótona en el lado izquierdo y la definición de límite inferior en el derecho dan el resultado.
\end{proof}
\begin{teorema}[Convergencia dominada]\label{teo:v2-dct}\index{convergencia!dominada}
Sean $f_n$ medibles reales con $f_n\to f$ casi en todas partes, $f$ medible. Si existe $g\ge0$ integrable tal que $|f_n|\le g$ casi en todas partes para cada $n$, entonces $f_n,f$ son integrables,
\[\int|f_n-f|\longrightarrow0,\qquad\int f_n\longrightarrow\int f.\]
\end{teorema}
\begin{proof}
Fuera de la unión numerable de los conjuntos excepcionales, $|f|\le g$ y $|f_n-f|\le2g$. Por monotonía todas esas funciones son integrables. Aplique Fatou a $2g-|f_n-f|$, anulando previamente todas las funciones en el conjunto excepcional y donde $g$ es infinita. Su límite es $2g$, por lo que
\[2\int g\le\liminf_n\left(2\int g-\int|f_n-f|\right)
=2\int g-\limsup_n\int|f_n-f|.\]
Como $\int g$ es finita, la resta es legítima y obtenemos $\limsup_n\int|f_n-f|\le0$. La no negatividad implica convergencia a cero. Finalmente $|\int f_n-\int f|\le\int|f_n-f|$ prueba la convergencia de las integrales. La función dominante es una sola $g$; el subíndice accidental $g_n$ en el enunciado original no debe interpretarse como una dominación variable sin control adicional.
\end{proof}

\section{Funciones complejas e integración}\label{lectura:v2-12}
\begin{proposicion}\index{funcion@función!compleja medible}
Para $f:X\to\C$, son equivalentes: sus preimágenes de abiertos son medibles; sus preimágenes de borelianos son medibles; sus partes real e imaginaria son medibles.
\end{proposicion}
\begin{proof}
La familia de subconjuntos de $\C$ cuya preimagen es medible es una $\sigma$-álgebra; si contiene los abiertos, contiene los borelianos. Los niveles de las partes real e imaginaria son preimágenes de semiplanos abiertos, por lo que la primera condición implica la tercera. Recíprocamente, la preimagen de un rectángulo abierto $(a,b)+i(c,d)$ es $\{a<\Re f<b\}\cap\{c<\Im f<d\}$. Los rectángulos de extremos racionales forman una base numerable de $\C$; todo abierto es unión numerable de ellos. Esto prueba la tercera implicación y cierra las equivalencias.
\end{proof}
Las combinaciones lineales complejas y los productos de funciones medibles son medibles: sus partes reales e imaginarias son sumas y productos de funciones reales medibles. También lo es $|f|=\sqrt{(\Re f)^2+(\Im f)^2}$, pues la raíz cuadrada es continua.

\begin{definicion}
Una función compleja medible es integrable si sus partes real e imaginaria lo son, y se define $\int f=\int\Re f+i\int\Im f$.
\end{definicion}
\begin{teorema}\label{teo:v2-complejas}
La condición anterior equivale a $\int|f|<\infty$. La integral compleja es lineal y satisface $|\int f|\le\int|f|$. El teorema de convergencia dominada vale también para funciones complejas, incluyendo convergencia en integral del error absoluto.
\end{teorema}
\begin{proof}
Use $|\Re f|,|\Im f|\le|f|\le|\Re f|+|\Im f|$ para la equivalencia de integrabilidad. La aditividad se verifica por componentes. Si $\alpha=a+ib$ y $f=u+iv$, entonces $\alpha f=(au-bv)+i(av+bu)$; la linealidad real aplicada a las dos componentes da $\int\alpha f=\alpha\int f$.

Si $z=\int f=0$, la desigualdad del valor absoluto es inmediata. Si $z\ne0$, elija $\omega=\overline z/|z|$, de modo que $|\omega|=1$ y $\omega z=|z|$. Por linealidad,
\[\left|\int f\right|=\Re\int\omega f=\int\Re(\omega f)\le\int|f|.\]
Para convergencia dominada, $\Re f_n\to\Re f$ e $\Im f_n\to\Im f$ casi en todas partes y ambas sucesiones están dominadas en valor absoluto por $g$. El teorema real da convergencia de sus integrales. Además $|f_n-f|\to0$ y está dominada por $2g$; aplicarlo a esta sucesión real no negativa prueba la convergencia del error absoluto.
\end{proof}
El teorema de Fubini complejo de la lectura 12 se demostrará junto con su versión real en el capítulo \ref{cap:v2-06}; no se necesita para los resultados de este capítulo.

\section{Ejercicios sobre funciones integrables}
\begin{ejercicio}[Lectura 9, ejercicio 1]\label{ej:v2-l09-1}
Pruebe: (i) linealidad de la integral sobre $A$; (ii) integrabilidad sobre $A\cup B$ y $A\cap B$, e inclusión--exclusión, si $f$ es integrable sobre $A$ y $B$; (iii) monotonía de la integral sobre $A$.
\end{ejercicio}

\begin{ejercicio}[Lectura 9, ejercicio 2]\label{ej:v2-l09-2}
Construya una función no medible cuyo valor absoluto sea integrable.
\end{ejercicio}

\begin{ejercicio}[Lectura 9, ejercicio 3]\label{ej:v2-l09-3}
Demuestre la desigualdad de Markov y deduzca: (i) una función integrable es finita casi en todas partes; (ii) $\int|f|=0$ implica $f=0$ casi en todas partes; (iii) si $f$ es real integrable y $\int_Af\ge0$ para todo $A$ medible, entonces $f\ge0$ casi en todas partes.
\end{ejercicio}

\begin{ejercicio}[Lectura 9, ejercicio 4]\label{ej:v2-l09-4}
Dé dos funciones integrables cuyo producto no sea integrable.
\end{ejercicio}

\begin{ejercicio}[Lectura 9, ejercicio 5]\label{ej:v2-l09-5}
Si $f$ es integrable, $g$ es medible y $f\le g$ casi en todas partes, pruebe que la integral extendida de $g$ existe y $\int f\le\int g$.
\end{ejercicio}

\begin{ejercicio}[Lectura 9, ejercicio 6]\label{ej:v2-l09-6}\index{Riemann y Lebesgue, comparacion@Riemann y Lebesgue, comparación}
Si $f:[a,b]\to\R$ es Riemann-integrable, pruebe que es Lebesgue-integrable y ambas integrales coinciden.
\end{ejercicio}


\section{Ejercicios sobre intercambio de límites e integrales}
\begin{ejercicio}[Lectura 10, ejercicio 1: Beppo Levi]\label{ej:v2-l10-1}\index{Beppo Levi, teorema de}
Si $f_n\ge0$ son medibles, pruebe $\int\sum_nf_n=\sum_n\int f_n$.
\end{ejercicio}

\begin{ejercicio}[Lectura 10, ejercicio 2]\label{ej:v2-l10-2}
Sea $A\subseteq\R$ y una familia $(f_t)_{t\in A}$ de funciones integrables tal que $f_t(x)\to f_a(x)$ cuando $t\to a$ en $A$, para casi todo $x$ y cada $a\in A$. Si $|f_t|\le g$ casi en todas partes para una sola $g\in L^1$, pruebe la continuidad de $t\mapsto\int f_t$.
\end{ejercicio}

\begin{ejercicio}[Lectura 10, ejercicio 3: transformada de Laplace]\label{ej:v2-l10-3}\index{Laplace, transformada de}
Para $f$ Lebesgue-medible en $[0,\infty)$, defina $Lf(s)=\int_0^\infty e^{-sx}f(x)\,dx$ cuando la función es integrable. Pruebe: (i) su dominio es estable hacia la derecha; (ii) $Lf$ es diferenciable en el interior del dominio; (iii) si está definida en todo $\R$, $Lf(s)\to0$ cuando $s\to\infty$; (iv) aditividad.
\end{ejercicio}

\begin{ejercicio}[Lectura 10, ejercicio 4]\label{ej:v2-l10-4}
Si $|f_n|\le g$ casi en todas partes con $g$ integrable, pruebe integrabilidad de $\limsup f_n$ y $\int\limsup f_n\ge\limsup\int f_n$.
\end{ejercicio}

\begin{ejercicio}[Lectura 10, ejercicio 5]\label{ej:v2-l10-5}
Para $f_n=-n^{-1}\ind{[0,n]}$ en $\R$, pruebe convergencia uniforme a cero y fallo de la desigualdad de Fatou sin no negatividad.
\end{ejercicio}

\begin{ejercicio}[Lectura 10, ejercicio 6]\label{ej:v2-l10-6}
Si $f_n\ge-h$ con $h\ge0$ integrable, pruebe $\int\liminf f_n\le\liminf\int f_n$. Explique por qué no se aplica al ejercicio anterior.
\end{ejercicio}

\begin{ejercicio}[Lectura 10, ejercicio 7]\label{ej:v2-l10-7}
Extienda convergencia monótona a una sucesión creciente cuyo primer término es integrable, aunque no sea no negativo.
\end{ejercicio}

\begin{ejercicio}[Lectura 10, ejercicio 8: núcleos]\label{ej:v2-l10-8}\index{nucleo de medidas@núcleo de medidas}
Sea $K:X\times\Bcal\to[0,\infty]$ tal que $K(x,\cdot)$ es una medida en $(Y,\Bcal)$ y $K(\cdot,B)$ es medible para cada $B\in\Bcal$. Si $\mu$ es una medida en $X$ y $f:Y\to[0,\infty]$ es medible, pruebe: (i) $\nu(B)=\int_XK(x,B)\,d\mu(x)$ es una medida; (ii) $x\mapsto\int_Yf(y)K(x,dy)$ es medible; (iii) su integral en $X$ es $\int_Yf\,d\nu$.
\end{ejercicio}

\begin{ejercicio}[Lectura 10, ejercicio 9]\label{ej:v2-l10-9}
En el ejercicio anterior suponga $\mu(X)<\infty$, $M=\sup_xK(x,Y)<\infty$ y $f$ medible acotada real. Pruebe medibilidad y acotación de $x\mapsto\int_Yf\,K(x,dy)$ y la misma identidad de integrales.
\end{ejercicio}


\section{Una aplicación: derivación bajo la integral}
\begin{teorema}\label{teo:v2-derivar-integral}\index{derivacion bajo el signo integral@derivación bajo el signo integral}
Sea $F(t)=\int_Xh(x,t)\,d\mu(x)$ para $t$ en un intervalo abierto $I$. Suponga que cada $h(\cdot,t)$ es integrable, que $\partial_th(x,t)$ existe para todo $(x,t)$ y que $|\partial_th(x,t)|\le g(x)$ para una $g\in L^1$, $g\ge0$. Entonces $F$ es diferenciable y $F'(t)=\int_X\partial_th(x,t)\,d\mu(x)$.
\end{teorema}
\begin{proof}
Fije $t\in I$ y una sucesión $t_n\to t$, $t_n\ne t$. Los cocientes $q_n(x)=(h(x,t_n)-h(x,t))/(t_n-t)$ son medibles por las hipótesis de integrabilidad. Convergen puntualmente a $\partial_th(x,t)$, que por ello es medible. El teorema del valor medio aplicado a $s\mapsto h(x,s)$ da $|q_n(x)|\le g(x)$. No hace falta que el punto intermedio elegido por ese teorema dependa mediblemente de $x$, porque solamente usamos la cota. Convergencia dominada prueba
\[\frac{F(t_n)-F(t)}{t_n-t}=\int_Xq_n\,d\mu\longrightarrow\int_X\partial_th(x,t)\,d\mu.\]
El límite es el mismo para toda sucesión $t_n\to t$; el criterio secuencial para límites reales prueba la diferenciabilidad y la fórmula.
\end{proof}

\chapter{Medidas producto, Tonelli y Fubini}\label{cap:v2-06}
Una integral iterada no es, por definición, una integral en un producto. Su igualdad requiere demostrar medibilidad de las secciones y justificar el intercambio de operaciones. Desarrollamos la lectura 11 y el teorema complejo de la lectura 12 de Valencia Bucio \cite{valencia2}.

\section{Rectángulos, secciones y clases monótonas}\label{lectura:v2-11}
\begin{definicion}\index{medida!producto}\index{seccion de un conjunto@sección de un conjunto}
Para espacios medibles $(X,\Acal)$ y $(Y,\Bcal)$, la $\sigma$-álgebra producto es $\Acal\otimes\Bcal=\sigma\{A\times B:A\in\Acal,B\in\Bcal\}$. Para $E\subseteq X\times Y$ definimos $E_x=\{y:(x,y)\in E\}$ y $E^y=\{x:(x,y)\in E\}$. Para una función $F$ se definen análogamente $F_x(y)=F(x,y)$ y $F^y(x)=F(x,y)$.
\end{definicion}
\begin{lema}\label{lema:v2-secciones}
Las secciones de conjuntos medibles en la $\sigma$-álgebra producto son medibles. Las secciones de funciones medibles son medibles.
\end{lema}
\begin{proof}
Para $x$ fijo, la familia de conjuntos $E$ con $E_x\in\Bcal$ contiene los rectángulos, pues $(A\times B)_x$ es $B$ si $x\in A$ y vacío en otro caso. Además $(E^c)_x=(E_x)^c$ y $(\bigcup_nE_n)_x=\bigcup_n(E_n)_x$; por ello esta familia es una $\sigma$-álgebra y contiene $\Acal\otimes\Bcal$. Para funciones reales o extendidas, $\{F_x>t\}=\{F>t\}_x$; para funciones complejas se aplican las partes real e imaginaria. El argumento para las otras secciones es idéntico intercambiando los factores.
\end{proof}
\begin{definicion}\index{clase monotona@clase monótona}
Una clase monótona es una familia no vacía estable bajo uniones crecientes numerables e intersecciones decrecientes numerables. La menor clase monótona que contiene una familia $\mathcal G$ se denota $M(\mathcal G)$.
\end{definicion}
\begin{teorema}[Clase monótona]\label{teo:v2-clasemonotona}
Si $\mathcal R$ es un álgebra, entonces $M(\mathcal R)=\sigma(\mathcal R)$.
\end{teorema}
\begin{proof}
Escriba $\mathcal M=M(\mathcal R)$. Toda $\sigma$-álgebra es una clase monótona, así que $\mathcal M\subseteq\sigma(\mathcal R)$. La familia $\{E\in\mathcal M:E^c\in\mathcal M\}$ es una clase monótona: los complementos intercambian uniones crecientes e intersecciones decrecientes. Contiene $\mathcal R$, luego por minimalidad coincide con $\mathcal M$. Así $\mathcal M$ es estable bajo complementos.

Fije $A\in\mathcal R$. La familia $\{E\in\mathcal M:E\cap A\in\mathcal M\}$ es monótona y contiene $\mathcal R$, por lo que es toda $\mathcal M$. Ahora fije $E\in\mathcal M$. La familia $\{F\in\mathcal M:E\cap F\in\mathcal M\}$ es monótona y, por el paso anterior, contiene $\mathcal R$; nuevamente coincide con $\mathcal M$. Hemos probado estabilidad bajo intersecciones binarias. Como $X\in\mathcal R$, $\mathcal M$ es un álgebra. Una unión numerable de sus elementos es el límite creciente de sus uniones finitas, luego también pertenece a $\mathcal M$. Por tanto $\mathcal M$ es una $\sigma$-álgebra que contiene $\mathcal R$, y la minimalidad de $\sigma(\mathcal R)$ da la inclusión restante.
\end{proof}

\section{Construcción de la medida producto}
\begin{proposicion}\label{prop:v2-producto}
Existe una medida $\pi$ en $\Acal\otimes\Bcal$ con $\pi(A\times B)=\mu(A)\nu(B)$. Si $\mu,\nu$ son $\sigma$-finitas, esta medida es única y también es $\sigma$-finita.
\end{proposicion}
\begin{proof}
En el álgebra de uniones finitas de rectángulos, defina $v(E)=\sum_j\mu(A_j)\nu(B_j)$ usando una partición rectangular disjunta de $E$. Una partición común en ambos factores prueba independencia de representación y aditividad finita, como en el capítulo \ref{cap:v2-03}. También
\[v(E)=\int_X\nu(E_x)\,d\mu(x),\]
pues $\nu(E_x)=\sum_j\ind{A_j}(x)\nu(B_j)$; si algún coeficiente es infinito, la igualdad se obtiene truncándolo y aplicando convergencia monótona.

Si $E=\bigcup_nE_n$ es disjunta y todos esos conjuntos pertenecen al álgebra, sus secciones son disjuntas y $\nu(E_x)=\sum_n\nu((E_n)_x)$. Beppo Levi da $v(E)=\sum_nv(E_n)$. Así $v$ es premedida y el teorema \ref{teo:v2-extension} produce la medida producto.

Para unicidad, primero suponga finitos los dos factores. Dos medidas que coinciden en rectángulos coinciden en el álgebra de sus uniones finitas disjuntas. La familia donde coinciden es una clase monótona: para uniones crecientes se usa continuidad por abajo, y para intersecciones decrecientes se usa finitud de la medida total. El teorema \ref{teo:v2-clasemonotona} da igualdad en toda la $\sigma$-álgebra producto. Si los factores son $\sigma$-finitos, elija $X_n\uparrow X$, $Y_n\uparrow Y$ de medida finita. Las restricciones a $X_n\times Y_n$ coinciden por el caso finito; continuidad por abajo sobre esos rectángulos prueba unicidad global. Esos mismos rectángulos muestran $\sigma$-finitud de $\pi$.
\end{proof}

\begin{lema}\label{lema:v2-cavalieri}
Si los factores son $\sigma$-finitos y $E\in\Acal\otimes\Bcal$, las funciones $x\mapsto\nu(E_x)$, $y\mapsto\mu(E^y)$ son medibles y
\[\pi(E)=\int_X\nu(E_x)\,d\mu(x)=\int_Y\mu(E^y)\,d\nu(y).\]
\end{lema}
\begin{proof}
Suponga primero $\mu(X),\nu(Y)<\infty$. La propiedad vale en rectángulos y, por aditividad, en su álgebra. La familia de conjuntos que la satisfacen es una clase monótona. En efecto, para $E_n\uparrow E$, las secciones crecen y continuidad por abajo de $\nu$ implica $\nu((E_n)_x)\uparrow\nu(E_x)$; los límites son medibles y convergencia monótona identifica sus integrales con $\lim_n\pi(E_n)=\pi(E)$. Si $E_n\downarrow E$, continuidad por arriba se aplica a las secciones porque $\nu(Y)<\infty$; las funciones están dominadas por la constante integrable $\nu(Y)$ en el espacio finito $X$. Convergencia dominada prueba la misma igualdad. La otra dirección de integración se trata igual. El teorema de clase monótona concluye el caso finito.

En el caso $\sigma$-finito, use los rectángulos $X_n\times Y_n$ de la prueba anterior. Las funciones
\[h_n(x)=\ind{X_n}(x)\nu(E_x\cap Y_n)\]
son medibles por el caso finito, extendidas por cero, y crecen a $\nu(E_x)$. Su integral vale $\pi(E\cap(X_n\times Y_n))$. Convergencia monótona y continuidad por abajo de $\pi$ dan la primera igualdad global; intercambiar factores da la segunda.
\end{proof}

\section{Integrar por secciones}
A partir de esta sección y en los ejercicios sobre integrales iteradas suponemos que ambos factores son $\sigma$-finitos. El ejercicio \ref{ej:v2-l11-5} abandona expresamente esta hipótesis para mostrar por qué no puede omitirse. La construcción de la medida producto de la sección anterior no necesitaba esa suposición para su existencia.
\begin{teorema}[Tonelli]\label{teo:v2-tonelli}\index{Tonelli, teorema de}
Sean $\mu,\nu$ $\sigma$-finitas. Si $F:X\times Y\to[0,\infty]$ es medible, las integrales de sus secciones son funciones medibles y
\[\int_{X\times Y}F\,d\pi=\int_X\left(\int_YF(x,y)\,d\nu(y)\right)d\mu(x)
=\int_Y\left(\int_XF(x,y)\,d\mu(x)\right)d\nu(y).\]
Se permiten valores infinitos.
\end{teorema}
\begin{proof}
Para indicadores es el lema \ref{lema:v2-cavalieri}; para simples no negativas es su combinación lineal finita. Tome simples $S_n\uparrow F$. Para cada $x$, convergencia monótona en $Y$ da $\int(S_n)_x\,d\nu\uparrow\int F_x\,d\nu$. El límite es medible. Convergencia monótona en $X$ y en $X\times Y$ da
\[\int_X\int_YF\,d\nu d\mu=\lim_n\int_X\int_YS_n\,d\nu d\mu
=\lim_n\int_{X\times Y}S_n\,d\pi=\int_{X\times Y}F\,d\pi.\]
El argumento con los factores invertidos prueba la otra igualdad.
\end{proof}
\begin{teorema}[Fubini real y complejo]\label{teo:v2-fubini}\index{Fubini, teorema de}
Sean $\mu,\nu$ medidas $\sigma$-finitas. Si $F:X\times Y\to\F$ es medible y $\int|F|\,d\pi<\infty$, entonces sus secciones son integrables casi en todas partes en el factor correspondiente. Las integrales interiores, definidas como cero en los conjuntos excepcionales, son integrables y satisfacen las igualdades de Tonelli.
\end{teorema}
\begin{proof}
Tonelli aplicado a $|F|$ muestra que $H(x)=\int_Y|F_x|\,d\nu$ es medible y tiene integral finita. Por ello $H(x)<\infty$ fuera de un conjunto medible nulo: allí $F_x$ es integrable. Si $F$ es real, aplique Tonelli a $F^+,F^-$ y reste; ambas integrales globales son finitas, así que las restas son legítimas. La función interior $f(x)=\int F_x\,d\nu$ es medible al ser la diferencia de las integrales positivas en la zona donde son finitas, extendida por cero en su complemento nulo. Además $|f(x)|\le H(x)$, luego es integrable. Para $F$ compleja se aplica lo anterior a $\Re F$ e $\Im F$ y se combinan sus integrales usando linealidad compleja. Intercambiar los factores completa la conclusión.
\end{proof}
\begin{observacion}[Productos completados]
Los enunciados anteriores se refieren a la $\sigma$-álgebra producto. Si se completa la medida producto, una función medible para la completación tiene un representante medible para el producto e igual casi en todas partes. Esto se comprueba primero para conjuntos, escribiendo cada conjunto de la completación como un medible más una modificación dentro de un nulo, y después para funciones mediante aproximaciones simples. Tonelli aplicado al nulo que contiene la discrepancia muestra que sus secciones son nulas para casi todo punto del otro factor. Por ello las conclusiones de Fubini se conservan, con secciones entendidas casi en todas partes y con medidas de los factores completadas. No debe afirmarse que todas las secciones de un conjunto de la completación son necesariamente medibles en las $\sigma$-álgebras originales.
\end{observacion}

\section{Ejercicios sobre productos e integrales iteradas}
\begin{ejercicio}[Lectura 11, ejercicio 1]\label{ej:v2-l11-1}
Pruebe que una intersección no vacía de clases monótonas es una clase monótona.
\end{ejercicio}

\begin{ejercicio}[Lectura 11, ejercicio 2: Tonelli--Hobson]\label{ej:v2-l11-2}
Pruebe que una función \emph{medible} $F$ en el producto es integrable si y solo si una de las integrales iteradas de $|F|$ es finita.
\end{ejercicio}

\begin{ejercicio}[Lectura 11, ejercicio 3]\label{ej:v2-l11-3}
Para $a_{m,n}=1$ si $m=n$, $a_{m,n}=-1$ si $m+1=n$ y cero en los demás casos, calcule las dos sumas iteradas.
\end{ejercicio}

\begin{ejercicio}[Lectura 11, ejercicio 4]\label{ej:v2-l11-4}
Si $f\in L^1(X)$ y $g\in L^1(Y)$, pruebe integrabilidad de $h(x,y)=f(x)g(y)$ y $\int h=(\int f)(\int g)$.
\end{ejercicio}

\begin{ejercicio}[Lectura 11, ejercicio 5]\label{ej:v2-l11-5}
En $X=Y=[0,1]$ con borelianos, tome Lebesgue en $X$ y conteo en $Y$. Para la diagonal $D=\{(x,x):x\in[0,1]\}$, pruebe medibilidad y desigualdad de sus integrales iteradas por secciones.
\end{ejercicio}

\begin{ejercicio}[Lectura 11, ejercicio 6: convolución]\label{ej:v2-l11-6}\index{convolucion@convolución}
Sean $f,g\in L^1(\R)$. Pruebe: (i) $(x,y)\mapsto f(x-y)g(y)$ es Lebesgue-medible en $\R^2$; (ii) $(f*g)(x)=\int_\R f(x-y)g(y)\,dy$ existe casi en todas partes; (iii) $\|f*g\|_1\le\|f\|_1\|g\|_1$.
\end{ejercicio}

\begin{observacion}[Conmutatividad y asociatividad]
Conmutatividad resulta del cambio $y\mapsto x-y$ en la integral; esta sustitución preserva Lebesgue por la invariancia bajo reflexión y traslación, probada primero en indicadores y extendida por aproximación simple. Para asociatividad, la integral del valor absoluto de $f(x-y-z)g(y)h(z)$ en $\R^3$ es $\|f\|_1\|g\|_1\|h\|_1$ por Tonelli. Fubini permite cambiar el orden y efectuar las traslaciones interiores, dando $f*(g*h)=(f*g)*h$ casi en todas partes. La estimación triple justifica el intercambio; las igualdades son de clases $L^1$ y no necesariamente en cada punto de representantes arbitrarios.
\end{observacion}
\begin{ejercicio}[Lectura 11, ejercicio 7]\label{ej:v2-l11-7}
Si $A,B\in\Acal\otimes\Bcal$ y $\nu(A_x)=\nu(B_x)$ para casi todo $x$, pruebe $\pi(A)=\pi(B)$.
\end{ejercicio}

\parte{Convergencia y aproximación}{Compararemos convergencia puntual, casi uniforme, en medida y en media. Los contraejemplos separan concentración, dispersión y desplazamiento de la masa. La construcción de $L^1$ convierte estos problemas en preguntas de distancia, completitud y densidad; sus prerrequisitos son los teoremas de convergencia de la parte anterior.}
\chapter{Modos de convergencia de funciones}\label{cap:v2-07}

Una sucesión de funciones puede aproximar bien a su límite en casi todos los puntos y, sin embargo, conservar errores grandes sobre conjuntos que se desplazan. Otra puede tener errores pequeños salvo en conjuntos de medida decreciente, pero visitar cada punto infinitas veces. Este capítulo separa esas situaciones. Las lecturas 14 y 15 de Valencia \cite{valencia2} proporcionan el hilo conductor; los ejemplos adicionales siguen la orientación de Maximenko \cite{maximenkoConvergencias}.

\section{Convergencia en medida y casi en todas partes}\label{lectura:v2-14}
Trabajamos con funciones medibles finitas, reales o complejas, en $(X,\Acal,\mu)$. Cuando una función es finita casi en todas partes, elegimos un representante finito modificándola en su conjunto excepcional medible.
\begin{definicion}
Se dice que $f_n\to f$ \emph{en medida}\index{convergencia!en medida} si para cada $a>0$,
\[\mu\{x:|f_n(x)-f(x)|>a\}\longrightarrow0.\]
Se dice que converge \emph{casi uniformemente}\index{convergencia!casi uniforme} si para cada $\eps>0$ existe $A\in\Acal$ tal que $\mu(X\setminus A)<\eps$ y la convergencia restringida a $A$ es uniforme.
\end{definicion}
En la primera definición puede reemplazarse $>a$ por $\ge a$: basta comparar los umbrales $a/2$ y $a$. El conjunto excepcional de la convergencia casi uniforme depende de $\eps$, pero, una vez elegido, debe servir para \emph{todos} los umbrales de error de la convergencia uniforme.

\begin{ejemplo}[Dos obstáculos distintos]
En $\N$ con medida de conteo, $f_n(k)=\ind{\{n,n+1,\ldots\}}(k)$ converge puntualmente a cero. Sin embargo, $\mu\{|f_n|>1/2\}=\infty$ para todo $n$.

En $[0,1)$ enumeremos por bloques los intervalos diádicos
\[I_{m,j}=[j2^{-m},(j+1)2^{-m}),\qquad 0\le j<2^m,\quad m=1,2,\ldots .\]
La sucesión de sus indicadores converge en medida a cero, pues las longitudes de los intervalos del bloque $m$ son $2^{-m}$. Cada punto pertenece exactamente a uno de los intervalos de cada bloque; por tanto recibe el valor uno infinitas veces y el valor cero infinitas veces. No hay convergencia puntual en ningún punto. El uso de intervalos semiabiertos evita ambigüedades en extremos diádicos.
\end{ejemplo}

\begin{teorema}\label{teo:v2-ae-medida}
Si $\mu(X)<\infty$ y $f_n\to f$ casi en todas partes, entonces $f_n\to f$ en medida.
\end{teorema}
\begin{proof}
Fijemos $a>0$ y pongamos $B_N=\bigcup_{n\ge N}\{|f_n-f|>a\}$. Los conjuntos $B_N$ decrecen; su intersección está contenida en el conjunto nulo donde falla la convergencia puntual. Como $\mu(B_1)\le\mu(X)<\infty$, la continuidad de la medida desde arriba da $\mu(B_N)\to0$. Si $n\ge N$, el conjunto malo para $f_n$ está contenido en $B_N$, y eso prueba la afirmación.
\end{proof}

\begin{teorema}[Extracción de una subsucesión]\label{teo:v2-medida-ae}
Si $f_n\to f$ en medida, existe una subsucesión que converge a $f$ casi en todas partes.
\end{teorema}
\begin{proof}
Elegimos inductivamente $n_k>n_{k-1}$ de modo que
\[\mu\{|f_{n_k}-f|>1/k\}<2^{-k}.
\]
La suma de las medidas es finita. Por Borel--Cantelli, casi ningún punto pertenece a infinitos conjuntos de esta lista. En cada punto restante existe $k_0(x)$ tal que $|f_{n_k}(x)-f(x)|\le1/k$ para $k\ge k_0(x)$; esto es precisamente convergencia puntual.
\end{proof}

\begin{teorema}[Egorov]\label{teo:v2-egorov}\index{Egorov, teorema de}
En un espacio de medida finita, la convergencia casi en todas partes de funciones medibles finitas implica convergencia casi uniforme.
\end{teorema}
\begin{proof}
Retiramos primero el conjunto nulo $N$ donde falla la convergencia. Para cada entero $k\ge1$, los conjuntos
\[E_{k,m}=\{x\notin N:\sup_{n\ge m}|f_n(x)-f(x)|>1/k\}
\]
decrecen a vacío cuando $m\to\infty$. Su medida tiende a cero por finitud de $\mu(X)$. Dado $\eps>0$, elegimos $m_k$ de manera que $\mu(E_{k,m_k})<\eps2^{-k-1}$. Definimos
\[A=X\setminus\left(N\cup\bigcup_{k\ge1}E_{k,m_k}\right).
\]
Entonces $\mu(X\setminus A)<\eps$. Para comprobar uniformidad, sea $\delta>0$ y elijamos $k$ con $1/k<\delta$. Si $n\ge m_k$, en \emph{todos} los puntos de $A$ se tiene $|f_n-f|\le1/k<\delta$. El conjunto $A$ se construyó antes de elegir este $\delta$, que es el orden correcto de los cuantificadores.
\end{proof}

\begin{ejemplo}[Concentración gaussiana]
La sucesión $g_n(x)=n e^{-n^2x^2}$ en $\R$ converge a cero para $x\ne0$. También converge casi uniformemente: dado $\eps>0$, fijemos $0<\eta<\eps/2$ y retiremos $(-\eta,\eta)$. En el complemento,
\[\sup_{|x|\ge\eta}g_n(x)=n e^{-n^2\eta^2}\longrightarrow0\qquad(n\to\infty).
\]
Por sustitución $u=nx$, su integral es la constante positiva $\int_\R e^{-u^2}\,du$, de modo que no converge a cero en media. No necesitamos evaluar esa integral para concluirlo. En el cálculo de \cite{maximenkoConvergencias}, la variable que debe tender a infinito es $n$, manteniendo $\eta$ fijo; hacer $\eta\to0$ no prueba uniformidad en el conjunto elegido.
\end{ejemplo}

\section{Ejercicios sobre conjuntos excepcionales}
\begin{ejercicio}[Lectura 14, ejercicio 1]\label{ej:v2-l14-1}
Demuestre que la convergencia uniforme implica convergencia en medida.
\end{ejercicio}
\begin{solucion}
Para un umbral $a>0$, desde cierto índice se cumple $|f_n(x)-f(x)|<a$ para todo $x$. El conjunto malo es entonces vacío, incluso cuando $\mu(X)=\infty$.
\end{solucion}
\begin{ejercicio}[Lectura 14, ejercicio 2]\label{ej:v2-l14-2}
En $\N$ con medida de conteo, demuestre que converger en medida equivale a converger uniformemente.
\end{ejercicio}
\begin{solucion}
La implicación directa del ejercicio anterior vale en cualquier espacio. Recíprocamente, para $a>0$ la medida del conjunto malo acaba siendo menor que uno. La medida de conteo sólo toma valores enteros no negativos o infinito, por lo que ese conjunto debe ser vacío. Así, $\sup_k|f_n(k)-f(k)|\le a$ desde cierto índice; usar $a=\eps/2$ da la condición uniforme con error $\eps$.
\end{solucion}
\begin{ejercicio}[Lectura 14, ejercicio 3]\label{ej:v2-l14-3}
Si $f_n\to0$ casi en todas partes en $[0,1]$, encuentre números $a_n>0$ con $a_n\to\infty$ y $a_nf_n\to0$ casi en todas partes.
\end{ejercicio}
\begin{solucion}
Para cada $k$ elegimos $N_k$ estrictamente creciente de modo que
\[m\{\sup_{n\ge N_k}|f_n|>k^{-2}\}<2^{-k}.
\]
Esto es posible por continuidad desde arriba, como en Egorov. Definimos $a_n=k$ si $N_k\le n<N_{k+1}$, y $a_n=1$ antes de $N_1$. Borel--Cantelli asegura que casi todo $x$ queda fuera de los conjuntos malos para todos los $k$ suficientemente grandes. En el bloque correspondiente, $|a_nf_n(x)|\le k/k^2=1/k$. Puesto que los índices de bloque tienden a infinito, se obtiene el resultado.
\end{solucion}
\begin{ejercicio}[Lectura 14, ejercicio 4]\label{ej:v2-l14-4}
Suponga $|f_n|\le g\in L^1(\mu)$ y $f_n\to f$ casi en todas partes. Demuestre convergencia casi uniforme sin suponer finita la medida del espacio.
\end{ejercicio}
\begin{solucion}
Retiremos un conjunto nulo común donde fallen las desigualdades o la convergencia. Allí también $|f|\le g$. Sean $E_j=\{g>1/j\}$; Markov da $\mu(E_j)\le j\int g<\infty$. Aplicamos Egorov en cada $E_j$, retirando $N_j\subset E_j$ de medida menor que $\eps2^{-j-1}$. Fuera de $N=\bigcup_jN_j$ y del conjunto nulo inicial, la convergencia es uniforme en cada $E_j$. Para un error $\delta>0$, elegimos $j$ con $2/j<\delta$. Fuera de $E_j$, $|f_n-f|\le2g\le2/j<\delta$ para todo $n$; en $E_j\setminus N$, el error es menor que $\delta$ desde un índice independiente del punto. Como $\mu(N)<\eps$, esto prueba la afirmación. La unión de excepciones se tomó una sola vez; no se cambia al escoger $\delta$.
\end{solucion}

\section{Sucesiones de Cauchy y existencia de límites}\label{lectura:v2-15}
\begin{definicion}
Una sucesión es \emph{de Cauchy en medida}\index{Cauchy!en medida} si, para cada $a>0$,
\[\lim_{m,n\to\infty}\mu\{|f_m-f_n|>a\}=0.
\]
Es \emph{de Cauchy casi uniformemente} si para cada $\eps>0$ existe $A$ con $\mu(X\setminus A)<\eps$ en el que la sucesión es uniformemente de Cauchy.
\end{definicion}
\begin{teorema}\label{teo:v2-completitud-medida}
Toda sucesión de Cauchy en medida tiene un límite medible finito en medida, único salvo igualdad casi en todas partes. Además, tiene una subsucesión que converge casi uniformemente a ese límite.
\end{teorema}
\begin{proof}
Escogemos índices crecientes $n_k$ tales que
\[\mu(A_k)<2^{-k},\qquad A_k=\{|f_{n_{k+1}}-f_{n_k}|>2^{-k}\}.
\]
Para ello basta elegir $n_k$ más allá de un índice de Cauchy que sirva al umbral $2^{-k}$ y error de medida $2^{-k}$, haciendo crecientes esos índices. Pongamos $B_j=\bigcup_{k\ge j}A_k$. Entonces $\mu(B_j)\le2^{1-j}$ y $\mu(\bigcap_jB_j)=0$. Si $x\notin B_j$ y $q>p\ge j$,
\[|f_{n_q}(x)-f_{n_p}(x)|\le\sum_{k=p}^{q-1}2^{-k}\le2^{1-p}.
\]
La completitud de $\R$ o $\C$ produce un límite finito en $X\setminus\bigcap_jB_j$; lo definimos como cero en el conjunto excepcional. Es medible porque el conjunto de convergencia es medible y allí el límite es límite de funciones medibles. Al pasar $q\to\infty$, la misma cota prueba convergencia uniforme en $X\setminus B_j$. Dado $\eps>0$, elegimos $j$ con $2^{1-j}<\eps$; así obtenemos convergencia casi uniforme, y por tanto en medida.

Para toda la sucesión, fijemos $a,\eps>0$. Si $n,n_p$ son suficientemente grandes, la propiedad de Cauchy da $\mu\{|f_n-f_{n_p}|>a/2\}<\eps/2$. Escogiendo además $p$ suficientemente grande obtenemos $\mu\{|f_{n_p}-f|>a/2\}<\eps/2$. La desigualdad triangular y la inclusión de conjuntos malos dan $\mu\{|f_n-f|>a\}<\eps$.

Si $f$ y $g$ son dos límites en medida, para cada $a>0$,
\[\mu\{|f-g|>a\}\le\mu\{|f-f_n|>a/2\}+\mu\{|f_n-g|>a/2\}\longrightarrow0.
\]
La unión sobre $a=1,1/2,1/3,\ldots$ demuestra $f=g$ casi en todas partes.
\end{proof}
\begin{proposicion}
Una sucesión es de Cauchy casi uniformemente si y sólo si converge casi uniformemente a una función medible finita. Tal convergencia implica convergencia casi en todas partes y en medida.
\end{proposicion}
\begin{proof}
Para cada $j$ escojamos $A_j$ con $\mu(X\setminus A_j)<2^{-j}$ donde la sucesión es uniformemente de Cauchy. La completitud del cuerpo escalar da un límite uniforme en $A_j$. Estos límites coinciden en intersecciones, pues son límites de la misma sucesión numérica. La unión de los $A_j$ tiene complemento de medida cero: su medida es menor que $2^{-j}$ para cada $j$. Definiendo el límite puntual allí y cero fuera obtenemos una función medible, con convergencia uniforme en cada $A_j$. Esto prueba convergencia casi uniforme. La implicación inversa es el criterio uniforme de Cauchy. Para una sucesión casi uniformemente convergente se repite la unión de conjuntos buenos con excepciones $2^{-j}$ y se obtiene convergencia casi en todas partes. Finalmente, fijados $a,\eps>0$, en un conjunto bueno de complemento menor que $\eps$ el conjunto malo para el umbral $a$ acaba vacío; por tanto su medida total acaba siendo menor que $\eps$.
\end{proof}

\section{Ejercicios sobre límites y composición}
\begin{ejercicio}[Lectura 15, ejercicio 1]\label{ej:v2-l15-1}
Complete la extracción casi uniforme de una sucesión de Cauchy en medida.
\end{ejercicio}
\begin{solucion}
Con los conjuntos $A_k,B_j$ del teorema \ref{teo:v2-completitud-medida}, fuera de $B_j$ la desigualdad $|f_{n_p}-f|\le2^{1-p}$ vale simultáneamente para todo $p\ge j$. Dado $\eps$, primero fijamos $j$ para que $\mu(B_j)<\eps$; dado después $\delta$, escogemos $p_0\ge j$ con $2^{1-p_0}<\delta$. Para $p\ge p_0$ la cota es uniforme en el conjunto fijo $X\setminus B_j$. Este orden de elecciones completa el argumento solicitado.
\end{solucion}
\begin{ejercicio}[Lectura 15, ejercicio 2]\label{ej:v2-l15-2}
¿La convergencia casi uniforme implica convergencia uniforme fuera de un conjunto nulo?
\end{ejercicio}
\begin{solucion}
No. En $(0,1)$, $f_n=\ind{(0,1/n)}$ converge casi uniformemente a cero: para una excepción de medida $\eps$, retiramos $(0,\eps/2)$, y la sucesión es eventualmente cero en el resto. Pero para cualquier conjunto nulo $N$, el intervalo $(0,1/n)\setminus N$ es no vacío, de modo que $\sup_{(0,1)\setminus N}f_n=1$ para todo $n$.
\end{solucion}
\begin{ejercicio}[Lectura 15, ejercicio 3]\label{ej:v2-l15-3}
Analice $f_n=\ind{[n,n+1]}$ en $\R$ con medida de Lebesgue.
\end{ejercicio}
\begin{solucion}
Cada punto pertenece a lo sumo a dos de estos intervalos, así que $f_n\to0$ puntualmente. El conjunto $\{|f_n|>1/2\}$ tiene medida uno; no hay convergencia en medida y, por la proposición anterior, tampoco casi uniforme.
\end{solucion}
\begin{ejercicio}[Lectura 15, ejercicio 4]\label{ej:v2-l15-4}
Demuestre que toda subsucesión de una sucesión de Cauchy en medida también lo es.
\end{ejercicio}
\begin{solucion}
Dados umbral $a$ y error $\eps$, elijamos $N$ que sirve para todos los índices originales $m,n\ge N$. Si $n_k$ son los índices de la subsucesión, desde cierto $k_0$ se tiene $n_k\ge N$. La misma estimación sirve entonces para cualesquiera $k,l\ge k_0$.
\end{solucion}
\begin{ejercicio}[Lectura 15, ejercicio 5]\label{ej:v2-l15-5}
Caracterice las funciones borelianas $\varphi:\R\to\R$ que preservan convergencia casi en todas partes por composición.
\end{ejercicio}
\begin{solucion}
Toda función continua la preserva, aplicando continuidad a la sucesión numérica en cada punto donde $f_n\to f$. Para la recíproca debemos exigir $\mu(X)>0$, o formular la propiedad para todos los espacios de medida. Sean $t_n\to t$ y tomemos las funciones constantes $f_n=t_n$, $f=t$. La convergencia casi en todas partes de las composiciones obliga a $\varphi(t_n)\to\varphi(t)$: si la sucesión numérica no convergiera, fallaría en todo $X$, que tiene medida positiva. La continuidad secuencial equivale a continuidad en $\R$. Sin la hipótesis $\mu(X)>0$, la recíproca original es falsa: si $\mu$ es la medida cero, cualquier función boreliana preserva la propiedad vacía de converger casi en todas partes.
\end{solucion}
\begin{ejercicio}[Lectura 15, ejercicio 6]\label{ej:v2-l15-6}
Si $\varphi$ es uniformemente continua y $f_n\to f$ uniformemente, pruebe convergencia uniforme de las composiciones.
\end{ejercicio}
\begin{solucion}
Para $\eps>0$ elija $\delta>0$ tal que $|u-v|<\delta$ implique $|\varphi(u)-\varphi(v)|<\eps$. Desde cierto índice la primera desigualdad vale para $u=f_n(x),v=f(x)$ simultáneamente en todos los puntos, y por ello también la segunda.
\end{solucion}
\begin{ejercicio}[Lectura 15, ejercicio 7]\label{ej:v2-l15-7}
Pruebe el resultado anterior para convergencia casi uniforme y para convergencia en medida.
\end{ejercicio}
\begin{solucion}
En el caso casi uniforme aplicamos el ejercicio \ref{ej:v2-l15-6} en cada conjunto bueno, conservando su excepción. En medida, para un umbral $\eps$ elegimos $\delta$ por continuidad uniforme. Entonces
\[\{|\varphi(f_n)-\varphi(f)|\ge\eps\}\subset\{|f_n-f|\ge\delta\},
\]
y la medida del segundo conjunto tiende a cero.
\end{solucion}
\begin{ejercicio}[Lectura 15, ejercicio 8]\label{ej:v2-l15-8}
Si $\varphi$ no es uniformemente continua, construya funciones $f_n\to f$ uniformemente tales que $\varphi(f_n)$ no converja en medida a $\varphi(f)$.
\end{ejercicio}
\begin{solucion}
Existen $\eps_0>0$ y números $a_k,b_k$ con $|a_k-b_k|<1/k$ y $|\varphi(a_k)-\varphi(b_k)|\ge\eps_0$. En $\N$ con medida de conteo definimos $f(k)=a_k$ y $f_n(k)=a_k$ si $k<n$, $f_n(k)=b_k$ si $k\ge n$. Todas las funciones son medibles y $\sup_k|f_n(k)-f(k)|\le1/n$. Sin embargo, el conjunto donde las composiciones difieren al menos $\eps_0$ contiene $\{n,n+1,\ldots\}$ y tiene medida infinita. Esta construcción identifica la función de la continuidad \emph{uniforme}: el mismo control debe funcionar aun cuando los valores $f(k)$ recorran regiones distintas del dominio.
\end{solucion}

\chapter{El espacio de funciones integrables}\label{cap:v2-08}

La integral de $|f-g|$ mide el error total de una aproximación. Para convertirla en distancia hay que identificar funciones iguales casi en todas partes. Esa identificación permite estudiar simultáneamente completitud, aproximación y paso al límite. Reunimos las lecturas 13, 16 y 17 de Valencia \cite{valencia2}; el argumento por series para completitud también se explica en las notas de Maximenko \cite{maximenkoLp}.

\section{Seminormas, cocientes e igualdad casi en todas partes}\label{lectura:v2-13}
\begin{definicion}
Una \emph{seminorma}\index{seminorma} en un espacio vectorial $V$ sobre $\F=\R$ o $\C$ es una aplicación $p:V\to[0,\infty)$ tal que $p(\lambda v)=|\lambda|p(v)$ y $p(u+v)\le p(u)+p(v)$. Es norma si además $p(v)=0$ implica $v=0$.
\end{definicion}
Las condiciones implican $p(0)=0$ y $|p(u)-p(v)|\le p(u-v)$. El conjunto $W=\{v:p(v)=0\}$ es un subespacio: la homogeneidad y la desigualdad triangular preservan el valor cero. En el cociente $V/W$, dos vectores representan la misma clase si su diferencia pertenece a $W$. Las operaciones $[u]+[v]=[u+v]$, $\lambda[v]=[\lambda v]$ no dependen de los representantes, porque $W$ es subespacio. La fórmula
\[\|[v]\|=p(v)
\]
está bien definida por la desigualdad inversa; es una norma porque una clase tiene norma cero exactamente cuando es la clase de $W$.

\begin{ejemplo}
En las funciones Riemann integrables sobre $[a,b]$, $p(f)=\int_a^b|f|$ es una seminorma: cambiar el valor en un punto no cambia la integral. Si $X$ es compacto y $A\subset X$, la restricción $C(X)\to C(A)$ tiene núcleo $W_A=\{f:f|_A=0\}$, de modo que $C(X)/W_A$ identifica funciones con la misma restricción. Si $X$ no es compacto, la expresión $\sup_X|f|$ puede ser infinita para $f\in C(X)$; para usarla como norma hay que trabajar en $C_b(X)$ o imponer compacidad. Las afirmaciones sobre ideales de $C(X)$ tampoco se trasladan sin hipótesis del caso compacto al general.
\end{ejemplo}

Sea $\mathscr L^1(\mu)$ el espacio de funciones medibles $f$ con $\int|f|<\infty$. Las propiedades de la integral muestran que es un espacio vectorial y que $p(f)=\int|f|$ es seminorma. Su núcleo está formado exactamente por las funciones nulas casi en todas partes. Definimos
\[L^1(\mu)=\mathscr L^1(\mu)/\{f:f=0\text{ casi en todas partes}\},\qquad
\|f\|_1=\int_X|f|\,d\mu.
\]
Por comodidad escribimos $f$ para su clase. Una afirmación puntual sobre elementos de $L^1$ exige escoger representantes; una afirmación de norma debe ser independiente de esa elección.

\begin{ejercicio}[Lectura 13, ejercicio 1]\label{ej:v2-l13-1}
Si $T:V\to U$ es lineal, pruebe $V/\ker T\cong\operatorname{im}T$.
\end{ejercicio}

\begin{ejercicio}[Lectura 13, ejercicio 2]\label{ej:v2-l13-2}
Demuestre $\dim V=\dim W+\dim(V/W)$ para un subespacio $W$.
\end{ejercicio}

\begin{ejercicio}[Lectura 13, ejercicio 3]\label{ej:v2-l13-3}
Si $W\subset\ker T$, construya la aplicación inducida $V/W\to U$ y pruebe su unicidad.
\end{ejercicio}


\section{Completitud y aproximación por funciones simples}\label{lectura:v2-16}
\begin{teorema}[Riesz--Fischer para $L^1$]\label{teo:v2-L1-Banach}\index{Riesz--Fischer, teorema de}\index{L1@$L^1$!completitud}
El espacio $L^1(\mu)$ es de Banach para cualquier espacio de medida.
\end{teorema}
\begin{proof}
Primero consideremos una serie $\sum_ku_k$ tal que $\sum_k\|u_k\|_1<\infty$. Elegimos representantes medibles y ponemos $H=\sum_k|u_k|$. La convergencia monótona da
\[\int H\,d\mu=\sum_k\int|u_k|\,d\mu<\infty.
\]
En consecuencia, $H$ es finita fuera de un conjunto nulo. Allí la serie converge absolutamente a una función $u$, que definimos como cero en la excepción. Es medible, $|u|\le H$ casi en todas partes y
\[\left\|u-\sum_{k=1}^nu_k\right\|_1\le\sum_{k>n}\|u_k\|_1\longrightarrow0.
\]
Así, toda serie absolutamente convergente en norma converge en $L^1$.

Ahora sea $(f_n)$ de Cauchy. Escogemos una subsucesión con $\|f_{n_{k+1}}-f_{n_k}\|_1<2^{-k}$. La serie de diferencias converge en $L^1$ por el argumento anterior, y sus sumas parciales telescópicas muestran que $f_{n_k}$ converge en $L^1$ a cierto $f$. Dado $\eps>0$, tomamos un índice de Cauchy $N$ con $\|f_n-f_m\|_1<\eps/2$ para $m,n\ge N$, y un $n_k\ge N$ con $\|f_{n_k}-f\|_1<\eps/2$. Para $n\ge N$, la desigualdad triangular da $\|f_n-f\|_1<\eps$. La sucesión completa converge.
\end{proof}

\begin{proposicion}\label{prop:v2-simples-densas}\index{densidad!de funciones simples}
Las funciones simples integrables son densas en $L^1(\mu)$.
\end{proposicion}
\begin{proof}
Si $f\ge0$ es integrable, escojamos simples $0\le s_n\uparrow f$. Entonces $\|f-s_n\|_1=\int f-\int s_n\to0$ por convergencia monótona; cada $s_n$ es integrable por comparación. Para $f$ real aplicamos el argumento a sus partes positiva y negativa. Para $f$ compleja lo aplicamos a las partes real e imaginaria y usamos $|u+iv|\le|u|+|v|$. En una representación canónica de una simple integrable, todo conjunto con coeficiente no nulo tiene medida finita, pues $|a|\mu(E)\le\int|s|<\infty$.
\end{proof}

\section{Separabilidad del espacio de medida}
Para $A,B\in\Acal$, pongamos $\rho(A,B)=\mu(A\mathbin\triangle B)$, permitiendo infinito, y
\[d_\mu(A,B)=\frac{\rho(A,B)}{1+\rho(A,B)},\qquad \frac{\infty}{1+\infty}:=1.
\]
Esta es una seudométrica. En efecto, $A\triangle C\subset(A\triangle B)\cup(B\triangle C)$, así que $\rho$ satisface la desigualdad triangular extendida. La función creciente $h(t)=t/(1+t)$ es subaditiva: para $s,t<\infty$,
\[h(s+t)=\frac{s+t}{1+s+t}\le\frac{s}{1+s}+\frac{t}{1+t}.
\]
Si un término es infinito, la desigualdad es inmediata porque el miembro derecho es al menos uno. Identificando conjuntos que difieren en un nulo se obtiene un espacio métrico, llamado aquí \emph{álgebra de medida}. Decimos que el espacio de medida es separable si esta métrica admite un conjunto numerable denso. La condición $d_\mu(A,B)<\eps/(1+\eps)$ equivale a $\mu(A\triangle B)<\eps$; no deben confundirse las dos cantidades.

\begin{teorema}\label{teo:v2-medida-separable}\index{separabilidad!del algebra de medida@del álgebra de medida}
Si el álgebra de medida es separable, entonces $L^1(\mu)$ es separable. Si además $\mu$ es semifinita, entonces es $\sigma$-finita.
\end{teorema}
\begin{proof}
Todo subespacio de un espacio métrico separable es separable: las bolas con centros en un conjunto numerable denso y radios racionales forman una base numerable; de cada elemento de la base que interseca el subespacio elegimos un punto del subespacio. Apliquemos esto al conjunto de clases representadas por conjuntos de medida finita. Obtenemos una familia numerable $(E_j)$, con $\mu(E_j)<\infty$, que aproxima en medida de diferencia simétrica a cualquier conjunto de medida finita.

Las combinaciones lineales finitas de los $\ind{E_j}$ con coeficientes racionales, o racionales complejos, forman un conjunto numerable. Aproximan a cada simple integrable: si $s=\sum_{i=1}^ra_i\ind{A_i}$ con $\mu(A_i)<\infty$, elegimos $E_{j_i}$ con $\sum_i|a_i|\mu(A_i\triangle E_{j_i})$ tan pequeño como se desee. Luego aproximamos cada $a_i$ por un racional $q_i$, haciendo pequeña $\sum_i|a_i-q_i|\mu(E_{j_i})$. Ambas estimaciones se suman por la desigualdad triangular. La proposición \ref{prop:v2-simples-densas} concluye separabilidad de $L^1$.

Para la última afirmación, recordemos que semifinitud significa que todo conjunto de medida infinita contiene un subconjunto medible de medida \emph{positiva y finita}; equivalentemente, la medida de cualquier conjunto es el supremo de las medidas finitas de sus subconjuntos. Sea $E=\bigcup_jE_j$. Si $X\setminus E$ contuviera un conjunto $B$ de medida positiva finita, entonces $\mu(B\triangle E_j)\ge\mu(B)>0$ para todo $j$, contradiciendo la densidad de $(E_j)$ entre los conjuntos de medida finita. Por semifinitud, $\mu(X\setminus E)=0$. Así, los $E_j$ y el complemento nulo constituyen una cubierta numerable por conjuntos de medida finita.
\end{proof}

\begin{observacion}
La separabilidad de $L^1$ no implica separabilidad del álgebra de medida, como muestra el ejercicio \ref{ej:v2-l16-4}. Tampoco la medida finita implica separabilidad de $L^1$. Para verlo, sea $I$ no numerable y $X=\{0,1\}^I$. Consideremos los conjuntos $\pi_J^{-1}(B)$, donde $J\subset I$ es numerable y $B$ es boreliano en $\{0,1\}^J$. Forman una $\sigma$-álgebra: una sucesión de tales conjuntos depende sólo de la unión numerable de sus coordenadas. Sobre cada espacio numerable de coordenadas existe la probabilidad de monedas independientes, construida con los dígitos binarios de un punto uniforme de $[0,1)$; un cilindro que prescribe $r$ dígitos tiene medida $2^{-r}$, y la unicidad de extensión determina la medida de todos los borelianos. Definimos $\mu(\pi_J^{-1}(B))$ mediante esa probabilidad. La definición es independiente de $J$, porque al pasar a la unión de dos conjuntos de coordenadas las marginales coinciden en cilindros y, por unicidad, en todos los borelianos. La aditividad numerable se comprueba llevando toda la sucesión a un único conjunto numerable de coordenadas. Tenemos así una probabilidad en $X$. Para $i\ne j$, las funciones coordenadas $x\mapsto x_i$ y $x\mapsto x_j$ distan $1/2$ en $L^1$. Una familia no numerable con separación uniforme positiva impide separabilidad: bolas disjuntas de radio menor que $1/4$ requerirían puntos distintos de cualquier conjunto denso.
\end{observacion}

\section{Ejercicios sobre cocientes, densidad y separabilidad}
\begin{ejercicio}[Lectura 16, ejercicio 1]\label{ej:v2-l16-1}
Identifique la completación de las funciones Riemann integrables en $[a,b]$ con la seminorma integral.
\end{ejercicio}

\begin{ejercicio}[Lectura 16, ejercicio 2]\label{ej:v2-l16-2}
Si $d$ es una seudométrica, pruebe que identificar puntos a distancia cero produce una métrica.
\end{ejercicio}

\begin{ejercicio}[Lectura 16, ejercicio 3]\label{ej:v2-l16-3}
Caracterice la separabilidad del cociente mediante una sucesión de representantes.
\end{ejercicio}

\begin{ejercicio}[Lectura 16, ejercicio 4]\label{ej:v2-l16-4}
Compare la separabilidad del álgebra de medida de $\N$ con medida de conteo y la de su espacio $L^1$.
\end{ejercicio}

\begin{ejercicio}[Lectura 16, ejercicio 5]\label{ej:v2-l16-5}
Pruebe que $L^1(\R,m)$ es separable.
\end{ejercicio}

\begin{ejercicio}[Lectura 16, ejercicio 6]\label{ej:v2-l16-6}
Examine la afirmación de que la medida de conteo en la $\sigma$-álgebra numerable--conumerable de $\R$ produce $L^1$ separable sin ser $\sigma$-finita. Dé una versión correcta.
\end{ejercicio}


\section{Convergencia en media y control de la masa}\label{lectura:v2-17}
\begin{definicion}
Decimos que $f_n\to f$ \emph{en media}\index{convergencia!en media} si $\|f_n-f\|_1\to0$.
\end{definicion}
Por Markov,
\[\mu\{|f_n-f|>a\}\le\frac{\|f_n-f\|_1}{a},\qquad a>0.
\]
Así, convergencia en media implica convergencia en medida y permite extraer una subsucesión que converge casi uniformemente y casi en todas partes. La desigualdad tiene dirección $\le$, que es esencial para esta deducción. Si $\mu(X)<\infty$, la convergencia uniforme de funciones integrables implica convergencia en media: primero $|f|\le|f_N|+1$ para algún $N$ muestra que $f$ es integrable, y luego $\|f_n-f\|_1\le\mu(X)\sup_X|f_n-f|$.

\begin{teorema}[Convergencia dominada en medida]\label{teo:v2-dominada-medida}
Si $f_n\to f$ en medida y $|f_n|\le g\in L^1$, entonces $f\in L^1$ y $\|f_n-f\|_1\to0$.
\end{teorema}
\begin{proof}
Una subsucesión converge a $f$ casi en todas partes, por lo que $|f|\le g$ casi en todas partes y $f$ es integrable. Si la conclusión fallara, existirían $\eps_0>0$ y una subsucesión con $\|f_{n_k}-f\|_1\ge\eps_0$ para todo $k$. Esa subsucesión aún converge en medida y admite otra que converge casi en todas partes. En esta última, $|f_{n_{k_j}}-f|\le2g$; convergencia dominada da normas tendiendo a cero, contradicción.
\end{proof}

\begin{lema}\label{lema:v2-control-integral}
Para $g\in L^1$, la integral de $|g|$ es absolutamente continua respecto de la medida: dado $\eps>0$, existe $\delta>0$ tal que $\mu(E)<\delta$ implica $\int_E|g|<\eps$. Además, existe un conjunto $A$ de medida finita con $\int_{X\setminus A}|g|<\eps$.
\end{lema}
\begin{proof}
Las funciones $|g|\ind{\{|g|>M\}}$ tienden a cero casi en todas partes cuando $M$ crece. Convergencia dominada permite elegir $M>0$ para que su integral sea menor que $\eps/2$. Tomemos $\delta=\eps/(2M)$. Si $\mu(E)<\delta$, entonces
\[\int_E|g|\le M\mu(E)+\int_{\{|g|>M\}}|g|<\eps.\]
Para la segunda afirmación tomamos $A_k=\{|g|\ge1/k\}$. Markov da $\mu(A_k)\le k\|g\|_1<\infty$, mientras $|g|\ind{X\setminus A_k}\to0$ puntualmente y está dominada por $|g|$. La convergencia dominada permite elegir $k$.
\end{proof}

\begin{teorema}[Criterio de Vitali en $L^1$]\label{teo:v2-vitali-L1}\index{Vitali!criterio de convergencia}
Sean $f_n,f\in L^1(\mu)$. Entonces $\|f_n-f\|_1\to0$ si y sólo si se cumplen:
\begin{enumerate}
\item $f_n\to f$ en medida;
\item para cada $\eps>0$ existe $A\in\Acal$ con $\mu(A)<\infty$ y $\sup_n\int_{X\setminus A}|f_n|<\eps$;
\item para cada $\eps>0$ existe $\delta>0$ tal que $\mu(E)<\delta$ implica $\sup_n\int_E|f_n|<\eps$.
\end{enumerate}
\end{teorema}
\begin{proof}
Si hay convergencia en $L^1$, ya probamos (i). Para (ii), elija $N$ con $\|f_n-f\|_1<\eps/2$ para $n\ge N$. Por el lema anterior hay un conjunto de medida finita que recoge toda la integral de $|f|$ salvo $\eps/2$. Añadiendo un conjunto de medida finita apropiado para cada una de las funciones $f_1,\ldots,f_{N-1}$, obtenemos una unión finita $A$ que sirve a todas. Para $n\ge N$, la integral exterior es menor que $\|f_n-f\|_1+\int_{X\setminus A}|f|<\eps$; para los primeros índices se impuso directamente. Para (iii), escogemos $N$ con la misma estimación y aplicamos la primera parte del lema a la familia finita $f,f_1,\ldots,f_{N-1}$; el mínimo de sus tolerancias sirve a toda la sucesión, usando de nuevo la desigualdad triangular.

Recíprocamente, fijemos $\eps>0$. Por (ii), y ampliando $A$ mediante el lema aplicado a $f$, podemos lograr $\mu(A)<\infty$ y
\[\sup_n\int_{X\setminus A}|f_n|<\eps,\qquad \int_{X\setminus A}|f|<\eps.
\]
Por (iii) y la continuidad absoluta de la integral de $|f|$, escogemos $\delta>0$ tal que cualquier $E$ de medida menor que $\delta$ satisfaga $\sup_n\int_E|f_n|<\eps$ y $\int_E|f|<\eps$. Sea $a=\eps/(1+\mu(A))>0$. Por (i), desde cierto índice $E_n=\{|f_n-f|>a\}$ tiene medida menor que $\delta$. Descomponiendo $X$ en $X\setminus A$, $A\cap E_n$ y $A\setminus E_n$, obtenemos
\[\|f_n-f\|_1<2\eps+2\eps+a\mu(A)<5\eps.
\]
Como $\eps$ es arbitrario, la norma tiende a cero.
\end{proof}
\begin{observacion}[Una hipótesis que no puede omitirse]
La condición $\mu(A)<\infty$ en (ii) es necesaria. Sin ella se podría tomar $A=X$, haciendo vacía la condición. En $\R$, $f_n=n^{-1}\ind{[0,n]}$ converge uniformemente a cero y cumple (iii), pues $\int_E|f_n|\le m(E)$. Sin embargo, $\|f_n\|_1=1$. La masa se dispersa sobre conjuntos cada vez mayores. En contraste, $n\ind{(0,1/n)}$ concentra su masa en conjuntos pequeños y viola (iii). Los dos controles de Vitali impiden fenómenos diferentes.
\end{observacion}

\section{Ejercicios sobre integrabilidad, convexidad y traslaciones}
\begin{ejercicio}[Lectura 17, ejercicio 1]\label{ej:v2-l17-1}
Demuestre la unicidad del límite en media.
\end{ejercicio}

\begin{ejercicio}[Lectura 17, ejercicio 2]\label{ej:v2-l17-2}
Analice $f_n=n\ind{[1/n,2/n]}$ en $[0,1]$.
\end{ejercicio}

\begin{ejercicio}[Lectura 17, ejercicio 3]\label{ej:v2-l17-3}
Si $f_n\to f$ en medida y $|f_n|\le g\in L^1$, pruebe $\int f_n\to\int f$.
\end{ejercicio}

\begin{ejercicio}[Lectura 17, ejercicio 4]\label{ej:v2-l17-4}
Compare las condiciones de control de masa de Vitali cuando $|f_n|\le g\in L^1$.
\end{ejercicio}

\begin{ejercicio}[Lectura 17, ejercicio 5: Jensen]\label{ej:v2-l17-5}\index{Jensen, desigualdad de}
Sea $\mu(X)=1$, $f\in L^1$ real y $\varphi:\R\to\R$ convexa. Pruebe continuidad de $\varphi$, existencia de una recta de soporte en cada punto y $\varphi(\int f)\le\int\varphi(f)$.
\end{ejercicio}

\begin{ejercicio}[Lectura 17, ejercicio 6]\label{ej:v2-l17-6}
Si $\mu(X)<\infty$, $\varphi$ es continua y $|\varphi(t)|\le K|t|$ para $|t|\ge K$, demuestre que $f\in L^1$ real implica $\varphi\circ f\in L^1$. Establezca la versión para medida infinita.
\end{ejercicio}

\begin{ejercicio}[Lectura 17, ejercicio 7]\label{ej:v2-l17-7}
Bajo las hipótesis del ejercicio \ref{ej:v2-l17-6}, pruebe que $f_n\to f$ en $L^1$ implica $\varphi(f_n)\to\varphi(f)$ en $L^1$.
\end{ejercicio}

\begin{ejercicio}[Lectura 17, ejercicio 8]\label{ej:v2-l17-8}
Si $f\in L^1(\R)$ y $f_x(t)=f(t-x)$, pruebe que $x\mapsto f_x$ es continua con valores en $L^1$.
\end{ejercicio}

\parte{Espacios de funciones y dualidad}{Las desigualdades de Hölder y Minkowski organizan los espacios $L^p$. Su geometría permite representar funcionales y distinguir los exponentes interiores de los extremos $1$ e infinito. Se utilizan espacios normados, Hahn--Banach y geometría de Hilbert del volumen I, además de la integral y la convergencia en medida.}
\chapter{Los espacios \texorpdfstring{$L^p$}{Lp} y la convergencia en norma}\label{v2:cap09}\label{cap:v2-09}
\section{Medir el tamaño medio de una función}\label{lectura:v2-18}
La integral de $|f|$ mide un tamaño total. Al integrar $|f|^p$, con $p>1$, los valores grandes reciben un peso mayor. Esta modificación produce una familia de geometrías, no una sola noción de convergencia. Las lecturas 18 y 19 de \cite{valencia2} se desarrollan aquí distinguiendo las funciones medibles de sus clases de equivalencia. Trabajamos sobre un espacio de medida arbitrario $(X,\Acal,\mu)$ y sobre $\F=\R$ o $\C$; no se presupone que la medida sea finita ni completa.

\begin{definicion}
Para $1\le p<\infty$, sea $\mathcal L^p(\mu)$ el espacio de funciones medibles $f:X\to\F$ tales que $\int_X|f|^p\,d\mu<\infty$. Escribimos
\[
 \norm f_p=\left(\int_X|f|^p\,d\mu\right)^{1/p}.
\]
Por ahora esta expresión es una seminorma: dos funciones que difieren solamente en un conjunto nulo tienen distancia cero.\index{espacio Lp@espacio $L^p$}
\end{definicion}
La convexidad de $t\mapsto t^p$ da $(a+b)^p\le2^{p-1}(a^p+b^p)$ para $a,b\ge0$: se aplica la definición de convexidad a su promedio y se multiplica por $2^p$. Por tanto $|f+g|^p\le2^{p-1}(|f|^p+|g|^p)$. Esta cota prueba que $\mathcal L^p$ es un espacio vectorial antes de utilizar una desigualdad triangular para su seminorma. Además, $\norm f_p=0$ equivale a $f=0$ casi en todas partes: si la integral se anula, cada conjunto $\{|f|\ge1/n\}$ tiene medida cero por monotonía, y su unión es $\{f\ne0\}$.

\section{Young, Hölder y Minkowski}
\begin{lema}[Young]\label{v2:young}\index{Young, desigualdad de}
Si $1<p<\infty$, $q=p/(p-1)$ y $a,b\ge0$, entonces $ab\le a^p/p+b^q/q$.
\end{lema}
\begin{proof}
Para $b$ fijo, la función $H(a)=a^p/p-ab+b^q/q$ tiene derivada $a^{p-1}-b$. Disminuye hasta $a=b^{1/(p-1)}$ y aumenta después. Su valor mínimo es cero, porque en ese punto $a^p=b^q=ab$ y $1/p+1/q=1$. Si $b=0$, la desigualdad es inmediata. Esta prueba incluye los extremos donde no se puede tomar un logaritmo.
\end{proof}
\begin{teorema}[Hölder y Cauchy--Schwarz]\label{v2:holder}\index{Holder, desigualdad de@Hölder, desigualdad de}
Para exponentes conjugados $p,q>1$, si $f\in\mathcal L^p$ y $g\in\mathcal L^q$, entonces $fg\in\mathcal L^1$ y
\[
 \int_X|fg|\,d\mu\le\norm f_p\norm g_q.
\]
En particular, $\left|\int f\overline g\,d\mu\right|\le\norm f_2\norm g_2$.
\end{teorema}
\begin{proof}
Si una norma es cero, su función se anula fuera de un conjunto medible nulo; allí también se anula el producto. En otro caso aplicamos Young a $|f|/\norm f_p$ y $|g|/\norm g_q$. La integral del lado derecho es $1/p+1/q=1$, así que el producto normalizado es integrable. Multiplicar por las normas da la desigualdad. La última afirmación usa $p=q=2$ y $|\int h|\le\int|h|$.
\end{proof}
\begin{teorema}[Minkowski]\label{v2:minkowski}\index{Minkowski, desigualdad de}
Para $1\le p<\infty$, $\norm{f+g}_p\le\norm f_p+\norm g_p$. En consecuencia, $\norm{\cdot}_p$ es una seminorma sobre $\mathcal L^p$.
\end{teorema}
\begin{proof}
Para $p=1$ se integra la desigualdad triangular puntual. Para $p>1$, la pertenencia $f+g\in\mathcal L^p$ ya está probada. Si $q=p/(p-1)$, entonces $|f+g|^{p-1}\in\mathcal L^q$, porque $(p-1)q=p$, y su norma es $\norm{f+g}_p^{p-1}$. Por Hölder,
\[
 \norm{f+g}_p^p\le\int |f+g|^{p-1}(|f|+|g|)\,d\mu
 \le\norm{f+g}_p^{p-1}(\norm f_p+\norm g_p).
\]
Cuando la norma del lado izquierdo es positiva dividimos por su potencia $p-1$; cuando es cero no hace falta dividir. La homogeneidad se sigue de $|\alpha f|^p=|\alpha|^p|f|^p$.
\end{proof}
\begin{proposicion}[Medida finita]\label{v2:inclusion-finita}
Si $1\le p<q<\infty$ y $\mu(X)<\infty$, entonces $\mathcal L^q\subset\mathcal L^p$ y
\[
 \norm f_p\le\mu(X)^{1/p-1/q}\norm f_q.
\]
Si $\mu(X)=0$, ambas seminormas son cero. Si $\mu(X)=1$, la constante es uno.
\end{proposicion}
\begin{proof}
Suponiendo $0<\mu(X)<\infty$, aplicamos Hölder a $|f|^p$ y $1$, con exponentes $q/p$ y $q/(q-p)$. Resulta $\int|f|^p\le(\int|f|^q)^{p/q}\mu(X)^{1-p/q}$; tomar la raíz $p$-ésima concluye. Otra prueba de la inclusión divide $X$ en $\{|f|<1\}$ y $\{|f|\ge1\}$: en el primer conjunto $|f|^p\le1$, y en el segundo $|f|^p\le|f|^q$. Esta segunda prueba no produce la mejor constante.
\end{proof}

\section{Ejercicios resueltos: exponentes y desigualdades}
\begin{ejercicio}[Lectura 18, ejercicio 1]\label{ej:v2-l18-1}
Para $a,b>0$ y exponentes conjugados $p,q>1$, demostrar que hay igualdad en Young si y solo si $a^p=b^q$.
\end{ejercicio}
\begin{solucion}
La derivada de $H$ en la prueba del lema \ref{v2:young} cambia de signo en un único punto. Por tanto $H(a)=0$ exactamente cuando $a^{p-1}=b$. Elevar a $q$ da $a^p=b^q$, y cada una de estas operaciones es reversible para números positivos.
\end{solucion}
\begin{ejercicio}[Lectura 18, ejercicio 2]\label{ej:v2-l18-2}
Demostrar $ab\le\min_{t>0}\bigl(t^pa^p/p+t^{-q}b^q/q\bigr)$ para $a,b>0$ y $1/p+1/q=1$.
\end{ejercicio}
\begin{solucion}
Young aplicado a $ta$ y $t^{-1}b$ da la desigualdad para cada $t$. Hay igualdad cuando $t^pa^p=t^{-q}b^q$, es decir, cuando $t=(b^q/a^p)^{1/(p+q)}$. Ese número existe y es positivo; por tanto el mínimo existe y en realidad es exactamente $ab$.
\end{solucion}
\begin{ejercicio}[Lectura 18, ejercicio 3]\label{ej:v2-l18-3}
Si $f\in\mathcal L^p\cap\mathcal L^q$, demostrar $f\in\mathcal L^r$ para todo $r$ entre $p$ y $q$. Deducir que $\{p\ge1:f\in\mathcal L^p\}$ es un intervalo, posiblemente vacío.
\end{ejercicio}
\begin{solucion}
Podemos suponer $p<r<q$. En $\{|f|\le1\}$ se tiene $|f|^r\le|f|^p$, y en su complemento $|f|^r\le|f|^q$. La suma de las dos integrales es finita. Por tanto, si dos exponentes pertenecen al conjunto indicado, todos los intermedios también pertenecen; esta es precisamente la propiedad que caracteriza a un intervalo de la recta, incluidos singletons e intervalos abiertos o semiabiertos.
\end{solucion}
\begin{ejercicio}[Lectura 18, ejercicio 4]\label{ej:v2-l18-4}
Si $\mu(X)=1$ y $1\le p<q<\infty$, demostrar $\norm f_p\le\norm f_q$ para $f\in\mathcal L^q$.
\end{ejercicio}
\begin{solucion}
La proposición \ref{v2:inclusion-finita} da la desigualdad con factor $1^{1/p-1/q}=1$. En concreto, Hölder aplicado a $|f|^p\cdot1$ da $\int|f|^p\le(\int|f|^q)^{p/q}$, y al tomar raíces se obtiene exactamente la afirmación.
\end{solucion}
\begin{ejercicio}[Lectura 18, ejercicio 5]\label{ej:v2-l18-5}
Sean $1\le p<q<\infty$. (i) Demostrar $\ell^p\subset\ell^q$. (ii) Encontrar una sucesión de $\ell^q\setminus\ell^p$. (iii) Encontrar una sucesión acotada que tienda a cero y no pertenezca a ningún $\ell^p$, $p\ge1$.
\end{ejercicio}
\begin{solucion}
(i) Si $\sum|x_n|^p<\infty$, entonces $x_n\to0$: en otro caso habría infinitos términos cuya potencia $p$ supera una constante positiva. Para $n$ grande $|x_n|\le1$, luego $|x_n|^q\le|x_n|^p$; los términos iniciales son finitos. También $\norm x_q\le\norm x_p$, porque $|x_n|\le\norm x_p$ implica $\sum|x_n|^q\le\norm x_p^{q-p}\sum|x_n|^p$.

(ii) Tomamos $x_n=n^{-1/p}$. Su serie de potencias $p$ es armónica y diverge, mientras que $\sum n^{-q/p}$ converge porque $q/p>1$.

(iii) Sirve $x_n=1/\log(n+1)$. Es acotada y tiende a cero. Para cada $p$ fijo, $\log(n+1)^p\le n^{1/2}$ cuando $n$ es suficientemente grande; esto se obtiene escribiendo $n=e^t$ y usando que $t^p e^{-t/2}\to0$. Así $x_n^p\ge n^{-1/2}$ eventualmente y la serie diverge.
\end{solucion}
\begin{ejercicio}[Lectura 18, ejercicio 6]\label{ej:v2-l18-6}
Demostrar que $\langle f,g\rangle=\int_X f\overline g\,d\mu$ es sesquilineal sobre $\mathcal L^2$ y que $\langle f,f\rangle=0$ si y solo si $f=0$ casi en todas partes.
\end{ejercicio}
\begin{solucion}
Hölder garantiza que todas las integrales existen. La linealidad de la integral da linealidad en el primer argumento y conjugado-linealidad en el segundo. Además $\overline{\int f\overline g}=\int g\overline f$ y $\langle f,f\rangle=\int|f|^2\ge0$. La caracterización de integral nula demostrada al inicio da el último enunciado. Sobre funciones aún puede haber vectores no idénticamente cero de longitud cero; el producto interno genuino aparece al pasar al cociente. En el caso complejo la barra es indispensable.
\end{solucion}
\begin{ejercicio}[Lectura 18, ejercicio 7: rectificación]\label{ej:v2-l18-7}
El enunciado original pide demostrar que $r\mapsto\norm f_r$ es convexa en $[p,q]$ si $f\in\mathcal L^p\cap\mathcal L^q$. Examinar esta afirmación y establecer la desigualdad de interpolación válida.
\end{ejercicio}
\begin{solucion}
La afirmación no es cierta en un espacio de medida general, ni siquiera en uno de probabilidad. Sea $f=\ind{(0,1/2)}$ en $(0,1)$ con medida de Lebesgue. Entonces $\norm f_r=2^{-1/r}$ y, poniendo $a=\log2$,
\[
 \frac{d^2}{dr^2}e^{-a/r}=e^{-a/r}\frac{a(a-2r)}{r^4}<0\qquad(r\ge1).
\]
Es estrictamente cóncava, por lo que cualquier intervalo no degenerado contradice la convexidad solicitada.

La afirmación correcta es la convexidad de $s\mapsto\log\norm f_{1/s}$ cuando $f$ no es cero casi en todas partes. Si $1/r=\theta/p+(1-\theta)/q$, $0<\theta<1$, aplicamos Hölder a $|f|^{\theta r}$ y $|f|^{(1-\theta)r}$ con exponentes $p/(\theta r)$ y $q/((1-\theta)r)$. Obtenemos
\[
 \norm f_r\le\norm f_p^{\theta}\norm f_q^{1-\theta}.
\]
Tomar logaritmos da la convexidad indicada; la desigualdad sigue siendo válida cuando $f=0$. Además, Young ponderada da $\norm f_r\le\theta\norm f_p+(1-\theta)\norm f_q$, con pesos referidos a los exponentes recíprocos, no a $p,q$ directamente.
\end{solucion}
\begin{ejercicio}[Lectura 18, ejercicio 8]\label{ej:v2-l18-8}
Supóngase $\mu(X)<\infty$, sea $f:X\to\F$ medible y sea $E_n=\{n-1\le|f|<n\}$. Para $1\le p<\infty$, demostrar que $f\in\mathcal L^p$ si y solo si $\sum_{n\ge1}n^p\mu(E_n)<\infty$.
\end{ejercicio}
\begin{solucion}
Los $E_n$ son disjuntos y cubren $X$, pues $f$ toma valores finitos. Por aditividad de la integral,
$\int|f|^p\le\sum_n n^p\mu(E_n)$. Recíprocamente, para $n\ge2$ tenemos $n\le2(n-1)\le2|f|$ en $E_n$. Por tanto
\[
 \sum_{n\ge1}n^p\mu(E_n)\le\mu(E_1)+2^p\int_X|f|^p\,d\mu<\infty.
\]
La finitud de $\mu(X)$ controla el primer estrato, donde $|f|$ puede ser cero.
\end{solucion}
\begin{ejercicio}[Lectura 18, ejercicio 9]\label{ej:v2-l18-9}
Para $f(x)=x^{-1/2}(1+|\log x|)^{-1}$ en $(0,\infty)$ con medida de Lebesgue, demostrar $f\in\mathcal L^p$ si y solo si $p=2$.
\end{ejercicio}
\begin{solucion}
En $(0,1)$ hacemos $x=e^{-t}$ y en $(1,\infty)$ hacemos $x=e^t$. Las integrales se convierten respectivamente en
\[
 \int_0^\infty\frac{e^{(p/2-1)t}}{(1+t)^p}\,dt,
 \qquad \int_0^\infty\frac{e^{(1-p/2)t}}{(1+t)^p}\,dt.
\]
Para $p=2$ ambas son $\int_0^\infty(1+t)^{-2}dt=1$. Si $p>2$, el primer integrando crece exponencialmente salvo un factor polinómico y su integral diverge; si $p<2$, ocurre lo mismo con el segundo. Para justificar esa comparación, si $c>0$, entonces $(1+t)^p\le e^{ct/2}$ eventualmente, así que $e^{ct}/(1+t)^p\ge e^{ct/2}$. Se han controlado por separado el origen y el infinito.
\end{solucion}

\section{Clases de equivalencia, completitud y densidad}\label{lectura:v2-19}
\begin{definicion}
Definimos $L^p(\mu)=\mathcal L^p(\mu)/\{f:f=0\text{ casi en todas partes}\}$. La norma de $[f]$ es $\norm f_p$. Si $f=g$ casi en todas partes, sus potencias absolutas tienen la misma integral; la definición no depende del representante. La distancia $\norm{f-g}_p$ se llama distancia en $p$-promedio. Más generalmente puede hablarse de $f_n\to f$ en $p$-promedio cuando $f_n-f\in\mathcal L^p$ y esas normas tienden a cero, aunque $f$ no sea integrable.
\end{definicion}
\begin{teorema}[Completitud de $L^p$]\label{v2:lp-completo}\index{espacio Lp@espacio $L^p$!completitud}
Para $1\le p<\infty$, $L^p(\mu)$ es un espacio de Banach.
\end{teorema}
\begin{proof}
Primero estudiamos una serie $\sum h_j$ con $\sum_j\norm{h_j}_p=M<\infty$, eligiendo representantes medibles. Para $G_n=\sum_{j\le n}|h_j|$, Minkowski da $\norm{G_n}_p\le M$. Por convergencia monótona,
\[
 G=\sum_j|h_j|,\qquad \int G^p\,d\mu\le M^p.
\]
Así $G<\infty$ fuera de un conjunto medible nulo $N$. La serie converge absolutamente allí. Definimos $h=\sum_jh_j$ fuera de $N$ y $h=0$ en $N$; es medible y $|h|\le G$ casi en todas partes. La cola satisface, por Fatou aplicado a sus sumas finitas,
\[
 \left\|h-\sum_{j=1}^n h_j\right\|_p\le\sum_{j>n}\norm{h_j}_p\longrightarrow0.
\]
Esto prueba convergencia en norma sin confundir $G$ con $G^p$ ni la norma de una potencia con la potencia de la norma.

Ahora, si $(f_n)$ es de Cauchy, elegimos índices crecientes $n_j$ con $\norm{f_{n_{j+1}}-f_{n_j}}_p<2^{-j}$. La serie $f_{n_1}+\sum_j(f_{n_{j+1}}-f_{n_j})$ converge por lo anterior; sus sumas parciales son la subsucesión $f_{n_j}$. Si su límite es $f$, para $n,n_j$ grandes la desigualdad $\norm{f_n-f}_p\le\norm{f_n-f_{n_j}}_p+\norm{f_{n_j}-f}_p$ prueba que toda la sucesión converge a $f$.
\end{proof}
\begin{teorema}[Funciones simples y separabilidad]\label{v2:lp-densidad}\index{separabilidad!de Lp@de $L^p$}
Las funciones simples integrables son densas en $L^p$, $1\le p<\infty$. Si la familia de conjuntos medibles de medida finita, identificados módulo conjuntos nulos y con distancia $\mu(A\mathbin\triangle B)$, es separable, entonces $L^p$ es separable.
\end{teorema}
\begin{proof}
Sea $f\in\mathcal L^p$. Los conjuntos $E_m=\{1/m\le|f|\le m\}$ tienen medida finita, porque $\mu(E_m)\le m^p\norm f_p^p$. Las funciones $f\ind{E_m}$ convergen a $f$ puntualmente y su error elevado a $p$ está dominado por $|f|^p$; por convergencia dominada convergen en norma. En $E_m$, aproximamos el recorrido acotado por una cuadrícula finita en $\R$ o $\C$: una función simple $s$ con $|f-s|\le\eta$ allí y cero fuera satisface $\norm{f\ind{E_m}-s}_p\le\eta\mu(E_m)^{1/p}$. Como su soporte tiene medida finita, $s$ es integrable.

Para separabilidad, elijamos una familia numerable densa $\{A_j\}$ de conjuntos de medida finita. La identidad
\[
 \norm{\ind A-\ind B}_p^p=\mu(A\mathbin\triangle B)
\]
permite aproximar cada indicador integrable por indicadores de esta familia. Las combinaciones finitas con coeficientes racionales, o con partes real e imaginaria racionales, forman un conjunto numerable. Aproximan las funciones simples y, por la primera parte, todo $L^p$. La aplicación de indicadores preserva las clases nulas; la distancia se transforma mediante la raíz $p$-ésima, lo cual preserva las nociones de densidad y convergencia.
\end{proof}

\section{Convergencia en medida, dominación y Vitali}
\begin{teorema}\label{v2:lp-convergencias}
Sea $1\le p<\infty$.
\begin{enumerate}
\item Si $\mu(X)<\infty$, $f_n\in L^p$ y $f_n\to f$ uniformemente, entonces $f\in L^p$ y $f_n\to f$ en norma.
\item Si $f_n\to f$ casi en todas partes o en medida y $|f_n|\le g$ casi en todas partes para un $g\in L^p$, entonces $f\in L^p$ y $\norm{f_n-f}_p\to0$.
\item La convergencia en $L^p$ implica convergencia en medida y admite una subsucesión convergente casi en todas partes.
\end{enumerate}
\end{teorema}
\begin{proof}
(i) Para $n$ grande, $\norm{f_n-f}_p\le\mu(X)^{1/p}\sup_X|f_n-f|$. La finitud de esta norma y $f_n\in L^p$ muestran primero que $f\in L^p$; la estimación prueba después la convergencia.

(ii) En el caso puntual, fuera de la unión numerable de los conjuntos excepcionales tenemos $|f|\le g$, y $|f_n-f|^p\le2^p g^p\in L^1$. La convergencia dominada da la conclusión. En el caso de convergencia en medida, toda subsucesión tiene una subsucesión convergente casi en todas partes: se eligen índices con $\mu(\{|f_{n_j}-f|>2^{-j}\})<2^{-j}$, y la medida de la unión de esos conjuntos para $j\ge m$ tiende a cero. Aplicando el caso puntual se obtiene $|f|\le g$ casi en todas partes y convergencia en norma a lo largo de esa subsucesión. Si las normas de error originales no tendieran a cero, alguna subsucesión las mantendría mayores que un $\eps>0$, contradicción.

(iii) La desigualdad de Chebyshev da $\mu(\{|f_n-f|>a\})\le a^{-p}\norm{f_n-f}_p^p\to0$ para cada $a>0$. La extracción que acabamos de describir proporciona la subsucesión puntual.
\end{proof}
\begin{observacion}[La hipótesis de dominación no puede debilitarse]
La lectura 19 escribe $g\in L^1$ en su teorema de convergencia dominada en $L^p$. Para $p>1$ eso no basta. En $(0,1)$, sea $g(x)=x^{-1/p}$ y $f_n=\min\{g,n\}$. Cada $f_n$ pertenece a $L^p$, $g\in L^1$ y $f_n\to g$, pero $g\notin L^p$ porque $\int_0^1x^{-1}dx=\infty$. La condición adecuada es $g\in L^p$, o equivalentemente una cota integrable para $|f_n|^p$.
\end{observacion}
\begin{teorema}[Vitali en $L^p$]\label{v2:vitali-p}
Sean $f_n,f\in L^p$, $1\le p<\infty$. Entonces $f_n\to f$ en norma si y solo si:
\begin{enumerate}
\item $f_n\to f$ en medida;
\item para cada $\eps>0$ existe $A\in\Acal$, $\mu(A)<\infty$, tal que $\int_{X\setminus A}|f_n|^p\,d\mu<\eps$ para todo $n$;
\item para cada $\eps>0$ existe $\delta>0$ tal que $\mu(B)<\delta$ implica $\int_B|f_n|^p\,d\mu<\eps$ para todo $n$.
\end{enumerate}
La condición (ii) equivale a exigir la misma cota para todo $B$ medible disjunto de $A$.
\end{teorema}
\begin{proof}
Una función $h\in L^1$, $h\ge0$, tiene integral pequeña en conjuntos de medida pequeña: se elige $M$ con $\int_{\{h>M\}}h<\eps/2$ y después $\mu(B)<\eps/(2M)$. También concentra su integral en algún conjunto de medida finita: $\{h>1/m\}$ tiene medida finita y las integrales sobre sus complementos tienden a cero. Estas dos observaciones se aplican a $|f|^p$ y a una cantidad finita de $|f_n|^p$.

Si $\norm{f_n-f}_p\to0$, usamos $|f_n|^p\le2^{p-1}(|f_n-f|^p+|f|^p)$. Para todos los índices suficientemente grandes, el primer término tiene integral tan pequeña como se desee en todo $X$; el segundo tiene las dos propiedades anteriores. Los índices iniciales se incorporan tomando una unión finita de los conjuntos de concentración y el mínimo de sus umbrales de medida. La convergencia en medida ya fue probada.

Recíprocamente, pongamos $h_n=|f_n-f|^p$. Por $h_n\le2^{p-1}(|f_n|^p+|f|^p)$, la familia $h_n$ satisface (ii) y (iii). Fijado $\eps>0$, elegimos $A$ de medida finita con $\int_{X\setminus A}h_n<\eps/3$ para todo $n$. Elegimos $\delta$ con $\mu(B)<\delta\Rightarrow\int_Bh_n<\eps/3$, y $a>0$ con $a^p\mu(A)<\eps/3$. La convergencia en medida implica que $B_n=A\cap\{|f_n-f|>a\}$ tiene medida menor que $\delta$ eventualmente. Entonces
\[
 \int_Xh_n\le\int_{X\setminus A}h_n+\int_{B_n}h_n+a^p\mu(A)<\eps.
\]
Cada una de las tres hipótesis controla una parte distinta de la integral.
\end{proof}
\begin{observacion}[El caso hilbertiano]
El ejercicio \ref{ej:v2-l18-6} y la completitud muestran que $L^2$ es un espacio de Hilbert. El teorema de representación de Hilbert del volumen I identifica su dual mediante $\varphi(g)=\langle g,h\rangle$. Todo operador compacto $T$ entre espacios de Hilbert es límite en norma de operadores de rango finito: una red finita de radio $\eps$ para $T$ de la bola unitaria genera un subespacio finito-dimensional $E$; si $P_E$ es la proyección ortogonal, entonces $\norm{T-P_ET}\le\eps$. No se necesita aquí separabilidad: la compacidad de esa imagen ya proporciona la red finita requerida.
\end{observacion}

\section{Ejercicios resueltos: traslaciones y geometría}
\begin{ejercicio}[Lectura 19, ejercicio 1]\label{ej:v2-l19-1}
Para $f\in L^p(\R)$, $1<p<\infty$, sea $f_x(t)=f(t-x)$. Demostrar $\norm{f_x-f_{x_0}}_p\to0$ cuando $x\to x_0$.
\end{ejercicio}
\begin{solucion}
La invariancia de la medida de Lebesgue da $\norm{h_x}_p=\norm h_p$. Si $h=\ind{(a,b)}$ con $a<b$ finitos, la diferencia de dos traslaciones tiene potencia $p$ igual al indicador de la diferencia simétrica de los intervalos, cuya medida es a lo sumo $2|x-x_0|$. La afirmación sigue para combinaciones finitas de estos indicadores por Minkowski.

Esas combinaciones son densas en $L^p(\R)$. En efecto, ya sabemos aproximar por indicadores de conjuntos $A$ de medida finita. La regularidad de Lebesgue proporciona $K\subset A\subset U$, con $K$ compacto, $U$ abierto y $m(U\setminus K)$ arbitrariamente pequeño. Finitos intervalos acotados contenidos en $U$ cubren $K$; su unión $J$ satisface $A\mathbin\triangle J\subset U\setminus K$. Por tanto $\norm{\ind A-\ind J}_p$ es pequeño. Dado ahora $f$, elegimos una combinación $h$ con $\norm{f-h}_p<\eps/3$. Entonces
\[
 \norm{f_x-f_{x_0}}_p\le2\norm{f-h}_p+\norm{h_x-h_{x_0}}_p<\eps
\]
para $x$ suficientemente próximo a $x_0$. Esta prueba vale también para $p=1$.
\end{solucion}
\begin{ejercicio}[Lectura 19, ejercicio 2: hipótesis necesaria]\label{ej:v2-l19-2}
El enunciado original afirma que para $p\ne2$ la norma de $L^p(X)$ no proviene de un producto interno. Examinar la afirmación y demostrarla cuando existen dos conjuntos medibles disjuntos de medida positiva finita.
\end{ejercicio}
\begin{solucion}
Sin esa condición la afirmación es falsa: en un espacio de un punto de medida uno, $L^p=\F$ y su norma es el valor absoluto para todo $p$. Supongamos ahora $A\cap B=\varnothing$ y $0<\mu(A),\mu(B)<\infty$. Sean $u=\mu(A)^{-1/p}\ind A$ y $v=\mu(B)^{-1/p}\ind B$. Sus normas son uno y $\norm{u+v}_p=\norm{u-v}_p=2^{1/p}$. La ley del paralelogramo exigiría
\[
 2\cdot2^{2/p}=2(1+1)=4,
\]
lo que equivale a $p=2$. Para $p\ne2$ se contradice una identidad necesaria de toda norma inducida por producto interno.
\end{solucion}
\section{Qué control aporta cada hipótesis}
La norma $L^p$ integra el error, no exige que sea pequeño en cada punto. La convergencia en medida detecta la frecuencia de los errores grandes; la uniformidad de las integrales y la concentración en conjuntos de medida finita impiden que esos errores acumulen masa o escapen al infinito. Las clases de equivalencia eliminan únicamente diferencias nulas: todas las construcciones puntuales anteriores se hicieron con representantes medibles y conjuntos excepcionales medibles identificados explícitamente.

\chapter{Dualidad de \texorpdfstring{$L^p$ para $1<p<\infty$}{Lp para 1<p<infinito}}\label{v2:cap10}\label{cap:v2-10}
\section{De Hölder a la representación de funcionales}\label{lectura:v2-20}
Hölder construye un funcional continuo a partir de una función integrable. La pregunta inversa es más profunda: ¿todo funcional tiene ese origen? Las lecturas 20 y 21 de \cite{valencia2} responden afirmativamente mediante la geometría de $L^p$, sin imponer finitud ni $\sigma$-finitud a la medida. Mantendremos el emparejamiento bilineal $\int fg$; por eso las conjugaciones aparecerán en las funciones que realizan la norma y no en la definición del funcional. Esta convención hace lineal la identificación del dual.

\begin{proposicion}\label{v2:dual-isometria}
Sean $1<p<\infty$ y $q=p/(p-1)$. Si $h\in L^q$, entonces $\Phi(h)(f)=\int_Xhf\,d\mu$ define un funcional lineal continuo sobre $L^p$ y $\norm{\Phi(h)}=\norm h_q$.
\end{proposicion}
\begin{proof}
Hölder prueba $|\Phi(h)(f)|\le\norm h_q\norm f_p$. Para la desigualdad contraria, si $h\ne0$ en $L^q$, ponemos $f=\overline h|h|^{q-2}$ donde $h\ne0$ y $f=0$ donde $h=0$. Como $(q-1)p=q$, tenemos $|f|^p=|h|^q$ y
\[
 \norm f_p=\norm h_q^{q/p}=\norm h_q^{q-1},\qquad
 \Phi(h)(f)=\int|h|^q=\norm h_q^q.
\]
El cociente entre el valor del funcional y la norma de $f$ es $\norm h_q$. Si $h=0$, ambas normas son cero. La definición por separado en los ceros evita escribir una potencia negativa de cero.
\end{proof}

\section{Las desigualdades escalares de Clarkson}
\begin{lema}\label{v2:clarkson-escalar-mayor}
Si $p\ge2$, entonces
\[
 \left(\frac{1+x}{2}\right)^p+\left(\frac{1-x}{2}\right)^p\le\frac{1+x^p}{2}\quad(0\le x\le1),
\]
y, para $z,w\in\C$,
\[
 \left|\frac{z+w}{2}\right|^p+\left|\frac{z-w}{2}\right|^p
 \le\frac{|z|^p+|w|^p}{2}.
\]
\end{lema}
\begin{proof}
Para $0<x\le1$, sea
\[
 F(x)=x^{-p}\bigl((1+x)^p+(1-x)^p-2^{p-1}(1+x^p)\bigr).
\]
La regla del producto da
\[
 F'(x)=\frac p{x^{p+1}}\bigl(2^{p-1}-(1+x)^{p-1}-(1-x)^{p-1}\bigr).
\]
La suma $(1+x)^{p-1}+(1-x)^{p-1}$ es creciente, pues su derivada es $(p-1)((1+x)^{p-2}-(1-x)^{p-2})\ge0$, y en $x=1$ vale $2^{p-1}$. Por tanto $F'\ge0$, y $F(x)\le F(1)=0$. El extremo $x=0$ se verifica directamente. Dividir entre $2^p$ da la primera desigualdad.

Si $z=0$ o $w=0$, la segunda se reduce a $2^{1-p}\le1/2$. En otro caso, por simetría suponemos $|w|\le|z|$ y escribimos $w/z=re^{i\theta}$, $0<r\le1$. La expresión que hay que maximizar, salvo el factor $|z|^p/2^p$, es
\[
 H(c)=(1+r^2+2rc)^{p/2}+(1+r^2-2rc)^{p/2},\qquad c=\cos\theta.
\]
Es par en $c$ y creciente para $0\le c\le1$, porque
$H'(c)=pr((1+r^2+2rc)^{p/2-1}-(1+r^2-2rc)^{p/2-1})\ge0$.
Su máximo está en $|c|=1$, donde se aplica la primera desigualdad con $x=r$. Los exponentes $p$ en los dos sumandos y el signo menos entre las potencias de la derivada son necesarios.
\end{proof}
\begin{lema}\label{v2:clarkson-escalar-menor}
Si $1<p\le2$ y $q=p/(p-1)$, entonces
\[
 (1+x)^q+(1-x)^q\le2(1+x^p)^{q/p}\quad(0\le x\le1),
\]
y $|z+w|^q+|z-w|^q\le2(|z|^p+|w|^p)^{q/p}$ para $z,w\in\C$.
\end{lema}
\begin{proof}
Para $p=2$ hay igualdad. Supongamos $1<p<2$. Los extremos $x=0,1$ se verifican directamente. Sustituyendo $x=(1-u)/(1+u)$, $0<u<1$, y elevando a $p-1$, la primera desigualdad equivale a
\[
 (1+u^q)^{p-1}\le\tfrac12((1+u)^p+(1-u)^p).
\]
Las series binomiales convergen absolutamente para $0<u<1$. Al restarlas y agrupar el término de grado par de la primera serie con dos términos consecutivos de la segunda, la diferencia es $\sum_{k\ge1}T_k$, donde
\[
 T_k=\binom p{2k}u^{2k}-\binom{p-1}{2k-1}u^{q(2k-1)}-\binom{p-1}{2k}u^{2kq}.
\]
La fórmula $\binom ak=a(a-1)\cdots(a-k+1)/k!$ permite factorizar
\[
 T_k=\frac{u^{2k}}{(2k-1)!}\prod_{j=2}^{2k}(j-p)
 \left(\frac{1-u^{a_k}}{a_k}-\frac{1-u^{b_k}}{b_k}\right),
 \quad a_k=\frac{2k-p}{p-1},\quad b_k=\frac{2k}{p-1}.
\]
Para comprobar la factorización se expanden los dos cocientes: el término constante es $p(p-1)/(2k(2k-p))$, y los otros coeficientes son $-(p-1)/(2k-p)$ y $(p-1)/(2k)$. Son exactamente los tres coeficientes binomiales después de extraer el producto escrito. Los factores exteriores son positivos. Además $0<a_k<b_k$ y
\[
 \frac{1-u^a}{a}=\int_u^1t^{a-1}\,dt
\]
es decreciente en $a>0$, porque $0<t<1$. Así $T_k\ge0$, y la suma es no negativa. No se requiere afirmar convergencia uniforme de estas series en los extremos.

Para números complejos se reduce por homogeneidad a $z=1$, $w=re^{i\theta}$, $0\le r\le1$. Como $q\ge2$, el argumento de $H(c)$ en el lema anterior, ahora con exponente $q$, prueba que el máximo se alcanza cuando los dos números están alineados. La primera desigualdad del presente lema concluye.
\end{proof}
\begin{lema}[Desigualdad inversa para funciones no negativas]\label{v2:minkowski-inversa}
Si $0<r\le1$ y $a,b\ge0$ son medibles con integrales de sus potencias finitas, entonces
\[
 \left(\int(a+b)^r\right)^{1/r}\ge
 \left(\int a^r\right)^{1/r}+\left(\int b^r\right)^{1/r}.
\]
\end{lema}
\begin{proof}
Llamemos $A,B$ a los dos términos de la derecha. Si alguno es cero, su función es cero casi en todas partes. En otro caso, por concavidad de $t^r$,
\[
 \left(\frac{a+b}{A+B}\right)^r\ge
 \frac A{A+B}\left(\frac aA\right)^r+
 \frac B{A+B}\left(\frac bB\right)^r.
\]
Al integrar, el lado derecho vale uno. Tomar raíces prueba la afirmación. Esta es una desigualdad inversa para funciones no negativas; no es la desigualdad triangular de una norma cuando $r<1$.
\end{proof}
\begin{teorema}[Clarkson]\label{v2:clarkson}\index{Clarkson, desigualdades de}
Para $f,g\in L^p$, se cumplen
\begin{align*}
 p\ge2:&\quad\left\|\frac{f+g}{2}\right\|_p^p+\left\|\frac{f-g}{2}\right\|_p^p
 \le\frac{\norm f_p^p+\norm g_p^p}{2},\\
 1<p\le2:&\quad\left\|\frac{f+g}{2}\right\|_p^q+\left\|\frac{f-g}{2}\right\|_p^q
 \le\left(\frac{\norm f_p^p+\norm g_p^p}{2}\right)^{q/p}.
\end{align*}
\end{teorema}
\begin{proof}
La primera desigualdad se obtiene integrando el lema \ref{v2:clarkson-escalar-mayor}. Para la segunda, ponemos $r=p-1$, $u=|(f+g)/2|^q$ y $v=|(f-g)/2|^q$. Como $qr=p$, el lema \ref{v2:minkowski-inversa} da
\[
 \left\|\frac{f+g}{2}\right\|_p^q+\left\|\frac{f-g}{2}\right\|_p^q
 \le\left(\int(u+v)^{p-1}\,d\mu\right)^{1/(p-1)}.
\]
El lema escalar menor implica $u+v\le2^{1-q}(|f|^p+|g|^p)^{1/(p-1)}$. Elevando a $p-1$, el factor es $2^{(1-q)(p-1)}=1/2$. Integrar produce la cota anunciada.
\end{proof}

\section{Un funcional alcanza su norma}\label{lectura:v2-21}
\begin{teorema}\label{v2:lp-alcanza}
Todo funcional continuo no nulo $\varphi\in(L^p)^*$, $1<p<\infty$, alcanza su norma: existe $v$ con $\norm v_p=1$ y $\varphi(v)=\norm\varphi$.
\end{teorema}
\begin{proof}
Elegimos $u_n$ de norma uno con $|\varphi(u_n)|\to\norm\varphi$ y valores no nulos. Multiplicando $u_n$ por $\overline{\varphi(u_n)}/|\varphi(u_n)|$ obtenemos $v_n$ de norma uno con $\varphi(v_n)=|\varphi(u_n)|>0$. Es esta fase conjugada la que elimina el argumento complejo.

Si $(v_n)$ no fuera de Cauchy, existirían $\eps>0$ y pares de índices $a_j,b_j\to\infty$ con $\norm{v_{a_j}-v_{b_j}}_p\ge\eps$. Clarkson, con $s=p$ si $p\ge2$ y $s=q$ si $p\le2$, implica
\[
 \left\|\frac{v_{a_j}+v_{b_j}}2\right\|_p\le
 \left(1-(\eps/2)^s\right)^{1/s}=c<1.
\]
Pero $\varphi((v_{a_j}+v_{b_j})/2)\to\norm\varphi$, mientras que su valor absoluto es a lo sumo $c\norm\varphi$. Es una contradicción. Por completitud, $v_n\to v$; la continuidad de la norma y de $\varphi$ da las dos igualdades requeridas.
\end{proof}
\begin{lema}[McShane]\label{v2:mcshane}
Sean $V$ normado complejo y $\varphi\ne0$ continuo, con $\norm v=1$ y $\varphi(v)=\norm\varphi$. Si $\psi_w(t)=\norm{v+tw}$ y $\psi_{-iw}$ son diferenciables en $0$, entonces
\[
 \frac{\varphi(w)}{\norm\varphi}=\psi_w'(0)+i\psi_{-iw}'(0).
\]
En el caso real basta la primera derivada.
\end{lema}
\begin{proof}
Normalizamos $\norm\varphi=1$. Para todo $t\in\R$, $\norm{v+tw}\ge\operatorname{Re}\varphi(v+tw)=1+t\operatorname{Re}\varphi(w)$. Restar uno y dividir por $t$ da una cota inferior si $t>0$ y superior si $t<0$. La existencia de la derivada obliga a $\psi_w'(0)=\operatorname{Re}\varphi(w)$. Sustituir $w$ por $-iw$ da $\psi_{-iw}'(0)=\operatorname{Im}\varphi(w)$. Esta prueba evita una factorización con denominadores complejos innecesaria para el argumento.
\end{proof}
\begin{teorema}[Representación y reflexividad de $L^p$]\label{v2:dual-lp}\index{reflexividad!de Lp@de $L^p$}
Para cualquier espacio de medida y $1<p<\infty$, la aplicación $\Phi:L^q\to(L^p)^*$, $\Phi(h)(f)=\int hf\,d\mu$, es una isometría lineal sobreyectiva. En consecuencia, $L^p$ es reflexivo.
\end{teorema}
\begin{proof}
Ya se probó que es una isometría. Sea $\varphi\ne0$ y elijamos $v$ como en el teorema \ref{v2:lp-alcanza}. Para $w\in L^p$, estudiamos $I(t)=\int|v+tw|^p$. La función escalar $t\mapsto|a+tb|^p$ tiene derivada
\[
 p|a+tb|^{p-2}\operatorname{Re}(\overline{a+tb}\,b),
\]
interpretada como cero cuando $a+tb=0$. En efecto, allí el cociente incremental es $|b|^p|h|^p/h\to0$ porque $p>1$. Para $0<|t|\le1$, el teorema del valor medio da la cota
\[
 \left|\frac{|v+tw|^p-|v|^p}{t}\right|
 \le p(|v|+|w|)^{p-1}|w|.
\]
El lado derecho es integrable por Hölder, pues $(|v|+|w|)^{p-1}\in L^q$. Aplicamos convergencia dominada a los cocientes incrementales mismos, que son medibles; no hace falta elegir mediblemente un punto intermedio del teorema del valor medio. Así
\[
 I'(0)=p\int|v|^{p-2}\operatorname{Re}(\overline v w)\,d\mu.
\]
Como $I(0)=1$, la regla de la cadena da
$\psi_w'(0)=\int|v|^{p-2}\operatorname{Re}(\overline v w)\,d\mu$.
McShane recupera las partes real e imaginaria y obtiene
\[
 \varphi(w)=\norm\varphi\int |v|^{p-2}\overline v w\,d\mu.
\]
La función $h=\norm\varphi\,|v|^{p-2}\overline v$, definida como cero en los ceros de $v$, pertenece a $L^q$ porque $(p-1)q=p$. Representa a $\varphi$; el funcional cero se representa por $h=0$.

Finalmente, si $\Lambda\in(L^p)^{**}$, el funcional $h\mapsto\Lambda(\Phi(h))$ sobre $L^q$ está representado por un $f\in L^p$, aplicando el resultado con exponente $q$. Por tanto $\Lambda(\Phi(h))=\int fh=\Phi(h)(f)$ para todo $h$. Como todos los funcionales de $(L^p)^*$ son de la forma $\Phi(h)$, se sigue $\Lambda=J(f)$, donde $J$ es el encaje canónico en el bidual. Esta es la reflexividad.
\end{proof}

\section{Ejercicios resueltos: geometría, multiplicadores y compacidad}
\begin{ejercicio}[Lectura 21, ejercicio 1: precisión de la igualdad]\label{ej:v2-l21-1}
Para $1<p<\infty$ y $f,g\in L^p$ de normas no nulas, demostrar
\[
 2^{-|2-p|/p}\le
 \frac{\norm{f+g}_p^2+\norm{f-g}_p^2}{2(\norm f_p^2+\norm g_p^2)}
 \le2^{|2-p|/p}.
\]
Encontrar pares que alcancen ambos extremos cuando existen dos conjuntos medibles \emph{disjuntos} de medida positiva finita. La disjunción falta en la condición original para esta última petición.
\end{ejercicio}
\begin{solucion}
Escribamos $r=p/2$ y, para $h\ge0$, $Q_r(h)=(\int h^r)^{1/r}$. Si $r\ge1$, Minkowski y $a^r+b^r\le(a+b)^r$ dan
\[
 Q_r(a+b)\le Q_r(a)+Q_r(b)\le2^{1-1/r}Q_r(a+b).
\]
La segunda cota usa primero la desigualdad entre normas de un par de números: $A+B\le2^{1-1/r}(A^r+B^r)^{1/r}$. Si $0<r\le1$, la desigualdad inversa y $a^r+b^r\le2^{1-r}(a+b)^r$ dan
\[
 Q_r(a)+Q_r(b)\le Q_r(a+b)\le2^{1/r-1}(Q_r(a)+Q_r(b)).
\]
Para la última desigualdad se usa $(a+b)^r\le a^r+b^r$ y la convexidad de $t^{1/r}$ sobre los dos números $Q_r(a)^r,Q_r(b)^r$.

Aplicamos estas comparaciones a $a=|f+g|^2$, $b=|f-g|^2$ y usamos la identidad escalar $a+b=2(|f|^2+|g|^2)$. Resulta
\[
 \norm{f+g}_p^2+\norm{f-g}_p^2
 \le2C(\norm f_p^2+\norm g_p^2),\qquad C=2^{|1-2/p|}.
\]
Sustituyendo ahora $(f,g)$ por $(f+g,f-g)$ en esta cota se obtiene la desigualdad inferior, pues las nuevas suma y diferencia son $2f,2g$.

Sean $u,v$ indicadores normalizados de dos conjuntos disjuntos de medida positiva finita. Para $(f,g)=(u,v)$ el cociente es $2^{2/p-1}$. Para $(f,g)=(u+v,u-v)$ es $2^{1-2/p}$. Son los dos extremos, cuyo orden depende de si $p<2$ o $p>2$. Si $p=2$, ambos valen uno.
\end{solucion}
\begin{ejercicio}[Lectura 21, ejercicio 2]\label{ej:v2-l21-2}
Sea $f:X\to\F$ medible. Supóngase que $\{f\ne0\}$ es de medida $\sigma$-finita y que $fg\in L^1$ para todo $g\in L^p$, $1<p<\infty$. Demostrar $f\in L^q$, $q=p/(p-1)$.
\end{ejercicio}
\begin{solucion}
Elegimos conjuntos crecientes $E_n$ de medida finita que cubran $\{f\ne0\}$ y ponemos $f_n=f\ind{E_n\cap\{|f|\le n\}}$. Entonces $f_n\in L^q$, $|f_n|\uparrow|f|$, y los funcionales $T_n(g)=\int f_ng$ tienen norma $\norm{f_n}_q$ por la proposición \ref{v2:dual-isometria}. Para cada $g\in L^p$,
$|T_n(g)|\le\int|fg|<\infty$, uniformemente en $n$. El principio de acotación uniforme, aplicado al Banach $L^p$, da $\sup_n\norm{f_n}_q=M<\infty$. La convergencia monótona implica $\int|f|^q=\lim_n\int|f_n|^q\le M^q$. Así $f\in L^q$.
\end{solucion}
\begin{ejercicio}[Lectura 21, ejercicio 3]\label{ej:v2-l21-3}
Supóngase que $\mu$ es semifinita: cada conjunto de medida infinita contiene uno de medida positiva finita. Si $f$ es medible y $fg\in L^1$ para todo $g\in L^p$, demostrar $f\in L^q$.
\end{ejercicio}
\begin{solucion}
Primero observemos que, en un conjunto $A$ de medida infinita, las medidas de sus subconjuntos de medida finita no pueden estar acotadas. Si su supremo fuera $M<\infty$, una sucesión que se acercara a $M$ tendría unión $B$ de medida $M$, tomando uniones finitas crecientes. El resto $A\setminus B$ tendría medida infinita y, por semifinitud, contendría un $D$ de medida positiva finita. Entonces $B\cup D$ contradice la definición de $M$. Por tanto podemos elegir $B_n\subset A$ de medida finita con $\mu(B_n)>n$; su unión es $\sigma$-finita e infinita.

Para $A_k=\{|f|\ge1/k\}$, si $\mu(A_k)=\infty$, elegimos un $B\subset A_k$ de medida $\sigma$-finita e infinita. La función $f\ind B$ satisface el ejercicio \ref{ej:v2-l21-2}, luego pertenece a $L^q$, pero $\int_B|f|^q\ge k^{-q}\mu(B)=\infty$, contradicción. Así todos los $A_k$ tienen medida finita. Su unión es $\{f\ne0\}$, y el ejercicio anterior se aplica directamente a $f$.
\end{solucion}
\begin{ejercicio}[Lectura 21, ejercicio 4]\label{ej:v2-l21-4}
En $\R$, sea $\Acal$ la $\sigma$-álgebra de conjuntos numerables o con complemento numerable, y sea $\mu(A)=0$ en el primer caso y $\mu(A)=1$ en el segundo. Explicar por qué la petición de alcanzar ambos extremos del ejercicio \ref{ej:v2-l21-1} falla en este espacio para $p\ne2$.
\end{ejercicio}
\begin{solucion}
Toda función medible real es constante fuera de un conjunto numerable. Para verlo, dividimos su recorrido en intervalos de longitud $2^{-n}$. De sus preimágenes, exactamente una puede ser conumerable y al menos una debe serlo, pues una unión numerable de conjuntos numerables no cubre $\R$. Fuera de la unión numerable de las preimágenes descartadas, los valores de la función están a distancia a lo sumo $2^{-n}$ para todo $n$, luego son iguales. Para funciones complejas se aplica el argumento a sus partes real e imaginaria. Por tanto cada $L^p$ es isométrico a $\F$ y el cociente del ejercicio es siempre uno, por el paralelogramo escalar. Las desigualdades siguen siendo verdaderas, pero ninguno de sus extremos estrictos se alcanza si $p\ne2$. Dos conjuntos de medida positiva pueden coincidir o intersectarse en medida uno; su mera existencia no sustituye la disjunción.
\end{solucion}
\begin{ejercicio}[Lectura 21, ejercicio 5]\label{ej:v2-l21-5}
Una sucesión $v_n$ converge débilmente a $v$ si $\varphi(v_n)\to\varphi(v)$ para todo $\varphi\in V^*$. Un operador continuo $T:V\to W$ es completamente continuo si transforma convergencia débil en convergencia en norma. Para espacios de Banach $V,W$, demostrar: (i) todo operador compacto es completamente continuo; (ii) si $V$ es reflexivo, todo operador completamente continuo es compacto.
\end{ejercicio}
\begin{solucion}
(i) Una sucesión débilmente convergente es acotada: se aplica acotación uniforme a los funcionales $J(v_n)$ sobre $V^*$ y se usa $\norm{J(v_n)}=\norm{v_n}$, consecuencia de Hahn--Banach. Si $T$ es compacto, toda subsucesión de $Tv_n$ admite una subsucesión convergente en norma. Su límite debe ser $Tv$, porque para $\psi\in W^*$, $\psi(Tv_n)=(T^*\psi)(v_n)\to\psi(Tv)$, y los funcionales separan puntos por Hahn--Banach. Si $Tv_n$ no convergiera en norma a $Tv$, una subsucesión quedaría a distancia al menos $\eps>0$ y su subsucesión convergente produciría una contradicción.

(ii) Justificamos el hecho de extracción que se necesita. En un Banach reflexivo toda sucesión acotada tiene una subsucesión débilmente convergente. Sea $Y$ el subespacio cerrado generado por los términos de la sucesión; es separable. También es reflexivo: dado $F\in Y^{**}$, la composición de $F$ con la restricción $V^*\to Y^*$ pertenece a $V^{**}$, luego es evaluación en algún $v\in V$. Todo funcional que se anula en $Y$ se anula en $v$; la separación de un punto de un subespacio cerrado mediante Hahn--Banach implica $v\in Y$. Así $F$ es evaluación en ese $v$.

El dual $Y^*$ es separable. En efecto, elijamos $y_j$ densos en la esfera de $Y$ y funcionales $\varphi_j$ de norma uno con $\varphi_j(y_j)=1$. La familia es normante: si $y\ne0$, aproximar $y/\norm y$ por $y_j$ muestra $\sup_j|\varphi_j(y)|=\norm y$. Si el espacio cerrado generado por los $\varphi_j$ no fuera $Y^*$, Hahn--Banach daría un funcional no nulo de $Y^{**}$ que se anula sobre todos ellos; por reflexividad sería evaluación en un $y\ne0$, contradiciendo que la familia sea normante. Las combinaciones racionales de los $\varphi_j$ forman, pues, un conjunto denso numerable de $Y^*$.

Una extracción diagonal hace converger las evaluaciones de una sucesión acotada sobre ese conjunto numerable. La cota uniforme de las normas extiende la convergencia a todo $\varphi\in Y^*$ y define un funcional continuo $F(\varphi)=\lim\varphi(v_{n_j})$. Reflexividad da $F=J(y)$, es decir, convergencia débil a $y$. La restricción de cualquier funcional de $V^*$ muestra que también es débil en $V$.

Aplicamos este hecho a cualquier sucesión de la bola unitaria de $V$. Su subsucesión débilmente convergente tiene, por completa continuidad, imagen convergente en norma. Así la imagen de la bola es relativamente secuencialmente compacta; en un espacio métrico completo eso equivale a compacidad de su cerradura. Por tanto $T$ es compacto.

La lectura menciona además un resultado atribuido a Odell acerca de espacios que no contienen $\ell^1$. Esa afirmación pertenece a la teoría de sucesiones débiles; no se usa como hipótesis ni como paso de estas dos demostraciones.
\end{solucion}
\begin{ejercicio}[Lectura 21, ejercicio 6]\label{ej:v2-l21-6}
Sean $1<p<\infty$, $q=p/(p-1)$ y $\mu$ $\sigma$-finita. Sea $k:X\times X\to\F$ medible para $\Acal\otimes\Acal$, con
\[
 \alpha=\int_X\left(\int_X|k(x,y)|^p\,d\mu(x)\right)^{q/p}d\mu(y)<\infty.
\]
Demostrar que $Kf(x)=\int_Xk(x,y)f(y)\,d\mu(y)$ define, casi en todas partes, un operador compacto en $L^p$, con $\norm K\le\alpha^{1/q}$.
\end{ejercicio}
\begin{solucion}
Escribamos $a(y)=\norm{k(\cdot,y)}_p$, de modo que $\norm a_q=\alpha^{1/q}$. La desigualdad integral de Minkowski da
\[
 \norm{Kf}_p\le\int_X a(y)|f(y)|\,d\mu(y)\le\norm a_q\norm f_p. \tag{*}
\]
Aquí también hay que justificar existencia. Para $H(x)=\int|k(x,y)f(y)|\,d\mu(y)$, Tonelli y Hölder implican $\int Hh\le\norm a_q\norm f_p\norm h_q$ para todo $h\ge0$ de $L^q$. Aplicamos esta desigualdad a las funciones que realizan la norma de $H_m=\min\{H,m\}\ind{E_m}$, donde $E_m\uparrow X$ y $\mu(E_m)<\infty$. Obtenemos $\norm{H_m}_p\le\norm a_q\norm f_p$; por convergencia monótona $H\in L^p$ con la misma cota. Así la integral definitoria es absolutamente convergente casi en todas partes y (*) queda probada. También queda claro que cambiar $f$ en un conjunto nulo no cambia $Kf$ en $L^p$.

Demostraremos compacidad aproximando el núcleo. La norma mixta
\[
 \norm h_{q,p}=\left(\int\norm{h(\cdot,y)}_p^q\,d\mu(y)\right)^{1/q}
\]
controla, por (*), la norma del operador asociado. Los núcleos $k_m=k\ind{E_m\times E_m}\ind{\{|k|\le m\}}$ convergen a $k$ en esta norma: primero, para casi todo $y$, convergencia dominada en $x$ da convergencia de las normas $L^p$ del error; después dominada en $y$, con cota $a(y)^q$, da la afirmación.

En un rectángulo de medida finita, todo núcleo acotado puede aproximarse en $L^r$, $r\ge\max\{p,q\}$, por combinaciones finitas de indicadores de rectángulos medibles. Explicamos el punto sobre conjuntos: la familia de conjuntos aproximables en medida por el álgebra de uniones finitas de rectángulos contiene esa álgebra, es estable por complementos relativos y por uniones numerables. Para esta última propiedad se aproxima primero una unión finita, pues la medida total es finita, y después cada uno de sus miembros. Por tanto contiene la $\sigma$-álgebra producto. Aproximar niveles de una función simple da la densidad en $L^r$. Además, si $h$ está soportado en $E\times F$ de medidas finitas positivas, dos aplicaciones de Hölder dan
\[
 \norm h_{q,p}\le\mu(E)^{1/p-1/r}\mu(F)^{1/q-1/r}\norm h_{L^r(E\times F)}.
\]
Los casos de medida cero son triviales. Así obtenemos núcleos $s_j(x,y)=\sum_{i=1}^{m_j}c_i\ind{A_i}(x)\ind{B_i}(y)$ que convergen a $k$ en norma mixta, con $A_i,B_i$ de medida finita. Su operador es
\[
 S_jf=\sum_i c_i\left(\int_{B_i}f\,d\mu\right)\ind{A_i},
\]
y tiene rango finito. Por (*), $\norm{K-S_j}\le\norm{k-s_j}_{q,p}\to0$. Un límite en norma de operadores compactos es compacto: una red finita para una aproximación, junto con el error uniforme sobre la bola unitaria, da acotación total de la imagen; la completitud del codominio da compacidad de la cerradura. Esto completa tanto la definición como la compacidad de $K$.
\end{solucion}
\section{Geometría y representación}
La convexidad cuantitativa de Clarkson obliga a que una sucesión casi maximizante sea de Cauchy. La diferenciación de la norma identifica después el funcional que apoya a la esfera en su límite. El argumento explica por qué el intervalo $1<p<\infty$ tiene una dualidad especialmente regular y por qué los extremos $p=1,\infty$ requieren estudiar propiedades adicionales de la medida.

\chapter{El espacio \texorpdfstring{$L^\infty$}{L infinito} y la dualidad de \texorpdfstring{$L^1$}{L1}}\label{v2:cap11}\label{cap:v2-11}
\section{Una cota que ignora conjuntos nulos}\label{lectura:v2-22}
La norma uniforme cambia si se modifica una función en un solo punto. Para trabajar con clases módulo igualdad casi en todas partes se necesita una cota esencial. Las lecturas 22 y 23 de \cite{valencia2} muestran que esta modificación produce un Banach, pero también que la representación de su dual depende de propiedades de la medida que no aparecen para $1<p<\infty$.
\begin{definicion}
Una función medible $f:X\to\F$ es esencialmente acotada si $\mu(\{|f|>M\})=0$ para algún $M<\infty$. Definimos
\[
 \norm f_\infty=\inf\{M\ge0:\mu(\{|f|>M\})=0\}.
\]
El espacio de estas funciones es $\mathcal L^\infty$; su cociente por igualdad casi en todas partes es $L^\infty$.\index{supremo esencial}
\end{definicion}
Siempre $\norm f_\infty\le\sup_X|f|$. La desigualdad puede ser estricta: la función de Thomae en $[0,1]$, igual a $1/q$ en los racionales irreducibles $a/q$ y cero en irracionales, tiene supremo uno y supremo esencial cero para Lebesgue. En cambio, para una función continua en $\R$ ambos supremos coinciden: si $|f(x_0)|>a$, la continuidad mantiene $|f|>a$ en un intervalo abierto de medida positiva.
\begin{proposicion}\label{v2:linfty-cota}
Si $f\in\mathcal L^\infty$, entonces $|f|\le\norm f_\infty$ casi en todas partes. La aplicación $\norm{\cdot}_\infty$ es una seminorma sobre $\mathcal L^\infty$, cuyo núcleo son las funciones nulas casi en todas partes; induce una norma sobre $L^\infty$.
\end{proposicion}
\begin{proof}
Para cada $n$ elegimos una cota esencial $M_n<\norm f_\infty+1/n$. Fuera de la unión numerable de los conjuntos nulos $\{|f|>M_n\}$ se cumple $|f|\le M_n$ para todo $n$, y al pasar al límite $|f|\le\norm f_\infty$. Por ello $|f+g|\le\norm f_\infty+\norm g_\infty$ fuera de un conjunto nulo, lo que da la desigualdad triangular. También $\norm{\alpha f}_\infty\le|\alpha|\norm f_\infty$; para $\alpha\ne0$ se aplica la misma desigualdad a $f=\alpha^{-1}(\alpha f)$ y se obtiene igualdad. El caso $\alpha=0$ se trata directamente. Finalmente, norma cero equivale a $\mu(\{|f|>0\})=0$. Todas estas afirmaciones son invariantes al cambiar el representante en un conjunto medible nulo.
\end{proof}
\begin{teorema}\label{v2:linfty-banach}
$L^\infty$ es un espacio de Banach. Las funciones simples son densas en su norma.
\end{teorema}
\begin{proof}
Sea $(f_n)$ de Cauchy y elijamos representantes. Quitamos el conjunto nulo que es unión de todos los conjuntos
\[
 \{|f_n|>\norm{f_n}_\infty\},\qquad
 \{|f_n-f_m|>\norm{f_n-f_m}_\infty\},\quad n,m\in\N.
\]
En su complemento, la sucesión es uniformemente de Cauchy y uniformemente acotada, porque las normas de una sucesión de Cauchy están acotadas. Tiene allí un límite uniforme $f$, que extendemos por cero al conjunto excepcional. Es medible como límite de las funciones medibles modificadas por cero en ese conjunto. La cota uniforme prueba $f\in\mathcal L^\infty$. Para $n,m$ grandes, $|f_n-f_m|\le\eps$ fuera del conjunto excepcional; al hacer $m\to\infty$, queda $|f_n-f|\le\eps$ allí. Por tanto $\norm{f_n-f}_\infty\le\eps$.

Para densidad, modificamos primero un representante para que $|f|\le M$ en todo $X$. Dividimos el intervalo $[-M,M]$, o el cuadrado $[-M,M]^2$ en el caso complejo, en una cantidad finita de celdas de diámetro menor que $\eps$. Sus preimágenes son medibles; asignamos a cada una un valor de la celda. La función simple resultante $s$ cumple $|f-s|<\eps$ en todo $X$. Es necesario aproximar el valor de $f$, no solamente $|f|$, para conservar los signos y argumentos complejos.
\end{proof}

\section{Ejercicios propuestos: tamaño esencial y límites de exponentes}
\begin{ejercicio}[Lectura 22, ejercicio 1]\label{ej:v2-l22-1}
Demostrar $\norm f_\infty=\inf\{\sup_{x\in A}|f(x)|:A\in\Acal,\ \mu(X\setminus A)=0\}$.
\end{ejercicio}

\begin{ejercicio}[Lectura 22, ejercicio 2]\label{ej:v2-l22-2}
Demostrar que $L^\infty\subset L^1$ si y solo si $\mu(X)<\infty$.
\end{ejercicio}

\begin{ejercicio}[Lectura 22, ejercicio 3]\label{ej:v2-l22-3}
Demostrar que $L^\infty(\R)$ con medida de Lebesgue no es separable, aunque esta medida sea separable.
\end{ejercicio}

\begin{ejercicio}[Lectura 22, ejercicio 4]\label{ej:v2-l22-4}
Si $f\in L^p$, $1\le p\le\infty$, y $g\in L^\infty$, demostrar $fg\in L^p$ y $\norm{fg}_p\le\norm f_p\norm g_\infty$.
\end{ejercicio}

\begin{ejercicio}[Lectura 22, ejercicio 5]\label{ej:v2-l22-5}
Si $\mu(X)<\infty$, demostrar que $f\in L^\infty$ equivale a: (i) $f\in L^p$ para todo $1\le p<\infty$; (ii) $\sup_{p\ge1}\norm f_p<\infty$. En tal caso, demostrar $\lim_{p\to\infty}\norm f_p=\norm f_\infty$.
\end{ejercicio}

\begin{ejercicio}[Lectura 22, ejercicio 6]\label{ej:v2-l22-6}
Si $f\in L^p(\R)$ para todo $p\ge p_0\ge1$ y además $f\in L^\infty(\R)$, ¿se cumple $\norm f_p\to\norm f_\infty$?
\end{ejercicio}

\begin{ejercicio}[Lectura 22, ejercicio 7]\label{ej:v2-l22-7}
Si una sucesión $x$ pertenece a $\ell^p$ para todo $p\ge p_0\ge1$, ¿se cumple $\norm x_p\to\norm x_\infty$?
\end{ejercicio}


\section{Los emparejamientos de los extremos}\label{lectura:v2-23}
Hölder en los extremos es el ejercicio \ref{ej:v2-l22-4} con $p=1$. Define aplicaciones lineales contractivas
\[
 \Phi:L^\infty\longrightarrow(L^1)^*,\quad\Phi(g)(f)=\int fg\,d\mu,
 \qquad
 \Psi:L^1\longrightarrow(L^\infty)^*,\quad\Psi(f)(g)=\int fg\,d\mu.
\]
En general no son ambas isometrías sobreyectivas. Los siguientes ejemplos separan tres dificultades distintas.
\begin{ejemplo}[Pérdida de inyectividad de $\Phi$]
En $\N$, definimos $\mu(A)=\#A$ si $1\notin A$ y $\mu(A)=\infty$ si $1\in A$. Es una medida: un conjunto que contiene $1$ tiene en cualquier partición disjunta exactamente una pieza que lo contiene, y fuera de $1$ se trata de conteo. Una función integrable debe anularse en $1$ y satisfacer $\sum_{n\ge2}|f(n)|<\infty$; por ello $L^1$ es isométrico a $\ell^1$ tras eliminar la primera coordenada. El único conjunto nulo es el vacío, así que $L^\infty=\ell^\infty$. Sin embargo, $g=\ind{\{1\}}$ tiene norma uno y $\Phi(g)=0$.
\end{ejemplo}
\begin{ejemplo}[Pérdida de sobreyectividad de $\Phi$]
En $X=(0,1)$ tomamos la $\sigma$-álgebra numerable-conumerable y medida de conteo. Toda función medible es constante fuera de un conjunto numerable, por el argumento del ejercicio \ref{ej:v2-l21-4}. Una función integrable tiene soporte numerable: para cada $n$, el conjunto donde $|f|\ge1/n$ debe ser finito. Por tanto
\[
 \varphi(f)=\sum_{x\in(0,1)}x f(x)
\]
está bien definida y $|\varphi(f)|\le\norm f_1$. El producto $x f(x)$ es medible porque está soportado en un conjunto numerable. Si $\varphi(f)=\int fg$ para una $g$ medible esencialmente acotada, al probar con $f=\ind{\{x\}}$ obtenemos $g(x)=x$ para todo $x$. Esa función no es medible para esta $\sigma$-álgebra, contradicción.
\end{ejemplo}
\begin{ejemplo}[Pérdida de sobreyectividad de $\Psi$]\label{v2:limite-extendido}
En el subespacio $c\subset\ell^\infty$ de sucesiones convergentes, el funcional $x\mapsto\lim_nx_n$ tiene norma uno. Hahn--Banach lo extiende a un funcional $L\in(\ell^\infty)^*$ de norma uno. Si $L(x)=\sum_nx_ny_n$ para algún $y\in\ell^1$, al evaluar en cada vector canónico $e_n$ se obtiene $y_n=L(e_n)=0$. Pero $L(1,1,\ldots)=1$, contradicción. Así la extensión existe, pero no está representada por una sucesión sumable.
\end{ejemplo}

\section{Semifinitud y localizabilidad}
\begin{definicion}
La medida $\mu$ es \emph{semifinita} si cada conjunto medible de medida infinita contiene uno de medida positiva finita. Un \emph{supremo esencial} de una familia $\mathcal E\subset\Acal$ es un $H\in\Acal$ tal que $\mu(E\setminus H)=0$ para todo $E\in\mathcal E$, y tal que cualquier $G$ con esa propiedad satisface $\mu(H\setminus G)=0$. La medida es \emph{localizable} si es semifinita y toda familia de conjuntos medibles admite un supremo esencial.\index{medida!semifinita}\index{medida!localizable}
\end{definicion}
\begin{proposicion}\label{v2:sigma-localizable}
Toda medida $\sigma$-finita es localizable.
\end{proposicion}
\begin{proof}
Escribamos $X=\bigcup_n E_n$ con $E_n$ disjuntos y de medida finita. Si $\mu(A)>0$, algún $A\cap E_n$ tiene medida positiva; si $\mu(A)=\infty$, esta observación da el subconjunto de medida positiva finita que requiere semifinitud. Definimos una medida finita
\[
 \lambda(B)=\sum_{n\ge1}\frac{2^{-n}}{1+\mu(E_n)}\mu(B\cap E_n).
\]
Tiene exactamente los mismos conjuntos nulos que $\mu$.

Dada $\mathcal E$, sea $\alpha$ el supremo de las medidas $\lambda$ de uniones numerables de miembros de $\mathcal E$, incluyendo la unión vacía. Elegimos uniones $H_n$ que se acerquen a $\alpha$ y ponemos $H=\bigcup_nH_n$. Sigue siendo una unión numerable de miembros de $\mathcal E$ y $\lambda(H)=\alpha$. Para $E\in\mathcal E$, $H\cup E$ es otra unión admisible, luego $\lambda(E\setminus H)=0$. Si $G$ contiene esencialmente a cada miembro, la unión numerable que forma $H$ satisface $\lambda(H\setminus G)=0$. Como los conjuntos nulos coinciden, $H$ es el supremo esencial para $\mu$.
\end{proof}
\begin{proposicion}\label{v2:l1-isometria}
$\Phi:L^\infty\to(L^1)^*$ es una isometría si y solo si $\mu$ es semifinita. Además $\Psi:L^1\to(L^\infty)^*$ siempre es una isometría.
\end{proposicion}
\begin{proof}
Si $\mu$ es semifinita y $0<a<\norm g_\infty$, el conjunto $\{|g|>a\}$ tiene medida positiva y contiene un $B$ de medida positiva finita. La función
$f=\mu(B)^{-1}\ind B\,\overline g/|g|$ tiene norma $L^1$ uno y $\int fg=\mu(B)^{-1}\int_B|g|>a$. Haciendo $a\uparrow\norm g_\infty$ obtenemos la isometría. Si falla semifinitud, hay un $A$ de medida infinita cuyos subconjuntos medibles de medida finita son nulos. Para $f\in L^1$, cada $A\cap\{|f|\ge1/n\}$ es finito en medida y por tanto nulo; así $f=0$ casi en todas partes en $A$. Se sigue $\Phi(\ind A)=0$ aunque $\norm{\ind A}_\infty=1$.

Para $f\in L^1$, la función $g=\overline f/|f|$ donde $f\ne0$, y cero en los demás puntos, tiene norma esencial a lo sumo uno y satisface $\int fg=\int|f|$. Esto da $\norm{\Psi(f)}\ge\norm f_1$; Hölder da la otra desigualdad.
\end{proof}
La afirmación completa es que $\Phi$ es una isometría sobreyectiva exactamente cuando $\mu$ es localizable. La demostración, incluida la necesidad de localizabilidad, se dará en el teorema \ref{v2:dual-l1-general}, una vez construidas las derivadas de Radon--Nikodym. Aquí ya se han probado sus ingredientes relativos a la norma y se han visto por qué no puede omitirse una condición sobre la medida.

\section{Ejercicios propuestos: funcionales y propiedades de la medida}
\begin{ejercicio}[Lectura 23, ejercicio 1]\label{ej:v2-l23-1}
Si $\mu$ es semifinita y $g$ es medible con $fg\in L^1$ para todo $f\in L^1$, demostrar $g\in L^\infty$.
\end{ejercicio}

\begin{ejercicio}[Lectura 23, ejercicio 2]\label{ej:v2-l23-2}
Para $f\in L^1$, construir $g\in L^\infty$ con $\norm g_\infty\le1$ y $\int fg=\norm f_1$.
\end{ejercicio}

\begin{ejercicio}[Lectura 23, ejercicio 3]\label{ej:v2-l23-3}
Si $\mu$ es semifinita y $g\in L^\infty$, demostrar
\[
 \norm g_\infty=\sup_{\norm f_1\le1}\left|\int fg\,d\mu\right|
 =\sup_{\norm f_1\le1}\norm{fg}_1.
\]
En el caso real se puede omitir el valor absoluto en el primer supremo; en el caso complejo un supremo de números complejos sin orden no está definido.
\end{ejercicio}

\begin{ejercicio}[Lectura 23, ejercicio 4]\label{ej:v2-l23-4}
Dar un espacio de medida finita y una $g\in L^\infty$ tal que $\norm{fg}_1<\norm g_\infty$ para toda $f\in L^1$ con $\norm f_1\le1$.
\end{ejercicio}

\begin{ejercicio}[Lectura 23, ejercicio 5]\label{ej:v2-l23-5}
Para un conjunto arbitrario $I$ con todas sus partes medibles y medida de conteo, demostrar $(\ell^1(I))^*\cong\ell^\infty(I)$ isométricamente.
\end{ejercicio}

\begin{ejercicio}[Lectura 23, ejercicio 6]\label{ej:v2-l23-6}
Supóngase $\mu$ $\sigma$-finita y $\varphi\in(L^1)^*$. Para cada $A$ de medida positiva finita supóngase dada $g_A\in L^\infty(A)$ que representa a $\varphi$ sobre funciones soportadas en $A$. Demostrar: (i) $g_A=g_B$ casi en todas partes en $A\cap B$; (ii) existe $g\in L^\infty(X)$ que representa a $\varphi$ en todo $L^1$.
\end{ejercicio}

\begin{ejercicio}[Lectura 23, ejercicio 7]\label{ej:v2-l23-7}
Sean $(X,\Acal,\mu)$, $(Y,\mathcal B,\nu)$ y $\phi:X\to Y$ medible, con $\mu(\phi^{-1}(A))=\nu(A)$ para todo $A\in\mathcal B$. Demostrar que $T_\phi(g)=g\circ\phi$ define una isometría lineal de $L^\infty(Y)$ en $L^\infty(X)$.
\end{ejercicio}

\section{Qué distingue a los exponentes extremos}
$L^\infty$ conserva el control uniforme después de eliminar un único conjunto nulo común a los datos numerables de una prueba. Su dual, sin embargo, contiene funcionales que no son integrales contra funciones de $L^1$. En la dirección contraria, semifinitud permite detectar la norma mediante funciones integrables y localizabilidad permitirá pegar representaciones construidas en conjuntos de medida finita.

\parte{Medidas con signo y representación}{La variación total controla medidas que admiten cancelaciones. Construiremos las descomposiciones de Hahn, Jordan y Lebesgue y la densidad de Radon--Nikodym. Después estudiaremos cómo unir representaciones locales y por qué el dual de $L^\infty$ incluye funciones de conjuntos sólo finitamente aditivas. Las hipótesis sobre la medida se mantendrán visibles.}
\chapter{Medidas con signo y medidas complejas}\label{v2:cap12}\label{cap:v2-12}
\section{Aditividad sin positividad}\label{lectura:v2-24}
Una integral indefinida $A\mapsto\int_A f\,d\mu$ puede tomar valores negativos o complejos. Sigue siendo aditiva, pero ya no es monótona. Las lecturas 24--26 de \cite{valencia2} explican cómo recuperar una medida positiva que controle todas sus oscilaciones: la variación total. Esta medida será la que determine qué funciones pueden integrarse y qué significa convergencia de medidas.
\begin{definicion}
Una \emph{medida con signo} sobre $(X,\Acal)$ es una función $\nu:\Acal\to[-\infty,\infty]$, con $\nu(\varnothing)=0$, numerablemente aditiva sobre familias disjuntas, que no toma simultáneamente los valores $+\infty$ y $-\infty$. Se exige que las sumas correspondientes estén definidas. Una \emph{medida compleja} toma valores en $\C$ y es numerablemente aditiva. En particular, todos sus valores son finitos.\index{medida!con signo}\index{medida!compleja}
\end{definicion}
Toda medida positiva es una medida con signo, pero no toda medida positiva es una medida compleja: para ello debe ser finita. La exclusión de los dos infinitos evita sumas indefinidas como $\infty-\infty$; no debe confundirse con la afirmación de que cada medida con signo tiene necesariamente un valor infinito.
\begin{proposicion}\label{v2:continuidad-signo}
Para una medida con signo $\nu$:
\begin{enumerate}
\item Si $B\subset A$ son medibles y $|\nu(A)|<\infty$, entonces $|\nu(B)|<\infty$.
\item Si $A_n\uparrow A$, entonces $\nu(A_n)\to\nu(A)$.
\item Si $A_n\downarrow A$ y algún $\nu(A_n)$ es finito, entonces $\nu(A_n)\to\nu(A)$.
\end{enumerate}
Para una medida compleja, sus partes real e imaginaria son medidas con signo finitas, y valen las mismas propiedades de continuidad, sin necesitar una hipótesis adicional de finitud.
\end{proposicion}
\begin{proof}
(i) En $\nu(A)=\nu(B)+\nu(A\setminus B)$ una suma definida solo puede ser finita si ambos términos lo son. (ii) Escribimos $A$ como unión disjunta de $A_1$ y $A_n\setminus A_{n-1}$, $n\ge2$. Las sumas parciales de la serie de aditividad son exactamente $\nu(A_n)$, y la definición de la suma da el límite. (iii) Descartamos los índices iniciales para tener $|\nu(A_1)|<\infty$. Los conjuntos $A_1\setminus A_n$ crecen a $A_1\setminus A$, de modo que (ii) y $\nu(A_n)=\nu(A_1)-\nu(A_1\setminus A_n)$ concluyen. Para medidas complejas, tomar partes real e imaginaria en una serie convergente prueba la aditividad de ambas; sus valores son reales finitos. Aplicamos los resultados anteriores a cada parte.
\end{proof}
\begin{ejemplo}[Integral indefinida]
Si $f\in L^1(\mu)$, la función $\nu_f(A)=\int_A f\,d\mu$ es una medida real con signo o compleja según $f$. Para conjuntos disjuntos $A_n$, los indicadores de sus uniones finitas convergen al de la unión total. Las funciones $f\ind{\cup_{j\le n}A_j}$ están dominadas por $|f|\in L^1$, por lo que convergencia dominada prueba la aditividad numerable. La medida del vacío es cero.
\end{ejemplo}
Las medidas con signo finitas forman un espacio vectorial real y las complejas forman uno complejo, con operaciones punto a punto. La aditividad de $a\nu+b\eta$ resulta de pasar al límite en las sumas parciales de las dos series convergentes. Esta restricción a medidas finitas en el caso real evita que una suma de medidas produzca $\infty-\infty$.

\section{Separación de las partes positiva y negativa}\label{lectura:v2-25}
\begin{definicion}
Un conjunto medible $P$ es $\nu$-positivo si $\nu(B)\ge0$ para todo $B\subset P$ medible. Análogamente $N$ es $\nu$-negativo si todos sus subconjuntos medibles tienen medida no positiva. Una \emph{descomposición de Hahn} es una partición medible $X=P\mathbin{\dot\cup}N$ de este tipo.\index{descomposicion de Hahn@descomposición de Hahn}
\end{definicion}
\begin{lema}\label{v2:subconjunto-positivo}
Si $0<\nu(A)<\infty$, existe un conjunto $\nu$-positivo $S\subset A$ con $\nu(S)>0$.
\end{lema}
\begin{proof}
Todos los subconjuntos medibles de $A$ tienen medida finita por la proposición \ref{v2:continuidad-signo}. Retiraremos subconjuntos de medida negativa. En la etapa $n$, mientras exista en el resto un conjunto $B$ con $\nu(B)<-1/n$, elegimos uno y lo retiramos. Esta etapa termina tras una cantidad finita de elecciones: si produjera infinitos conjuntos disjuntos de medida menor que $-1/n$, su unión, contenida en $A$, tendría medida $-\infty$ por aditividad, contradiciendo finitud.

La totalidad de los conjuntos retirados, al recorrer $n=1,2,\ldots$, forma una familia numerable disjunta de medidas negativas. Su unión $B$ tiene medida finita y no positiva. Pongamos $S=A\setminus B$. Entonces $\nu(S)=\nu(A)-\nu(B)\ge\nu(A)>0$. Si existiera $D\subset S$ medible con $\nu(D)<0$, podríamos escoger $n$ con $1/n<-\nu(D)$. El conjunto $D$ seguía disponible al terminar la etapa $n$, contradiciendo el criterio de terminación. Así $S$ es positivo.
\end{proof}
\begin{teorema}[Hahn]\label{v2:hahn}
Toda medida con signo admite una descomposición de Hahn. Si $(P,N)$ y $(P',N')$ son dos, entonces para todo $A\in\Acal$,
$\nu(A\cap P)=\nu(A\cap P')$ y $\nu(A\cap N)=\nu(A\cap N')$.
\end{teorema}
\begin{proof}
Podemos suponer que $\nu$ no toma el valor $+\infty$; en el caso contrario aplicamos el argumento a $-\nu$ e intercambiamos los conjuntos. Sea $\alpha=\sup\{\nu(E):E\text{ es positivo}\}$, inicialmente permitido infinito. Elegimos positivos $P_n$ cuyas medidas se acercan a $\alpha$ y ponemos $P=\bigcup_nP_n$. Una unión numerable de positivos es positiva: se disjunta la familia y se aplica aditividad a cualquier subconjunto medible de la unión. Así $\nu(P)\ge\nu(P_n)$ y $\nu(P)=\alpha$. Como $\nu$ no toma $+\infty$, $\alpha<\infty$.

Si $N=X\setminus P$ no fuera negativo, contendría un $A$ con $0<\nu(A)<\infty$. El lema anterior daría un positivo $S\subset A$ de medida positiva. Entonces $P\cup S$ sería positivo y tendría medida $\alpha+\nu(S)>\alpha$, contradicción. Esto prueba existencia.

Para unicidad, cualquier subconjunto medible de $P\cap N'$ tiene medida a la vez no negativa y no positiva, luego cero. Lo mismo ocurre en $P'\cap N$. Al descomponer $A\cap P$ y $A\cap P'$ según su parte común, sus diferencias tienen medida cero y las medidas coinciden. El argumento para las partes negativas es idéntico.
\end{proof}
\begin{definicion}[Jordan y variación real]
Dada una descomposición de Hahn, definimos
\[
 \nu^+(A)=\nu(A\cap P),\qquad \nu^-(A)=-\nu(A\cap N),\qquad |\nu|(A)=\nu^+(A)+\nu^-(A).
\]
Son medidas positivas: sus valores son no negativos y la aditividad numerable resulta de la de $\nu$ en cada parte. Además $\nu=\nu^+-\nu^-$, y al menos una de las dos medidas es finita, porque $\nu$ no toma ambos infinitos. La unicidad de Hahn demuestra que estas medidas no dependen de la partición elegida.\index{descomposicion de Jordan@descomposición de Jordan}\index{variacion total@variación total}
\end{definicion}
\begin{proposicion}\label{v2:variacion-real}
Para todo $A\in\Acal$,
\[
 |\nu|(A)=\sup\left\{\sum_{j=1}^m|\nu(A_j)|:A=\mathop{\dot\bigcup}_{j=1}^m A_j,\ A_j\in\Acal\right\}.
\]
\end{proposicion}
\begin{proof}
Para cada partición, $|\nu(A_j)|\le\nu^+(A_j)+\nu^-(A_j)$; al sumar se obtiene la cota $|\nu|(A)$. En la partición formada por $A\cap P$ y $A\cap N$, la suma es exactamente $\nu^+(A)+\nu^-(A)$. El signo en la cota es una suma de las variaciones, nunca su diferencia.
\end{proof}

\section{Ejercicios resueltos: orden e integración de medidas reales}
\begin{ejercicio}[Lectura 25, ejercicio 1]\label{ej:v2-l25-1}
Si $\nu=\nu_1-\nu_2$ con $\nu_1,\nu_2$ positivas y la diferencia definida, demostrar $\nu_1(A)\ge\nu^+(A)$ y $\nu_2(A)\ge\nu^-(A)$.
\end{ejercicio}
\begin{solucion}
En una parte positiva $P$, $\nu^+(A)=\nu(A\cap P)=\nu_1(A\cap P)-\nu_2(A\cap P)\le\nu_1(A\cap P)\le\nu_1(A)$. Si el primer miembro es infinito, la igualdad de la diferencia obliga a $\nu_1(A\cap P)=\infty$ y la desigualdad sigue siendo válida. En una parte negativa $N$ se aplica el mismo argumento a $-\nu$ y se obtiene $\nu^-(A)\le\nu_2(A)$.
\end{solucion}
\begin{ejercicio}[Lectura 25, ejercicio 2: rectificación del ínfimo]\label{ej:v2-l25-2}
Para medidas con signo finitas $\nu_1,\nu_2$, definir
\[
 \nu_1\vee\nu_2=\nu_1+(\nu_2-\nu_1)^+,
 \qquad \nu_1\wedge\nu_2=\nu_1-(\nu_1-\nu_2)^+.
\]
(i) Demostrar que la primera es el menor mayorante de ambas en el orden de medidas. (ii) Demostrar que la segunda es el mayor minorante. La lectura escribe $\nu_1-(\nu_2-\nu_1)^+$ para el ínfimo; examinar esa fórmula.
\end{ejercicio}
\begin{solucion}
El orden $\eta\ge\nu$ significa que $\eta-\nu$ es una medida positiva. Escribamos $d=\nu_2-\nu_1=d^+-d^-$. Entonces $\nu_1+d^+\ge\nu_1$ y $(\nu_1+d^+)-\nu_2=d^-\ge0$. Si $\xi\ge\nu_1,\nu_2$, tenemos $d=(\xi-\nu_1)-(\xi-\nu_2)$, una diferencia de positivas. El ejercicio \ref{ej:v2-l25-1} implica $d^+\le\xi-\nu_1$, luego $\nu_1\vee\nu_2\le\xi$.

Aplicamos este resultado a $-\nu_1,-\nu_2$ y cambiamos de signo: el mayor minorante es $-((-\nu_1)\vee(-\nu_2))=\nu_1-(\nu_1-\nu_2)^+$. La fórmula impresa es falsa: sobre un espacio de un punto, con $\nu_1=0$ y $\nu_2=\delta$, daría $-\delta$, mientras que el mayor minorante es cero.
\end{solucion}
\begin{ejercicio}[Lectura 25, ejercicio 3]\label{ej:v2-l25-3}
Si $|\nu(A)|\le\mu(A)$ para todo medible $A$, con $\mu$ positiva, demostrar $|\nu|(A)\le\mu(A)$.
\end{ejercicio}
\begin{solucion}
Para toda partición finita de $A$, $\sum_j|\nu(A_j)|\le\sum_j\mu(A_j)=\mu(A)$. Tomar el supremo sobre las particiones y usar la proposición \ref{v2:variacion-real} da la conclusión, incluso si $\mu(A)=\infty$.
\end{solucion}
\begin{ejercicio}[Lectura 25, ejercicio 4]\label{ej:v2-l25-4}
Demostrar que $\norm\nu=|\nu|(X)$ es una norma sobre el espacio de medidas con signo finitas.
\end{ejercicio}
\begin{solucion}
La descomposición de Hahn muestra que $\nu^+(X)=\nu(P)$ y $\nu^-(X)=-\nu(N)$ son finitos, por tanto la expresión está definida. Si $\norm\nu=0$, para todo $A$ tenemos $|\nu(A)|\le|\nu|(A)\le|\nu|(X)=0$, luego $\nu=0$. El supremo por particiones da $|a\nu|(X)=|a||\nu|(X)$. Para cada partición, $\sum_j|\nu(A_j)+\eta(A_j)|\le\sum_j|\nu(A_j)|+\sum_j|\eta(A_j)|\le\norm\nu+\norm\eta$; tomar supremo prueba la desigualdad triangular.
\end{solucion}
\begin{ejercicio}[Lectura 25, ejercicio 5]\label{ej:v2-l25-5}
Para medidas con signo finitas, demostrar $\norm{\nu_n-\nu}\to0$ si y solo si $\nu_n(A)\to\nu(A)$ uniformemente en $A\in\Acal$.
\end{ejercicio}
\begin{solucion}
Sea $\eta=\nu_n-\nu$. Por una parte $\sup_A|\eta(A)|\le\norm\eta$. Por otra, si $(P,N)$ es una partición de Hahn de $\eta$,
\[
 \norm\eta=|\eta(P)|+|\eta(N)|\le2\sup_A|\eta(A)|.
\]
Estas dos cotas prueban ambas implicaciones y muestran que las dos formas de medir la convergencia son equivalentes.
\end{solucion}
\begin{ejercicio}[Lectura 25, ejercicio 6]\label{ej:v2-l25-6}
Definir $\int f\,d\nu=\int f\,d\nu^+-\int f\,d\nu^-$ cuando $f$ es integrable para ambas medidas. Demostrar que tales funciones forman un espacio vectorial y decidir si $f\mapsto\int|f|\,d|\nu|$ es una seminorma.
\end{ejercicio}
\begin{solucion}
La condición equivale a $\int|f|\,d|\nu|=\int|f|\,d\nu^++\int|f|\,d\nu^-<\infty$, primero por aditividad para simples y después por convergencia monótona. Es, pues, exactamente $\mathcal L^1(|\nu|)$, un espacio vectorial. La expresión solicitada es su seminorma $L^1$: tiene homogeneidad y desigualdad triangular al integrar las correspondientes desigualdades puntuales. Su núcleo consta de funciones nulas $|\nu|$-casi en todas partes. La integral con signo es lineal y satisface $|\int f\,d\nu|\le\int|f|\,d|\nu|$; las dos integrales que se restan son finitas, por lo que no aparece una indeterminación.
\end{solucion}

\section{Variación de una medida compleja}\label{lectura:v2-26}
Si $\nu=\nu_1+i\nu_2$, las dos medidas reales son finitas y tienen sus variaciones de Jordan. Así
$\nu=\nu_1^+-\nu_1^-+i\nu_2^+-i\nu_2^-$. Esta expresión descompone las partes real e imaginaria en medidas positivas; su suma no es en general la variación total compleja.
\begin{definicion}
Para una medida compleja definimos
\[
 |\nu|(A)=\sup\left\{\sum_{j=1}^m|\nu(A_j)|:A=\mathop{\dot\bigcup}_{j=1}^mA_j,\ A_j\in\Acal\right\}.
\]
\end{definicion}
\begin{teorema}\label{v2:variacion-compleja}
La variación $|\nu|$ es una medida positiva finita. Para todo $A\in\Acal$,
\begin{gather*}
 |\nu|(A)\le\nu_1^+(A)+\nu_1^-(A)+\nu_2^+(A)+\nu_2^-(A),\\
 \sup_{B\subset A,\ B\in\Acal}|\nu(B)|\le|\nu|(A)
 \le4\sup_{B\subset A,\ B\in\Acal}|\nu(B)|,\\
  \nu_j^\pm(A)\le|\nu|(A).
\end{gather*}
En particular, $|\nu(A)|\le|\nu|(X)<\infty$.
\end{teorema}
\begin{proof}
Para cada partición, $|\nu(A_j)|\le|\nu_1(A_j)|+|\nu_2(A_j)|$. Sumando y usando las variaciones reales se obtiene la primera cota, que además prueba finitud, porque las cuatro medidas de Jordan son finitas. No basta, por sí solo, afirmar que cada valor de $\nu$ es un número complejo: hace falta esta cota uniforme para el supremo.

Para aditividad, sea $A=\dot\bigcup_n E_n$. Dada una partición finita $A_1,\ldots,A_m$ de $A$, la aditividad de $\nu$ y la desigualdad triangular dan
\[
 \sum_j|\nu(A_j)|\le\sum_n\sum_j|\nu(A_j\cap E_n)|\le\sum_n|\nu|(E_n).
\]
Por tanto $|\nu|(A)\le\sum_n|\nu|(E_n)$. Para la cota inversa, fijamos $N$ y elegimos en cada $E_n$, $n\le N$, una partición cuya suma se acerque a $|\nu|(E_n)$. Juntas, con el resto de $A$, forman una partición finita de $A$, de donde $|\nu|(A)\ge\sum_{n\le N}|\nu|(E_n)-\eps$. Hacemos $\eps\downarrow0$ y $N\to\infty$. Esto prueba aditividad numerable.

La partición $B,A\setminus B$ da la primera cota con el supremo. Para la segunda, sean $(P_j,N_j)$ particiones de Hahn de $\nu_j$. La primera cota del enunciado acota $|\nu|(A)$ por
\[
 |\nu_1(A\cap P_1)|+|\nu_1(A\cap N_1)|+
 |\nu_2(A\cap P_2)|+|\nu_2(A\cap N_2)|,
\]
y cada término es a lo sumo $\sup_{B\subset A}|\nu(B)|$. Finalmente $\nu_j^+(A)=\nu_j(A\cap P_j)\le|\nu(A\cap P_j)|\le|\nu|(A)$, y lo mismo vale para la parte negativa.
\end{proof}
\begin{teorema}\label{v2:medidas-banach}
El espacio $\mathcal M(X,\Acal;\F)$ de medidas reales finitas o complejas, con $\norm\nu=|\nu|(X)$, es un espacio de Banach.
\end{teorema}
\begin{proof}
Las propiedades de norma se demuestran como en el ejercicio \ref{ej:v2-l25-4}, utilizando particiones. Si $(\nu_n)$ es de Cauchy en variación, entonces $|\nu_n(A)-\nu_m(A)|\le\norm{\nu_n-\nu_m}$; existe $\nu(A)=\lim_n\nu_n(A)$ y la convergencia es uniforme en $A$. Pasar al límite en sumas finitas demuestra aditividad finita y $\nu(\varnothing)=0$.

Si $A_j\downarrow\varnothing$, fijamos $n$ con $\sup_A|\nu_n(A)-\nu(A)|<\eps/2$ y después $j$ con $|\nu_n(A_j)|<\eps/2$. Entonces $|\nu(A_j)|<\eps$. Así $\nu$ es continua en el vacío. Para conjuntos disjuntos $B_j$ con unión $B$, el resto $B\setminus\bigcup_{j\le m}B_j$ decrece al vacío; por aditividad finita, $\nu(B)-\sum_{j\le m}\nu(B_j)$ es su medida y tiende a cero. Por tanto $\nu$ es una medida.

Para una partición finita $A_1,\ldots,A_k$ y $n,m$ grandes, $\sum_j|\nu_n(A_j)-\nu_m(A_j)|\le\eps$. Hacer $m\to\infty$ y tomar después el supremo sobre todas las particiones da $\norm{\nu_n-\nu}\le\eps$. La misma cota, con un $n$ fijo, prueba que $\nu$ tiene variación finita. Este orden de los límites evita intercambiar indebidamente un supremo con una convergencia puntual.
\end{proof}
\begin{proposicion}[Integración compleja]\label{v2:integral-compleja}
Una función medible es integrable respecto a $|\nu|$ si y solo si lo es respecto a las cuatro medidas $\nu_j^\pm$. En ese caso definimos
\[
 \int f\,d\nu=\int f\,d\nu_1^+-\int f\,d\nu_1^-+
 i\int f\,d\nu_2^+-i\int f\,d\nu_2^-.
\]
La integral es lineal y $|\int f\,d\nu|\le\int|f|\,d|\nu|$.
\end{proposicion}
\begin{proof}
Las cotas entre medidas del teorema \ref{v2:variacion-compleja} se integran primero para funciones simples no negativas y después para funciones medibles no negativas, mediante aproximación creciente. Aplicadas a $|f|$ prueban la equivalencia de integrabilidad. Para una función simple $s=\sum_jc_j\ind{A_j}$ sobre conjuntos disjuntos, la definición da $\int s\,d\nu=\sum_jc_j\nu(A_j)$ y
$|\int s\,d\nu|\le\sum_j|c_j||\nu|(A_j)=\int|s|\,d|\nu|$.
Aproximamos $f$ en $L^1(|\nu|)$ por simples. La misma aproximación converge en las cuatro normas $L^1(\nu_j^\pm)$, por las cotas de medidas; por tanto sus cuatro integrales convergen a las de $f$. Pasando al límite se obtiene la desigualdad. La linealidad resulta de las cuatro integrales ordinarias. Si una función está definida solo casi en todas partes, se elige un representante medible y se extiende por cero sobre un conjunto excepcional medible; la desigualdad asegura independencia del representante.
\end{proof}

\section{Ejercicios resueltos: continuidad de medidas y norma de la integral}
\begin{ejercicio}[Lectura 26, ejercicio 1]\label{ej:v2-l26-1}
Sea $\nu:\Acal\to\C$ finitamente aditiva y nula en el vacío. Demostrar que es una medida si cumple cualquiera de estas condiciones: (i) continuidad sobre sucesiones crecientes; (ii) continuidad sobre sucesiones decrecientes con intersección vacía.
\end{ejercicio}
\begin{solucion}
Si $A_n$ son disjuntos, las uniones finitas $B_m=\bigcup_{n\le m}A_n$ crecen a $B=\bigcup_nA_n$. Bajo (i), $\nu(B)=\lim_m\nu(B_m)=\sum_n\nu(A_n)$. Bajo (ii), los restos $B\setminus B_m$ decrecen al vacío, y la identidad $\nu(B)-\sum_{n\le m}\nu(A_n)=\nu(B\setminus B_m)$ da la misma conclusión. No basta la convergencia de la serie sin una de estas condiciones de continuidad.
\end{solucion}
\begin{ejercicio}[Lectura 26, ejercicio 2: excepción de la medida cero]\label{ej:v2-l26-2}
Demostrar que $I_\nu(f)=\int f\,d\nu$ es un funcional continuo sobre $L^1(|\nu|)$, de norma uno cuando $\nu\ne0$. Determinar su norma si $\nu=0$.
\end{ejercicio}
\begin{solucion}
La proposición \ref{v2:integral-compleja} da $\norm{I_\nu}\le1$. Si $M=|\nu|(X)>0$, elegimos una partición $A_1,\ldots,A_m$ con $\sum_j|\nu(A_j)|>M-\eps$. Tomamos $c_j=\overline{\nu(A_j)}/|\nu(A_j)|$ cuando la medida es no nula y $c_j=1$ en otro caso. La simple $s=\sum_jc_j\ind{A_j}$ tiene $|s|=1$, luego $\norm s_{L^1(|\nu|)}=M$, y $I_\nu(s)=\sum_j|\nu(A_j)|>M-\eps$. Aplicar el funcional a $s/M$ y hacer $\eps\downarrow0$ da norma al menos uno. Para $\nu=0$, $L^1(|\nu|)$ es el espacio cero y el funcional tiene norma cero. La excepción no puede omitirse.
\end{solucion}
\section{La variación como medida de control}
Hahn separa signos; Jordan transforma esa separación en medidas positivas canónicas; la variación total controla todas las particiones a la vez. Con ella, la integración compleja se reduce a integración positiva y el espacio de medidas adquiere una norma completa. La aditividad numerable sigue siendo esencial: en el capítulo sobre el dual de $L^\infty$ veremos qué cambia al conservar únicamente aditividad finita.

\chapter{Radon--Nikodym y descomposición de Lebesgue}\label{v2:cap13}\label{cap:v2-13}
\section{Cuándo una medida tiene densidad}\label{lectura:v2-27}
Si $f\in L^1(\mu)$, la integral indefinida $\nu_f(A)=\int_A f\,d\mu$ es una medida finita. El problema inverso consiste en recuperar $f$ a partir de $\nu_f$. Sobre espacios generales, anularse en conjuntos nulos no garantiza que una medida finita provenga de una densidad integrable: puede haber una parte que no se detecta en ningún conjunto de medida $\mu$ finita. Las lecturas 27 y 30 de \cite{valencia2} estudian precisamente esta dificultad.
\begin{definicion}
Para una función finitamente aditiva $\nu:\Acal\to\F$, distinguimos:
\begin{enumerate}
\item \emph{Continuidad absoluta en conjuntos nulos}: $\mu(A)=0$ implica $\nu(A)=0$; escribimos $\nu\ll\mu$.
\item \emph{Continuidad absoluta uniforme}: para cada $\eps>0$ existe $\delta>0$ tal que $\mu(A)<\delta$ implica $|\nu(A)|<\eps$.
\item \emph{Continuidad verdadera}: para cada $\eps>0$ existen $F\in\Acal$, $\mu(F)<\infty$, y $\delta>0$ tales que $\mu(B\cap F)<\delta$ implica $|\nu(B)|<\eps$ para todo $B\in\Acal$.
\end{enumerate}
La lectura denomina continuidad absoluta a (ii). Para medidas finitas numerablemente aditivas, (i) y (ii) coinciden; para funciones solamente finitamente aditivas pueden diferir.\index{continuidad absoluta de medidas}\index{continuidad verdadera}
\end{definicion}
\begin{lema}\label{v2:absoluta-equivalencia}
Si $\nu$ es una medida finita real o compleja, entonces $\nu\ll\mu$ equivale a continuidad absoluta uniforme. Además $\nu\ll\mu$ equivale a $|\nu|\ll\mu$.
\end{lema}
\begin{proof}
Si $\mu(A)=0$, todo subconjunto medible de $A$ es nulo; por tanto, si $\nu\ll\mu$, todas las sumas de particiones que definen $|\nu|(A)$ son cero. La implicación inversa se sigue de $|\nu(A)|\le|\nu|(A)$.

La continuidad uniforme implica anulación sobre nulos haciendo $\eps\downarrow0$. Para la otra dirección, si falla la condición uniforme, existen $\eps>0$ y conjuntos $A_n$ con $\mu(A_n)<2^{-n}$ pero $|\nu(A_n)|\ge\eps$. Sean $B_m=\bigcup_{n\ge m}A_n$ y $N=\bigcap_m B_m$. Entonces $\mu(B_m)\le2^{1-m}$, luego $\mu(N)=0$ y $|\nu|(N)=0$. Por continuidad decreciente de la medida finita $|\nu|$, $|\nu|(B_m)\to0$, contradiciendo $|\nu|(B_m)\ge|\nu(A_m)|\ge\eps$.
\end{proof}
\begin{proposicion}\label{v2:verdadera-caracterizacion}
Una función finitamente aditiva real o compleja $\nu$ es verdaderamente continua si y solo si:
\begin{enumerate}
\item es una medida finita y $\nu\ll\mu$;
\item para cada $A$ con $\nu(A)\ne0$, existe $F$ de medida $\mu$ finita con $\nu(A\cap F)\ne0$.
\end{enumerate}
Equivalentemente, es una medida finita absolutamente continua y existe un conjunto $S$ de medida $\mu$ $\sigma$-finita tal que $|\nu|(X\setminus S)=0$.
\end{proposicion}
\begin{proof}
Supongamos continuidad verdadera. Implica continuidad absoluta uniforme, pues $\mu(B\cap F)\le\mu(B)$. Si $\nu(A)\ne0$, elegimos el control verdadero con $\eps<|\nu(A)|$. Entonces $|\nu(A\setminus F)|<\eps$, de modo que $\nu(A\cap F)=\nu(A)-\nu(A\setminus F)\ne0$.

Para aditividad numerable, sean $A_n$ disjuntos y $A=\bigcup_nA_n$. Dado $\eps>0$, elegimos $F,\delta$ como en la definición. Puesto que $\mu(F)<\infty$, las colas satisfacen $\mu(F\cap\bigcup_{n>m}A_n)\to0$. Por aditividad finita,
\[
 \left|\nu(A)-\sum_{n\le m}\nu(A_n)\right|
 =\left|\nu\left(\bigcup_{n>m}A_n\right)\right|<\eps
\]
para $m$ grande. Así $\nu$ es una medida. Sus valores son finitos por el codominio $\F$, y por el capítulo anterior tiene variación finita.

Recíprocamente, supongamos (i)--(ii). Sea $\alpha=\sup\{|\nu|(F):\mu(F)<\infty\}$. Elegimos conjuntos crecientes $F_n$ de medida finita con $|\nu|(F_n)\to\alpha$ y ponemos $S=\bigcup_nF_n$. Si $B\subset X\setminus S$ y $\nu(B)\ne0$, (ii) proporciona un $F$ de medida finita con $\nu(B\cap F)\ne0$, luego $|\nu|(F\setminus S)>0$. Para $n$ grande,
$|\nu|(F_n\cup(F\setminus S))=|\nu|(F_n)+|\nu|(F\setminus S)>\alpha$, contradicción. Por tanto $\nu$ se anula en todos los subconjuntos medibles de $X\setminus S$, es decir, $|\nu|(X\setminus S)=0$.

Dado $\eps>0$, elegimos $n$ con $|\nu|(S\setminus F_n)<\eps/2$. El lema \ref{v2:absoluta-equivalencia} da $\delta>0$ para que $\mu(E)<\delta$ implique $|\nu(E)|<\eps/2$. Si $\mu(B\cap F_n)<\delta$, entonces
\[
 |\nu(B)|\le|\nu(B\cap F_n)|+|\nu|(S\setminus F_n)+|\nu|(X\setminus S)<\eps.
\]
Esto es continuidad verdadera. Finalmente, si existe tal $S=\bigcup_nF_n$ con $F_n$ crecientes de medida finita, la misma estimación prueba continuidad verdadera; además, $\nu(A)=\lim_n\nu(A\cap F_n)$ demuestra (ii).
\end{proof}
\begin{proposicion}\label{v2:verdaderas-subespacio}
Las medidas finitas verdaderamente continuas forman un subespacio cerrado $\mathcal M_\mu\subset\mathcal M(X,\Acal;\F)$. Si $\nu\in\mathcal M_\mu$, también pertenecen a él $|\nu|$ y las medidas de Jordan de sus partes real e imaginaria.
\end{proposicion}
\begin{proof}
Para una combinación lineal se toma la unión de los dos conjuntos de control finito y el mínimo de sus umbrales, ajustando $\eps$ a los coeficientes. Para variación, si $\mu(B\cap F)<\delta$, la misma desigualdad se cumple para todo subconjunto $C\subset B$. La cota $|\nu|(B)\le4\sup_{C\subset B}|\nu(C)|$ prueba continuidad verdadera de $|\nu|$, empezando con $\eps/4$. Las partes de Jordan están dominadas por $|\nu|$ y tienen la misma propiedad. Si $\norm{\nu_j-\nu}\to0$, elegimos $j$ con esa norma menor que $\eps/2$ y aplicamos el control verdadero de $\nu_j$ con $\eps/2$; entonces $|\nu(B)|\le\norm{\nu-\nu_j}+|\nu_j(B)|<\eps$. Esto demuestra cerradura.
\end{proof}
\begin{proposicion}\label{v2:densidad-variacion}
Para $f\in L^1(\mu)$, la medida $\nu_f(A)=\int_A f\,d\mu$ pertenece a $\mathcal M_\mu$ y
\[
 |\nu_f|(A)=\int_A|f|\,d\mu,\qquad \norm{\nu_f}=\norm f_1.
\]
La aplicación $f\mapsto\nu_f$ es lineal e inyectiva sobre clases casi en todas partes.
\end{proposicion}
\begin{proof}
La integral de $|f|$ es uniformemente pequeña en conjuntos de medida pequeña y se concentra en conjuntos de medida finita, como se probó en Vitali. Elegimos $F$ con $\int_{X\setminus F}|f|<\eps/2$ y $\delta$ para el control en $B\cap F$. Entonces $|\nu_f(B)|\le\int_{B\cap F}|f|+\int_{X\setminus F}|f|<\eps$.

Para toda partición de $A$, $\sum_j|\int_{A_j}f|\le\int_A|f|$, luego $|\nu_f|(A)\le\int_A|f|$. Si $s$ es simple integrable, su partición por valores muestra igualdad, pues en cada pieza la fase es constante. Elegimos simples $s_n\to f$ en $L^1$. La desigualdad ya probada da $\norm{\nu_f-\nu_{s_n}}\le\norm{f-s_n}_1$. Por la desigualdad triangular, $\norm{\nu_f}\ge\norm{s_n}_1-\norm{f-s_n}_1\to\norm f_1$. Así hay igualdad en $X$, y aplicando el mismo argumento a $f\ind A$ se obtiene en todo $A$. La igualdad de normas prueba inyectividad sobre $L^1$.
\end{proof}

\section{Construcción de la derivada}
\begin{teorema}[Radon--Nikodym en la forma general de las lecturas]\label{v2:radon-nikodym}\index{Radon--Nikodym, teorema de}
Una función finitamente aditiva $\nu:\Acal\to\F$ es verdaderamente continua con respecto a $\mu$ si y solo si existe $f\in L^1(\mu)$ tal que
\[
 \nu(A)=\int_A f\,d\mu\qquad(A\in\Acal).
\]
La clase de $f$ es única; se denota $d\nu/d\mu$.\index{teorema de Radon--Nikodym}
\end{teorema}
\begin{proof}
Una dirección es la proposición anterior. Para la otra, comencemos con $\nu\ge0$. La caracterización de continuidad verdadera garantiza que es una medida positiva finita. Si $\nu=0$, tomamos $f=0$.

\emph{Un primer fragmento de densidad.} Si $\nu\ne0$, existe un conjunto $A$ de medida $\mu$ finita con $\nu(A)>0$. Como $\nu\ll\mu$, también $\mu(A)>0$. Pongamos $c=\nu(A)/(2\mu(A))$ y $\eta=\nu-c\mu_A$, donde $\mu_A(B)=\mu(A\cap B)$. Es una medida con signo finita y $\eta(A)=\nu(A)/2>0$. Sea $(P,N)$ una partición de Hahn para $\eta$. Entonces $\eta(A\cap P)>0$, pues $\eta(A\cap N)\le0$. En particular $\nu(A\cap P)>0$, y por continuidad absoluta $\mu(A\cap P)>0$. Para $\psi=c\ind{A\cap P}$ tenemos $\int\psi>0$ y, para todo $B$,
\[
 \int_B\psi=c\mu(A\cap B\cap P)\le\nu(B\cap P)\le\nu(B).
\]

\emph{Aproximar una densidad maximal.} Sea $\mathcal S_\nu$ el conjunto de simples no negativas $s$ con $\int_Bs\,d\mu\le\nu(B)$ para todo $B$. Contiene cero. Si $s,t$ pertenecen a él, también $\max\{s,t\}$: dividimos cada $B$ en $B\cap\{s\ge t\}$ y su complemento y aplicamos las dos desigualdades. Sea $\alpha=\sup_{s\in\mathcal S_\nu}\int s\le\nu(X)<\infty$. Elijamos $s_n$ aproximando el supremo y sustituyámoslas por sus máximos sucesivos. Así $s_n\uparrow f$ y convergencia monótona da $\int f=\alpha$ y $\int_B f\le\nu(B)$. El conjunto donde el límite es infinito es medible y nulo; definimos $f=0$ allí para obtener un representante finito sin modificar integrales.

La medida residual $\rho=\nu-\nu_f$ es positiva y verdaderamente continua, por la proposición \ref{v2:verdaderas-subespacio}. Si no fuera cero, el primer paso daría una simple $\psi\ge0$ con $\int\psi>0$ y $\int_B\psi\le\rho(B)$. Entonces $s_n+\psi\in\mathcal S_\nu$ para todo $n$. Para $n$ suficientemente grande, $\int s_n+\int\psi>\alpha$, contradicción. Por tanto $\rho=0$ y $\nu=\nu_f$.

Para una medida real con signo, aplicamos el caso positivo a $\nu^+,\nu^-$, que son verdaderamente continuas, y restamos sus densidades. Para una compleja aplicamos el procedimiento a las cuatro medidas de Jordan y combinamos las densidades con coeficientes $1,-1,i,-i$. La unicidad y la igualdad de normas se siguen de la proposición \ref{v2:densidad-variacion}.
\end{proof}
\begin{corolario}\label{v2:rn-sigmafinita}
Si $\mu$ es $\sigma$-finita y $\nu$ es una medida finita real o compleja, entonces $\nu\ll\mu$ si y solo si $\nu$ tiene una densidad en $L^1(\mu)$.
\end{corolario}
\begin{proof}
En la caracterización de continuidad verdadera podemos tomar $S=X$, que es $\sigma$-finito. Aplicamos el teorema. La implicación desde una densidad a la anulación sobre nulos es inmediata de la integral.
\end{proof}

\section{Ejercicios propuestos: hipótesis de continuidad absoluta}
\begin{ejercicio}[Lectura 27, ejercicio 1: aditividad numerable necesaria]\label{ej:v2-l27-1}
Examinar la afirmación original: para una función solamente finitamente aditiva, continuidad absoluta uniforme equivale a anularse en conjuntos $\mu$-nulos. Dar la versión correcta.
\end{ejercicio}

\begin{ejercicio}[Lectura 27, ejercicio 2: aclaración de notación]\label{ej:v2-l27-2}
Si $\mu$ es $\sigma$-finita, caracterizar las funciones finitamente aditivas verdaderamente continuas. La lectura pide que esto equivalga a pertenecer a $\mathcal M_\mu$, espacio que ella misma había definido mediante continuidad verdadera.
\end{ejercicio}

\begin{ejercicio}[Lectura 27, ejercicio 3]\label{ej:v2-l27-3}
Demostrar que $\mathcal M_\mu$ es un subespacio cerrado de $\mathcal M(X,\Acal;\R)$ para cualquier medida $\mu$.
\end{ejercicio}


\section{Singularidad y las tres partes de una medida}\label{lectura:v2-30}
\begin{definicion}
Dos medidas positivas, con signo o complejas son \emph{mutuamente singulares}, $\eta\perp\nu$, si existe $N\in\Acal$ con $|\eta|(N)=0$ y $|\nu|(X\setminus N)=0$; para una medida positiva su variación es ella misma. La relación es simétrica, reemplazando $N$ por su complemento. Fijada $\mu\ge0$, definimos
\begin{align*}
 \mathcal M_\mu^*&=\{\nu\in\mathcal M:\nu\ll\mu,\ \nu(F)=0\text{ si }\mu(F)<\infty\},\\
 \mathcal M_\mu^\perp&=\{\nu\in\mathcal M:\nu\perp\mu\}.
\end{align*}
\end{definicion}
\begin{proposicion}\label{v2:singulares-cerrados}
Estos dos conjuntos son subespacios cerrados. La singularidad de dos medidas equivale a la de sus variaciones. En el caso real equivale a la singularidad de cada una de las dos partes de Jordan de la primera con cada parte de la segunda; en el caso complejo se exige la condición para todos los pares de sus cuatro componentes positivas.
\end{proposicion}
\begin{proof}
Para $\mathcal M_\mu^*$ las condiciones son identidades sobre conjuntos determinados y pasan al límite, porque $|\nu_n(A)-\nu(A)|\le\norm{\nu_n-\nu}$. También se preservan por combinaciones lineales. Si $\nu_n\perp\mu$, elegimos conjuntos nulos $N_n$ que concentran sus variaciones. La unión $N=\bigcup_nN_n$ es nula y $|\nu|(X\setminus N)\le\norm{\nu-\nu_n}\to0$, pues $\nu_n$ se anula allí con toda su variación. La misma unión, para dos medidas, prueba estabilidad lineal de las singulares.

La equivalencia con las variaciones es parte de la definición y de la caracterización de los conjuntos donde una medida y todas sus restricciones se anulan. Si las variaciones son singulares, cada componente de Jordan, dominada por su variación, también lo es frente a todas las componentes de la otra medida. Recíprocamente, supongamos que unas familias finitas positivas $\alpha_i,\beta_j$ son singulares por pares. Para cada $i,j$ elegimos $N_{ij}$ con $\alpha_i(N_{ij})=0$ y $\beta_j(X\setminus N_{ij})=0$. Ponemos $N_j=\bigcap_iN_{ij}$ y $N=\bigcup_jN_j$. Entonces todas las $\alpha_i$ se anulan en $N$ y todas las $\beta_j$ se concentran en $N$. Las cotas de las variaciones por las sumas de componentes dan la singularidad de las medidas originales. Esto explicita que se requieren \emph{todos} los pares de signos y componentes.
\end{proof}
\begin{teorema}[Descomposición de Lebesgue general]\label{v2:lebesgue-descomposicion}\index{descomposicion de Lebesgue@descomposición de Lebesgue}
Para cualquier medida positiva $\mu$,
\[
 \mathcal M(X,\Acal;\F)=\mathcal M_\mu\oplus\mathcal M_\mu^*\oplus\mathcal M_\mu^\perp.
\]
Es decir, toda medida finita se descompone de manera única como $\nu=\nu_c+\nu_e+\nu_s$, donde $\nu_c$ tiene densidad $L^1(\mu)$, $\nu_e$ es absolutamente continua pero se anula en todo conjunto de medida $\mu$ finita, y $\nu_s$ es singular.
\end{teorema}
\begin{proof}
\emph{Construcción para $\nu\ge0$ finita.} Consideremos las funciones de conjuntos que aparecen en la lectura:
\[
 \nu_s(A)=\sup_{\mu(F)=0}\nu(A\cap F),\qquad
 \nu_t(A)=\sup_{\mu(F)<\infty}\nu(A\cap F).
\]
Demostraremos a la vez que son medidas y que se obtienen por restricción a conjuntos fijos. Para la primera, elegimos nulos $N_j$ con $\nu(N_j)$ acercándose al supremo $\nu_s(X)$ y ponemos $N=\bigcup_jN_j$. Entonces $\mu(N)=0$ y $\nu(N)=\nu_s(X)$. Si $F$ es nulo, también lo es $F\cup N$, por lo que $\nu(F\setminus N)=0$. Así $\nu(A\cap F)\le\nu(A\cap N)$ para todo $F$ nulo, y tomando $F=N$ obtenemos $\nu_s(A)=\nu(A\cap N)$.

Para la segunda, elegimos conjuntos crecientes $F_j$ de medida finita con $\nu(F_j)\to\nu_t(X)$ y ponemos $S=\bigcup_jF_j$. Si $F$ tiene medida finita, también la tiene $F\cup F_j$, de modo que $\nu(F\setminus S)=0$ al pasar al límite. Por tanto $\nu_t(A)\le\nu(A\cap S)$. La desigualdad inversa se sigue de $\nu(A\cap S)=\lim_j\nu(A\cap F_j)\le\nu_t(A)$. Podemos agregar $N$ a cada $F_j$ sin cambiar su finitud de medida; así suponemos $N\subset S$.

En particular, $\nu_s,\nu_t$ son medidas positivas, pues son restricciones de $\nu$. Definimos
\[
 \nu_c=\nu\big|_{S\setminus N}=\nu_t-\nu_s,\qquad
 \nu_e=\nu\big|_{X\setminus S}=\nu-\nu_t.
\]
La medida $\nu_s$ es singular, concentrada en $N$. Si $F$ es $\mu$-nulo, la maximalidad de $N$ da $\nu(F\setminus N)=0$, luego $\nu_c\ll\mu$. Está concentrada en el conjunto $\sigma$-finito $S$, así que es verdaderamente continua por la proposición \ref{v2:verdadera-caracterizacion}. Por construcción $\nu_e(F)=0$ para todo $F$ de medida $\mu$ finita; en particular se anula sobre nulos y pertenece a $\mathcal M_\mu^*$. La suma de las tres restricciones es $\nu$.

\emph{Medidas reales y complejas.} Descomponemos las cuatro medidas positivas de Jordan por el procedimiento anterior y reunimos sus partes con los coeficientes $1,-1,i,-i$. La estabilidad lineal de los tres subespacios prueba existencia en general.

\emph{Unicidad.} Si una medida es a la vez absolutamente continua y singular, su variación está concentrada en un conjunto $\mu$-nulo donde también debe anularse; por tanto es cero. De dos descomposiciones se sigue primero $\nu_s=\nu_s'$, pues su diferencia es singular y coincide con una medida absolutamente continua. Queda $\nu_c-\nu_c'=\nu_e'-\nu_e$. El lado izquierdo se concentra en una unión numerable $S=\bigcup_jF_j$ de conjuntos de medida finita; el derecho se anula en todos los subconjuntos medibles de cada $F_j$, y por aditividad numerable se anula en todos los subconjuntos de $S$. Fuera de $S$ el lado izquierdo se anula. Así ambos lados son cero, y coinciden las otras dos partes. Esto demuestra que la suma es directa.
\end{proof}
\begin{observacion}[El caso habitual y la parte invisible]
Si $\mu$ es $\sigma$-finita, $\mathcal M_\mu^*=\{0\}$: una medida de ese espacio se anula sobre cada pieza de una cobertura numerable de medida finita, y por aditividad se anula en todo conjunto. Se recupera la descomposición usual en una parte con densidad y una parte singular. Para ver que el término adicional puede existir, tome $X$ de un punto, $\mu(X)=\infty$ y $\nu(X)=1$. El único conjunto $\mu$-nulo o de medida finita es el vacío. Así $\nu\in\mathcal M_\mu^*$, pero ninguna función de $L^1(\mu)$ puede producirla: toda función integrable debe anularse en el punto.
\end{observacion}
\section{Densidad, singularidad y detección local}
La derivada de Radon--Nikodym representa exactamente la parte que es absolutamente continua y se concentra en regiones de medida finita numerablemente reunidas. La descomposición general identifica por separado la masa singular y la masa que ningún conjunto de medida finita puede detectar. Esta distinción será decisiva para representar funcionales sobre $L^1$ en espacios localizables.

\chapter{Representación del dual de \texorpdfstring{$L^1$}{L1} y del dual de \texorpdfstring{$L^\infty$}{L infinito}}\label{v2:cap14}\label{cap:v2-14}
\section{Pegar representaciones locales}\label{lectura:v2-28}
En un conjunto de medida finita, un funcional sobre $L^1$ determina una medida mediante sus valores en indicadores. Radon--Nikodym transforma esa medida en una función. Para trabajar en un espacio que no es $\sigma$-finito, debemos reunir una familia posiblemente no numerable de representaciones locales sin tomar una unión no numerable de excepciones nulas. La localizabilidad proporciona precisamente la herramienta adecuada. Este capítulo desarrolla las lecturas 28 y 29 de \cite{valencia2}.
\begin{lema}[Pegado medible]\label{v2:pegado-medible}
Supóngase que $\mu$ es localizable. Sean $f_\alpha:D_\alpha\to\R$ funciones medibles sobre dominios medibles, con $f_\alpha=f_\beta$ casi en todas partes en $D_\alpha\cap D_\beta$. Entonces existe una función medible $f:X\to\R$ cuya restricción a cada $D_\alpha$ coincide casi en todas partes con $f_\alpha$. La afirmación vale también para funciones complejas.
\end{lema}
\begin{proof}
Para cada racional $r$, sea $E_{\alpha,r}=\{x\in D_\alpha:f_\alpha(x)>r\}$ y elijamos un supremo esencial $E_r$ de esta familia. Fijemos $\beta$. La compatibilidad implica
\[
 \mu\bigl(E_{\alpha,r}\setminus(X\setminus(D_\beta\setminus E_{\beta,r}))\bigr)=0
\]
para todo $\alpha$: la diferencia consiste en puntos comunes a ambos dominios donde $f_\alpha>r$ y $f_\beta\le r$. La minimalidad del supremo esencial da $\mu((E_r\cap D_\beta)\setminus E_{\beta,r})=0$. La propiedad de mayorante da $\mu(E_{\beta,r}\setminus E_r)=0$. Así ambas secciones coinciden salvo un conjunto nulo.

Definimos $f_*(x)=\sup\{r\in\Q:x\in E_r\}$ en la recta extendida, con supremo del vacío igual a $-\infty$. Es medible, pues $\{f_*>a\}=\bigcup_{r\in\Q,\ r>a}E_r$. Fuera de la unión numerable de las diferencias simétricas anteriores, para cada $x\in D_\beta$ se cumple
\[
 x\in E_r\quad\Longleftrightarrow\quad f_\beta(x)>r\qquad(r\in\Q),
\]
y por densidad de los racionales $f_*(x)=f_\beta(x)$. Sustituimos $f_*$ por cero en el conjunto medible donde vale $\pm\infty$; obtenemos la función finita requerida. En cada dominio el cambio afecta solo a un conjunto nulo, aunque no afirmamos que la unión de todas esas excepciones sea nula. Para funciones complejas pegamos por separado partes real e imaginaria. Deben usarse superniveles $f_\alpha>r$ con un supremo: subniveles $f_\alpha\le r$ producirían un supremo infinito y no reconstruirían la función.
\end{proof}
\begin{teorema}[Dual de $L^1$ en un espacio general]\label{v2:dual-l1-general}
La aplicación $\Phi:L^\infty(\mu)\to(L^1(\mu))^*$, $\Phi(g)(f)=\int fg\,d\mu$, es una isometría si y solo si $\mu$ es semifinita. Es una isometría lineal sobreyectiva si y solo si $\mu$ es localizable.
\end{teorema}
\begin{proof}
La primera equivalencia ya está demostrada en la proposición \ref{v2:l1-isometria}. Probamos las dos direcciones de la segunda.

\emph{Localizabilidad implica representación.} Sea $\varphi\in(L^1)^*$ y $M=\norm\varphi$. Para cada $A$ de medida finita, definimos $\nu_A(B)=\varphi(\ind{A\cap B})$. Es finitamente aditiva, y $|\nu_A(B)|\le M\mu(A\cap B)$ demuestra directamente continuidad verdadera. Radon--Nikodym da una $g_A\in L^1(\mu)$ con $\nu_A(B)=\int_Bg_A\,d\mu$. Se anula casi en todas partes fuera de $A$, por unicidad de las densidades.

Afirmamos $|g_A|\le M$ casi en todas partes. Para toda función simple $h$ soportada en $A$, la representación da $\varphi(h)=\int hg_A$. Se extiende a funciones acotadas soportadas en $A$ por aproximación uniforme: esta converge en $L^1(\mu)$, porque $\mu(A)<\infty$, y también hace converger la integral contra $g_A\in L^1$. Si $E=\{x\in A:|g_A(x)|>M+\eps\}$ tuviera medida positiva, la función acotada $h=\ind E\,\overline{g_A}/|g_A|$ daría
\[
 (M+\eps)\mu(E)<\int_E|g_A|=\varphi(h)\le M\norm h_1=M\mu(E),
\]
contradicción. Por tanto la cota está probada. Gracias a ella, la representación se extiende a todo $L^1(A)$ por densidad de las simples integrables.

En $A\cap B$ las dos densidades representan al mismo funcional sobre indicadores. Su diferencia integra cero sobre cada subconjunto de la intersección; por unicidad de la densidad, $g_A=g_B$ casi en todas partes allí. Aplicamos el lema de pegado a $g_A|_A$ y obtenemos una función medible $g$ en $X$. En cada conjunto de medida finita se cumple $|g|\le M$ casi en todas partes. Si $\{|g|>M\}$ tuviera medida positiva, la semifinitud produciría un subconjunto de medida positiva finita, contradicción. Así $g\in L^\infty$ y $\norm g_\infty\le M$. Para todo indicador integrable, $\varphi(\ind A)=\int_Ag$; por linealidad y densidad, $\varphi(f)=\int fg$ para todo $f\in L^1$. La isometría previamente demostrada da igualdad de normas.

\emph{Una isometría sobreyectiva implica localizabilidad.} Ya sabemos que $\mu$ es semifinita. Sea $\mathcal E\subset\Acal$ una familia y sea $\mathcal D$ la familia de uniones finitas de sus miembros, incluyendo el vacío. Para $f\ge0$ de $L^1$, ponemos
\[
 P(f)=\sup_{D\in\mathcal D}\int_D f\,d\mu.
\]
Tenemos $0\le P(f)\le\norm f_1$. Es aditiva en el cono positivo. La cota $P(f+h)\le P(f)+P(h)$ es inmediata; para la inversa elegimos $D_1,D_2$ que aproximen los dos supremos y usamos $D_1\cup D_2\in\mathcal D$, obteniendo un error arbitrariamente pequeño. También es homogénea para escalares no negativos. Definimos $P(f)=P(f^+)-P(f^-)$ para funciones reales; la aditividad en el cono garantiza que esto es lineal, pues cualquier descomposición $f=a-b$ con $a,b\ge0$ da el mismo valor: de $a-b=c-d$ se sigue $a+d=c+b$. Extendemos linealmente al caso complejo. La positividad da $|P(f)|\le P(|f|)\le\norm f_1$: en el caso complejo multiplicamos $P(f)$ por una fase que lo vuelva real positivo y usamos $\operatorname{Re}(e^{i\theta}f)\le|f|$.

Por sobreyectividad existe $g\in L^\infty$ que representa a $P$. La positividad sobre indicadores de medida finita y la semifinitud implican que $g$ es real y no negativa casi en todas partes: una parte imaginaria no nula o una parte real negativa se detectaría en un conjunto de medida positiva finita y contradiría los valores reales no negativos de $P$. Elegimos tal representante y ponemos $H=\{g>0\}$.

Si $E\in\mathcal E$ y $F\subset E\setminus H$ tiene medida finita, entonces $P(\ind F)=\mu(F)$, porque $E\in\mathcal D$, mientras que $\int_Fg=0$. Así $\mu(F)=0$. Por semifinitud $\mu(E\setminus H)=0$. Si $G$ contiene esencialmente a todos los miembros de $\mathcal E$, también contiene a cada $D\in\mathcal D$. Para $F\subset H\setminus G$ de medida finita tenemos $P(\ind F)=\sup_D\mu(F\cap D)=0$, y por representación $\int_Fg=0$. Como $g>0$ en $F$, se sigue $\mu(F)=0$; semifinitud implica $\mu(H\setminus G)=0$. Por tanto $H$ es el supremo esencial requerido.
\end{proof}

\section{Cargas de variación acotada}\label{lectura:v2-29}
\begin{definicion}
Una \emph{carga} es una función finitamente aditiva $\nu:\Acal\to\F$ con $\nu(\varnothing)=0$. Definimos su variación por
\[
 |\nu|(A)=\sup\left\{\sum_{j=1}^m|\nu(A_j)|:A=\mathop{\dot\bigcup}_{j=1}^mA_j,\ A_j\in\Acal\right\}.
\]
Sea $\operatorname{ba}(\Acal)$ el espacio de cargas de variación total finita y sea $\operatorname{ba}_\mu(\Acal)$ el subespacio de las que se anulan sobre conjuntos $\mu$-nulos. Esta última condición es continuidad absoluta en nulos, no continuidad absoluta uniforme. La lectura usa las notaciones $T$ y $T_\mu$ para estos espacios.\index{carga}\index{medida!finitamente aditiva}
\end{definicion}
\begin{proposicion}\label{v2:ba-banach}
La variación de una carga es una función positiva finitamente aditiva, posiblemente infinita. Para una carga de rango acotado,
\[
 \sup_{B\subset A}|\nu(B)|\le|\nu|(A)\le4\sup_{B\subset A}|\nu(B)|.
\]
Por tanto rango acotado equivale a variación finita. Con la norma $\norm\nu=|\nu|(X)$, $\operatorname{ba}(\Acal)$ es Banach y $\operatorname{ba}_\mu(\Acal)$ es un subespacio cerrado.
\end{proposicion}
\begin{proof}
Si $A,B$ son disjuntos, toda partición de $A\cup B$ se refina intersectando sus piezas con $A,B$, y la desigualdad triangular da $|\nu|(A\cup B)\le|\nu|(A)+|\nu|(B)$. Uniendo particiones que aproximan los dos supremos se obtiene la desigualdad contraria. Esto prueba aditividad finita; no afirma aditividad numerable.

La cota inferior usa la partición $B,A\setminus B$. Para la superior, dada una partición $A_j$ de $A$, agrupamos las piezas donde $\operatorname{Re}\nu(A_j)$ es no negativa y donde es negativa. Por aditividad finita, la suma de los valores absolutos de estas partes reales es a lo sumo $2\sup_{B\subset A}|\nu(B)|$. Agrupando del mismo modo las partes imaginarias y usando $|z|\le|\operatorname{Re}z|+|\operatorname{Im}z|$, obtenemos la constante cuatro. En el caso real basta dos.

Las propiedades de norma proceden de las particiones, como para medidas. Si $(\nu_n)$ es de Cauchy en esta norma, los valores convergen uniformemente sobre $\Acal$ a una función $\nu$ finitamente aditiva. Para cada partición finita, pasamos al límite en $\sum_j|\nu_n(A_j)-\nu_m(A_j)|\le\eps$ y tomamos después supremo; resulta $\norm{\nu_n-\nu}\le\eps$. Esto prueba variación finita y completitud sin exigir continuidad sobre sucesiones de conjuntos. Si todas las $\nu_n$ se anulan sobre un conjunto nulo $A$, el límite también lo hace. La condición es asimismo estable por combinaciones lineales, de donde la cerradura del subespacio.
\end{proof}

\section{Integración contra una carga}
Para una función simple $s=\sum_jc_j\ind{A_j}$ definimos $I_\nu(s)=\sum_jc_j\nu(A_j)$. Una refinación común de las particiones muestra que no depende de la representación y que
\[
 |I_\nu(s)-I_\nu(t)|\le\norm\nu\,\sup_X|s-t|,\qquad
 |I_\nu(s)|\le I_{|\nu|}(|s|).
\]
Toda función medible acotada $f$ se aproxima uniformemente por simples, mediante particiones finitas de su recorrido. Por la primera cota, sus integrales forman una sucesión de Cauchy; definimos $I_\nu(f)$ como su límite. Si se desean aproximaciones con módulos crecientes, se aproxima $|f|$ por los escalones diádicos inferiores $r_n=2^{-n}\lfloor2^n|f|\rfloor$ y su argumento por ángulos de una cuadrícula de paso $2\pi2^{-n}$. La simple $s_n=r_ne^{i\theta_n}$ tiene $|s_n|=r_n\uparrow|f|$ y error uniforme a lo sumo $2^{-n}+2\pi\norm f_{\mathrm{sup}}2^{-n}$. En el caso real se conserva el signo y solo se discretiza el módulo.

La cota anterior prueba independencia de la aproximación y continuidad para la norma uniforme. La linealidad pasa al límite desde las simples. También pasa al límite la desigualdad
\[
 \left|\int f\,d\nu\right|\le\int|f|\,d|\nu|,
\]
pues $s_n\to f$ uniformemente implica $|s_n|\to|f|$ uniformemente y $|\nu|$ tiene masa finita. Si $\nu\in\operatorname{ba}_\mu$ y $A$ es nulo, todos sus subconjuntos medibles son nulos, de modo que $|\nu|(A)=0$. Por ello una función acotada nula casi en todas partes tiene integral cero. Para una clase $f\in L^\infty$, elegimos un representante acotado suprimiendo sus valores mayores que $\norm f_\infty$ en el conjunto nulo correspondiente. La integral obtenida es independiente del representante acotado. En general esta integral no satisface convergencia monótona ni dominada; únicamente se ha utilizado continuidad uniforme.
\begin{teorema}[Dual de $L^\infty$]\label{v2:dual-linfty}
La aplicación $\nu\mapsto I_\nu$ es una isometría lineal sobreyectiva de $\operatorname{ba}_\mu(\Acal)$ sobre $(L^\infty(\mu))^*$.
\end{teorema}
\begin{proof}
La construcción anterior da $|I_\nu(f)|\le\norm\nu\norm f_\infty$, luego $\norm{I_\nu}\le\norm\nu$. Dada una partición $A_1,\ldots,A_m$, elegimos $c_j=\overline{\nu(A_j)}/|\nu(A_j)|$ cuando $\nu(A_j)\ne0$ y cero en otro caso. Entonces $s=\sum_jc_j\ind{A_j}$ tiene norma esencial a lo sumo uno y
\[
 I_\nu(s)=\sum_j|\nu(A_j)|.
\]
Tomar el supremo sobre particiones da $\norm{I_\nu}\ge\norm\nu$.

Si $\varphi\in(L^\infty)^*$, definimos $\nu(A)=\varphi(\ind A)$. La linealidad prueba aditividad finita. Si $A$ es nulo, su indicador es el vector cero de $L^\infty$, luego $\nu(A)=0$. Para toda partición y las fases anteriores,
$\sum_j|\nu(A_j)|=\varphi(s)\le\norm\varphi$; por tanto $\nu\in\operatorname{ba}_\mu$. El funcional y la integral coinciden sobre indicadores, después sobre simples y finalmente sobre todo $L^\infty$ por densidad en norma. Así la aplicación es sobreyectiva. Su isometría implica inyectividad.
\end{proof}

\section{Ejercicios propuestos: aditividad finita y ultrafiltros}
\begin{ejercicio}[Lectura 29, ejercicio 1: enunciado falso]\label{ej:v2-l29-1}
La lectura afirma: si $\nu$ es finitamente aditiva y $(A_n)$ es disjunta, la convergencia de $\sum_n\nu(A_n)$ implica que su suma es $\nu(\bigcup_nA_n)$. Examinarla.
\end{ejercicio}

\begin{ejercicio}[Lectura 29, ejercicio 2]\label{ej:v2-l29-2}
Para una carga de variación acotada, demostrar: (i) la integral de una simple no depende de su representación; (ii) la integral de una función medible acotada no depende de la sucesión de simples que converge uniformemente a ella.
\end{ejercicio}

\begin{ejercicio}[Lectura 29, ejercicio 3]\label{ej:v2-l29-3}
Sea $X\ne\varnothing$ y sea $\nu:\mathcal P(X)\to\{0,1\}$ finitamente aditiva, con $\nu(X)=1$. Para $\mathcal U=\{A:\nu(A)=1\}$, demostrar: (i) $\varnothing\notin\mathcal U$; (ii) $A\in\mathcal U$, $A\subset B$ implica $B\in\mathcal U$; (iii) $A,B\in\mathcal U$ implica $A\cap B\in\mathcal U$; (iv) si $A\notin\mathcal U$, entonces $X\setminus A\in\mathcal U$. Deducir una correspondencia biyectiva entre estas cargas y los ultrafiltros sobre $X$.
\end{ejercicio}

\begin{ejercicio}[Lectura 29, ejercicio 4: excepción del funcional cero]\label{ej:v2-l29-4}
Para $\nu\in\operatorname{ba}_\mu(\Acal)$, caracterizar la identidad
\[
 \int fg\,d\nu=\left(\int f\,d\nu\right)\left(\int g\,d\nu\right)
 \qquad(f,g\in L^\infty).
\]
La lectura afirma equivalencia con una carga cero-uno de masa uno.
\end{ejercicio}

\begin{ejercicio}[Lectura 29, ejercicio 5]\label{ej:v2-l29-5}
Un ultrafiltro $\mathcal U$ es libre si $\bigcap\mathcal U=\varnothing$, y fijo en otro caso. Demostrar: (i) si es fijo, existe $x$ con $\mathcal U=\{A:x\in A\}$; (ii) en un conjunto finito todos son fijos, y en uno infinito existen libres; (iii) la carga asociada a uno libre se anula en los conjuntos finitos; (iv) para un ultrafiltro libre en $\N$ y una sucesión real acotada $x$,
\[
 \liminf_nx_n\le\int_\N x\,d\nu\le\limsup_nx_n.
\]
Deducir que la integral extiende el límite de las sucesiones convergentes.
\end{ejercicio}

\section{Dos representaciones, dos clases de objetos}
El dual de $L^1$ se representa mediante funciones esencialmente acotadas cuando la medida permite pegar datos locales. El dual de $L^\infty$ se representa siempre mediante cargas de variación acotada que respetan los conjuntos nulos. Las cargas de ultrafiltros muestran por qué sustituir aditividad finita por aditividad numerable cambiaría ese dual y perdería funcionales continuos.

\parte{Medidas de Borel y funcionales continuos}{La topología permite aproximar conjuntos por abiertos y compactos y funciones medibles por continuas. Tras establecer las herramientas de separación necesarias, construiremos las medidas que representan funcionales continuos. Se distinguirán los resultados válidos en espacios localmente compactos generales de los que requieren segundo numerabilidad o $\sigma$-compacidad.}
\chapter{Medidas de Borel y regularidad}\label{v2:cap15}\label{cap:v2-15}
\section{Relacionar la medida con la topología}\label{lectura:v2-31}
La medida y la topología son estructuras diferentes. La regularidad permite aproximar conjuntos medibles por abiertos y compactos, y con ello aproximar funciones medibles por funciones continuas. Desarrollamos las lecturas 31 y 32 de \cite{valencia2}. Salvo indicación contraria, $X$ es localmente compacto y Hausdorff. Los resultados que necesitan segundo numerabilidad o $\sigma$-compacidad lo indicarán expresamente.
\begin{definicion}
$\mathcal B(X)$ es la $\sigma$-álgebra generada por los abiertos. Una medida de Borel está definida sobre ella. En una $\sigma$-álgebra $\Acal\supset\mathcal B(X)$, llamaremos \emph{regular} a una medida positiva $\mu$ que cumple:
\begin{enumerate}
\item $\mu(K)<\infty$ para todo compacto $K$;
\item $\mu(A)=\inf\{\mu(U):A\subset U,\ U\text{ abierto}\}$ para todo $A\in\Acal$;
\item $\mu(U)=\sup\{\mu(K):K\subset U,\ K\text{ compacto}\}$ para todo abierto $U$.
\end{enumerate}
Estas son finitud sobre compactos, regularidad exterior y regularidad interior sobre abiertos. No sustituimos la última condición por regularidad interior sobre todos los medibles sin justificar una hipótesis adicional.\index{medida!de Borel}\index{regularidad de una medida}
\end{definicion}
Lebesgue sobre $\R$ es regular. En un espacio discreto, los compactos son finitos: la cubierta por singletons demuestra la necesidad y los finitos son compactos. Por ello la medida de conteo es finita sobre compactos, exteriormente regular porque todos los conjuntos son abiertos e interiormente regular porque el cardinal de un conjunto es el supremo de los cardinales de sus partes finitas. La medida de Dirac $\delta_x$ también es regular: es finita; si $x\notin A$, el abierto $X\setminus\{x\}$ contiene a $A$ y tiene medida cero; si $x\in A$, cualquier abierto que lo contiene tiene medida uno. Para regularidad interior de un abierto que contiene $x$ basta el compacto $\{x\}$.

\section{Las herramientas topológicas necesarias}
\begin{proposicion}\label{v2:topologia-lch}
En un espacio localmente compacto Hausdorff:
\begin{enumerate}
\item dos compactos disjuntos tienen vecindades abiertas disjuntas;
\item si $K\subset U$, con $K$ compacto y $U$ abierto, existe $V$ abierto con $K\subset V\subset\cl V\subset U$ y $\cl V$ compacto;
\item existe $h\in C_c(X)$ con $0\le h\le1$, $h=1$ en $K$ y $\supp h\subset U$;
\item una función continua sobre un compacto $K$ se extiende a una función de $C_c(X)$, con soporte contenido en cualquier abierto $U\supset K$ y sin aumentar su supremo absoluto;
\item si $X$ es segundo numerable, todo abierto es unión numerable de compactos y $X$, así como su compactificación unipuntual, es metrizable.
\end{enumerate}
Aquí $C_c(X)$, también denotado $C_{00}(X)$ en las lecturas, es el espacio de funciones continuas de soporte compacto.
\end{proposicion}
\begin{proof}
\emph{Separación y reducción de vecindades.} Para un punto $x$ y un compacto que no lo contiene, la propiedad Hausdorff da vecindades disjuntas de $x$ y de cada punto del compacto. Extraer una cubierta finita y tomar una intersección para la vecindad de $x$ separa el punto del compacto. Repetir este argumento sobre un segundo compacto y extraer otra cubierta finita prueba (i).

Para $x\in U$, tomemos una vecindad compacta $L$ de $x$. El compacto $L\setminus U$ se separa de $x$ por abiertos disjuntos $O,W$, con $L\setminus U\subset W$. Entonces $V_x=O\cap\interior L$ contiene a $x$ y $\cl{V_x}\subset L\setminus W\subset U$; su cerradura es compacta. Cubrir $K$ por finitos $V_x$ y unirlos da (ii).

\emph{Separación mediante funciones.} Recordemos por qué un compacto Hausdorff admite las funciones que necesitamos. Sus cerrados son compactos, y (i) lo hace normal. En un espacio normal, si un cerrado $A$ está contenido en un abierto $O$, hay un abierto $V$ con $A\subset V\subset\cl V\subset O$: se separan $A$ y el cerrado complementario. Usando esto inductivamente, para dos cerrados disjuntos $A,B$ se construyen abiertos $U_r$, indexados por racionales diádicos de $[0,1]$, con $A\subset U_0$, $U_1=X\setminus B$ y $\cl{U_r}\subset U_s$ si $r<s$. Entre dos índices consecutivos de una etapa se inserta el abierto obtenido por la reducción anterior. La función
\[
 u(x)=\inf\{r:x\in U_r\},\qquad\inf\varnothing=1,
\]
vale cero en $A$ y uno en $B$. Es continua: $\{u<t\}=\bigcup_{r<t}U_r$ y $\{u>t\}=\bigcup_{r>t}(X\setminus\cl{U_r})$. Esto prueba el lema de Urysohn en el compacto.

La compactificación unipuntual $X^*=X\cup\{\infty\}$ tiene como vecindades de $\infty$ los complementos de compactos de $X$. Es compacta: un abierto de una cubierta que contenga a $\infty$ deja por cubrir un compacto. Es Hausdorff: (ii) separa un punto de $X$ del punto añadido mediante una vecindad de cerradura compacta. Si $X$ ya es compacto, el punto añadido es aislado y las mismas afirmaciones valen. Para probar (iii), elegimos $K\subset V\subset\cl V\subset U$ como en (ii), y aplicamos Urysohn en $X^*$ a $K$ y $X^*\setminus V$. La función que vale uno en el primero y cero en el segundo tiene soporte contenido en $\cl V$, como se requiere.

\emph{Extensión desde un compacto.} Para una función real $f$ continua en $K$, con $|f|\le M$, Urysohn aplicado a los cerrados de $K$ donde $f\le-M/3$ y $f\ge M/3$ proporciona una función continua $g_0$ en $X^*$ con $\norm{g_0}_{\mathrm{sup}}\le M/3$ y $|f-g_0|\le2M/3$ en $K$. Aplicamos el mismo procedimiento al residuo y obtenemos $g_n$ con norma a lo sumo $M(2/3)^n/3$ y residuo después de la etapa $n$ a lo sumo $M(2/3)^{n+1}$. La serie converge uniformemente a una función continua que extiende $f$ y tiene norma a lo sumo $M$. Para una función compleja extendemos partes real e imaginaria y componemos con la proyección radial $z\mapsto z\min\{1,M/|z|\}$, definida como cero en cero; conserva los valores de módulo a lo sumo $M$ y no aumenta el módulo. Multiplicamos finalmente por una función de corte del punto (iii), igual a uno en $K$ y soportada en $U$. Esto prueba (iv). No se ha supuesto que todo espacio localmente compacto Hausdorff sea normal; la extensión se hizo en el compacto $X^*$.

\emph{Numerabilidad y métrica.} De una base numerable de $X$ retenemos los miembros de cerradura compacta. Por (ii) todavía forman una base, y cualquier abierto es unión de una subfamilia numerable de ellos con cerraduras contenidas en ese abierto. Sus uniones finitas de cerraduras dan la cubierta compacta creciente. En particular existen compactos crecientes $K_n$ que cubren $X$ y cuyas partes interiores también lo cubren. Sus complementos en $X^*$ forman una base numerable en $\infty$: cualquier compacto de $X$ está contenido en algún $K_n$ por una subcubierta finita de los interiores. Así $X^*$ es compacto Hausdorff y segundo numerable.

Para construir una métrica, tomamos todos los pares de elementos $B_i,B_j$ de una base numerable de $X^*$ con $\cl{B_i}\subset B_j$, y una función continua $u_{ij}$ igual a uno en $\cl{B_i}$ y cero fuera de $B_j$. Estas funciones separan puntos: para $x\ne y$ se reducen vecindades y se elige un par con $x\in B_i$, $y\notin B_j$. La aplicación $x\mapsto(u_{ij}(x))$ es continua e inyectiva en el cubo numerable, metrizado por $\sum_n2^{-n}|a_n-b_n|$. Como su dominio es compacto y el codominio Hausdorff, es un homeomorfismo sobre su imagen: las imágenes de cerrados compactos son cerradas. Transportar la métrica a $X^*$ y restringirla a $X$ termina la prueba.
\end{proof}
\begin{lema}[Partición finita cerca de un compacto]\label{v2:particion-compacta}
Si abiertos $U_1,\ldots,U_m$ cubren un compacto $K$, existen $\theta_j\in C_c(X)$, $0\le\theta_j\le1$, $\supp\theta_j\subset U_j$, con $\sum_j\theta_j=1$ en $K$ y $\sum_j\theta_j\le1$ en todo $X$.
\end{lema}
\begin{proof}
Para cada punto de $K$ elegimos un abierto de cerradura compacta contenido en algún $U_j$ y una función de corte igual a uno en esa cerradura, con soporte en $U_j$. Finitos de esos abiertos cubren $K$. Sumamos las funciones elegidas según el $U_j$ que les corresponde, obteniendo $g_j\ge0$ con soporte compacto en $U_j$. Si $S=\sum_jg_j$, entonces $S\ge1$ en $K$. Definimos $\theta_j=g_j/\max\{1,S\}$. El denominador es continuo y positivo; las cotas, los soportes y la identidad en $K$ se comprueban directamente.
\end{proof}

\section{Regularidad y conjuntos de medida finita}
\begin{lema}\label{v2:borel-finita-regular}
En un espacio localmente compacto Hausdorff segundo numerable, una medida de Borel finita satisface, para todo Borel $A$,
\[
 \mu(A)=\inf_{A\subset U\text{ abierto}}\mu(U)
 =\sup_{F\subset A\text{ cerrado}}\mu(F).
\]
\end{lema}
\begin{proof}
Sea $\mathcal R$ la familia de Borel $A$ que para cada $\eps>0$ admiten un cerrado $F\subset A$ y un abierto $U\supset A$ con $\mu(U\setminus F)<\eps$. Todo abierto pertenece: por la proposición \ref{v2:topologia-lch} es unión creciente de compactos, y continuidad de la medida finita aproxima su masa desde dentro. La familia es estable por complementos, intercambiando $F,U$ y tomando sus complementos.

También es estable por uniones numerables. Para $A=\bigcup_nA_n$, elegimos abiertos $U_n\supset A_n$ con $\mu(U_n\setminus A_n)<\eps2^{-n-2}$; su unión es una aproximación exterior. Como $\mu$ es finita, una unión finita de los $A_n$ deja fuera una parte de $A$ de medida menor que $\eps/4$. Aproximamos cada uno de esos finitos miembros desde dentro por cerrados, con error total menor que $\eps/4$; su unión cerrada aproxima $A$ desde dentro. El error conjunto es menor que $\eps$. Así $\mathcal R$ es una $\sigma$-álgebra que contiene los abiertos, por lo que contiene todos los Borel. Las igualdades se siguen de la monotonía y de hacer $\eps\downarrow0$.
\end{proof}
\begin{teorema}\label{v2:borel-localmente-finita}
En un espacio localmente compacto Hausdorff segundo numerable, toda medida de Borel finita sobre compactos es regular.
\end{teorema}
\begin{proof}
Cada abierto es unión creciente de compactos, luego la continuidad creciente da regularidad interior en abiertos. Para regularidad exterior elegimos una cubierta numerable por abiertos $U_n$ de cerradura compacta; cada uno tiene medida finita. Las medidas $\mu_n(A)=\mu(A\cap U_n)$ son finitas y satisfacen el lema anterior. Para un Borel $A$, elegimos abiertos $V_n\supset A\cap U_n$ tales que $\mu((V_n\cap U_n)\setminus A)<\eps2^{-n}$. La unión $V=\bigcup_n(V_n\cap U_n)$ es abierta, contiene a $A$ y satisface $\mu(V\setminus A)<\eps$. Por tanto $\mu(V)\le\mu(A)+\eps$, con la interpretación habitual si $\mu(A)=\infty$.
\end{proof}
\begin{proposicion}\label{v2:regular-interior-sigma}
Para una medida regular en el sentido de este capítulo, si $A\in\Acal$ es de medida $\sigma$-finita, entonces
\[
 \mu(A)=\sup\{\mu(K):K\subset A,\ K\text{ compacto}\}.
\]
\end{proposicion}
\begin{proof}
Si $\mu(A)<\infty$, elegimos un abierto $U\supset A$ con $\mu(U\setminus A)<\eps$. Elegimos un compacto $L\subset U$ con $\mu(U\setminus L)<\eps$. La regularidad exterior aplicada a $U\setminus A$ da un abierto $W$ que lo contiene y cumple $\mu(W)<2\eps$. El compacto $K=L\setminus W$ está contenido en $A$, y $\mu(A\setminus K)\le\mu(U\setminus L)+\mu(W)<3\eps$. Si $A$ es $\sigma$-finito de medida infinita, escribimos $A_n\uparrow A$ con medidas finitas; estas crecen a infinito. Para cualquier $M$ elegimos $n$ con $\mu(A_n)>M+1$ y un compacto $K\subset A_n$ con $\mu(K)>M$. El supremo es infinito.
\end{proof}

\section{Ejercicios resueltos: conjuntos de Borel y regularidad}
\begin{ejercicio}[Lectura 31, ejercicio 1]\label{ej:v2-l31-1}
Demostrar que una función continua entre espacios de Hausdorff es Borel-medible.
\end{ejercicio}
\begin{solucion}
La familia de subconjuntos $A$ del codominio cuya preimagen es Borel es una $\sigma$-álgebra, porque la preimagen respeta complementos y uniones numerables. Contiene los abiertos por continuidad; por tanto contiene la $\sigma$-álgebra que generan, esto es, todos los Borel. La propiedad Hausdorff no es necesaria para este argumento.
\end{solucion}
\begin{ejercicio}[Lectura 31, ejercicio 2]\label{ej:v2-l31-2}
Para un subespacio $Y\subset X$, demostrar $\mathcal B(Y)=\{B\cap Y:B\in\mathcal B(X)\}$.
\end{ejercicio}
\begin{solucion}
Las trazas de los Borel de $X$ forman una $\sigma$-álgebra en $Y$ que contiene todos sus abiertos relativos, así que contiene $\mathcal B(Y)$. Para la inclusión inversa, la inclusión $Y\hookrightarrow X$ es continua, y por el ejercicio anterior la preimagen de cada Borel $B$ de $X$, que es $B\cap Y$, pertenece a $\mathcal B(Y)$. No se necesita que $Y$ sea Borel en $X$.
\end{solucion}
\begin{ejercicio}[Lectura 31, ejercicio 3]\label{ej:v2-l31-3}
Demostrar el lema de regularidad por abiertos y cerrados para una medida de Borel finita en un espacio localmente compacto Hausdorff segundo numerable.
\end{ejercicio}
\begin{solucion}
Es el lema \ref{v2:borel-finita-regular}, cuya prueba completa se incluyó antes de utilizarlo. Su paso decisivo es que cada abierto es $F_\sigma$, en realidad unión creciente de compactos. La finitud de la medida permite hacer pequeño el resto de una unión creciente y permite intercambiar aproximación interior y exterior tomando complementos. La familia de conjuntos con ambas aproximaciones es entonces una $\sigma$-álgebra; por eso la propiedad pasa de abiertos a todos los Borel.
\end{solucion}
\begin{ejercicio}[Lectura 31, ejercicio 4]\label{ej:v2-l31-4}
Si $A\in\Acal$ es de medida $\sigma$-finita, demostrar que para cada $\eps>0$ existe un abierto $U\supset A$ con $\mu(U\setminus A)<\eps$, incluso cuando $\mu(A)=\infty$.
\end{ejercicio}
\begin{solucion}
Escribimos $A=\bigcup_nA_n$ con cada $\mu(A_n)<\infty$. La regularidad exterior proporciona abiertos $U_n\supset A_n$ con $\mu(U_n\setminus A_n)<\eps2^{-n}$. Para $U=\bigcup_nU_n$, tenemos $U\setminus A\subset\bigcup_n(U_n\setminus A_n)$ y por subaditividad $\mu(U\setminus A)<\eps$. No hemos restado dos cantidades infinitas.
\end{solucion}
\begin{ejercicio}[Lectura 31, ejercicio 5: alcance de la hipótesis]\label{ej:v2-l31-5}
Para $A\in\Acal$ de medida $\sigma$-finita, construir Borel $E\subset A\subset F$ con $\mu(F\setminus E)=0$. Deducir que, si $\mu$ es $\sigma$-finita, $\Acal$ está contenida en la completación de $\mathcal B(X)$ para la restricción de $\mu$. Examinar la omisión de la condición sobre $A$ en el primer enunciado original.
\end{ejercicio}
\begin{solucion}
Si $A=\bigcup_nA_n$ con medidas finitas, aproximamos cada $A_n$ por compactos $K_{n,j}\subset A_n$ con $\mu(A_n\setminus K_{n,j})<1/j$. La unión Borel $E=\bigcup_{n,j}K_{n,j}$ está en $A$ y $\mu(A\setminus E)=0$. Por el ejercicio \ref{ej:v2-l31-4}, elegimos abiertos $U_j\supset A$ con $\mu(U_j\setminus A)<1/j$ y ponemos $F=\bigcap_jU_j$. Entonces $\mu(F\setminus A)=0$, de donde $\mu(F\setminus E)=0$. Si todo $X$ es $\sigma$-finito, el argumento vale para cada $A\in\Acal$ y da precisamente su pertenencia a la completación de los Borel.

La condición no puede eliminarse en general. Sea $D$ discreto no numerable, $X=\R\times D$, y sea $\Acal$ la familia de conjuntos cuyas secciones $A_d\subset\R$ son Lebesgue-medibles. Definimos $\mu(A)=\sum_d m(A_d)$ si $A$ toca solo numerables componentes, y $\mu(A)=\infty$ en otro caso. Es una medida: una unión numerable toca no numerables componentes solo si alguno de sus miembros lo hace; en el otro caso la aditividad se reduce a las sumas numerables de las medidas de las secciones. Es regular: los compactos tocan finitas componentes y son acotados en ellas; los conjuntos de medida finita se aproximan por abiertos sección a sección con errores sumables; para los de medida infinita la regularidad exterior es automática. La regularidad interior de abiertos se sigue de que toda sección abierta no vacía tiene medida positiva, aproximable por compactos.

Elijamos un conjunto Lebesgue-medible nulo $V\subset\R$ que no sea Borel, y pongamos $A=V\times D$. Tal $V$ existe: el conjunto de Cantor es nulo, tiene $2^{\aleph_0}$ puntos y posee más subconjuntos que conjuntos Borel en $\R$. Si Borel $E\subset A\subset F$ cumplieran $\mu(F\setminus E)=0$, fuera de numerables componentes sus secciones tendrían que coincidir exactamente con $V$. Pero las secciones de un Borel son Borel, por continuidad de la inclusión de cada componente, contradicción.
\end{solucion}
\begin{ejercicio}[Lectura 31, ejercicio 6]\label{ej:v2-l31-6}
Demostrar que la completación de una medida regular es regular.
\end{ejercicio}
\begin{solucion}
Un conjunto de la completación está entre dos conjuntos originales $B\subset A\subset C$ con $\mu(C\setminus B)=0$ y tiene medida $\mu(B)=\mu(C)$. La regularidad exterior de $C$ proporciona la misma aproximación exterior para $A$. Las medidas de abiertos y compactos no cambian al completar, así que se conservan finitud en compactos y regularidad interior de abiertos. Esto prueba exactamente las tres condiciones de la definición, sin exigir regularidad interior de conjuntos de medida no $\sigma$-finita.
\end{solucion}
\begin{ejercicio}[Lectura 31, ejercicio 7]\label{ej:v2-l31-7}
Supóngase $\mu(\{x\})=0$ para todo $x$. Para un Borel $B$ de medida finita y $0<\alpha<\mu(B)$, construir un Borel $A\subset B$ con $\mu(A)=\alpha$.
\end{ejercicio}
\begin{solucion}
Primero, para cada $\eta>0$, $B$ admite una partición finita en Borel de medida menor que $\eta$. Elegimos un compacto $K\subset B$ con $\mu(B\setminus K)<\eta$. Cada $x\in K$ tiene un abierto $U_x$ de medida menor que $\eta$, por regularidad exterior de $\{x\}$. Una subcubierta finita de $K$, disjuntada e intersectada con $K$, y la pieza $B\setminus K$ dan la partición.

Construimos $A_n\uparrow A$ con $\mu(A_n)\le\alpha$ y $\alpha-\mu(A_n)<2^{-n}$. Dado $A_{n-1}$, particionamos $B\setminus A_{n-1}$ en piezas de medida menor que $2^{-n}$. Añadimos piezas mientras su suma no supere el déficit $\alpha-\mu(A_{n-1})$. Si el déficit final fuera al menos $2^{-n}$, cada pieza restante cabría, contradiciendo que dejamos de añadir; además hay masa total suficiente porque $\mu(B)>\alpha$. La unión finita elegida define $A_n$. La unión numerable $A=\bigcup_nA_n$ es Borel y, por continuidad creciente, tiene medida $\alpha$.
\end{solucion}
\begin{ejercicio}[Lectura 31, ejercicio 8]\label{ej:v2-l31-8}
Para $X$ compacto Hausdorff, sea $\mathcal B_0(X)$ la menor $\sigma$-álgebra que hace medibles todas las funciones continuas reales; sus miembros son conjuntos de Baire. Demostrar: (i) está generada por los cerrados $G_\delta$; (ii) si $X$ es segundo numerable, $\mathcal B_0(X)=\mathcal B(X)$.
\end{ejercicio}
\begin{solucion}
(i) Para $f$ continua, $\{f\le a\}=\bigcap_n\{f<a+1/n\}$ es cerrado y $G_\delta$. Estos subniveles generan la medibilidad de $f$, por lo que $\mathcal B_0$ está contenida en la $\sigma$-álgebra indicada. Si $F=\bigcap_nU_n$ es cerrado y $G_\delta$, Urysohn proporciona $h_n:X\to[0,1]$ continua, igual a uno en $F$ y cero fuera de $U_n$. La serie uniformemente convergente $h=\sum_n2^{-n}(1-h_n)$ es continua y $F=\{h=0\}$. Así cada generador pertenece a $\mathcal B_0$, y hay igualdad.

(ii) En un compacto segundo numerable, todo abierto es $F_\sigma$ por la proposición \ref{v2:topologia-lch}; tomando complementos, todo cerrado es $G_\delta$. El punto (i) muestra que $\mathcal B_0$ contiene todos los cerrados y, por tanto, todos los Borel. La inclusión contraria sigue de que toda función continua es Borel-medible. Una medida de Baire es una medida finita definida en $\mathcal B_0$.
\end{solucion}

\section{Aproximar funciones por funciones continuas}\label{lectura:v2-32}
\begin{teorema}\label{v2:cc-denso-lp}
Si $\mu$ es regular, $C_c(X)$ es denso en $L^p(\mu)$ para $1\le p<\infty$.
\end{teorema}
\begin{proof}
Las simples integrables son densas; basta aproximar $\ind A$ para $\mu(A)<\infty$. Por regularidad exterior e interior en conjuntos finitos, elegimos $K\subset A\subset U$, compacto y abierto respectivamente, con $\mu(U\setminus K)<\eps$. La función de corte $h$, igual a uno en $K$, entre cero y uno y soportada en $U$, satisface $|h-\ind A|^p\le\ind{U\setminus K}$. Así $\norm{h-\ind A}_p^p<\eps$. Aproximar por separado los finitos indicadores de una simple y usar Minkowski concluye.
\end{proof}
El resultado falla en $L^\infty$. En $[0,1]$ con Lebesgue, ninguna función continua aproxima $\ind{[0,1/2)}$ con error esencial menor que $1/2$: la continuidad obligaría a tener, en el punto $1/2$, simultáneamente un límite por la izquierda mayor que $1/2$ y un límite por la derecha menor que $1/2$. Más precisamente, una cota esencial estricta se convierte en la misma cota puntual dentro de cada intervalo abierto por continuidad y positividad de la medida de sus vecindades. No basta que un abierto y su complemento tengan interiores no vacíos para producir este fenómeno: un abierto y cerrado tiene indicador continuo.
\begin{teorema}[Lusin]\label{v2:lusin}\index{Lusin, teorema de}
Sea $f:X\to\F$ medible y supóngase $A=\{f\ne0\}$ de medida finita. Para cada $\eps>0$ existen un compacto $K\subset A$ y $h\in C_c(X)$ tales que $\mu(A\setminus K)<\eps$ y $h=f$ en $K$.
\end{teorema}
\begin{proof}
Si $f$ es acotada, pertenece a $L^1$ porque su soporte medible tiene medida finita. Por densidad existe una sucesión de $C_c$ que converge a $f$ en $L^1$; extraemos una subsucesión convergente casi en todas partes. Egorov en $A$ produce un conjunto medible $B\subset A$ con $\mu(A\setminus B)<\eps/2$ donde la convergencia es uniforme. La regularidad interior sobre conjuntos finitos da un compacto $K\subset B$ con $\mu(B\setminus K)<\eps/2$. El límite uniforme sobre $K$ es continuo, de modo que $f|_K$ es continua. La extensión desde compactos de la proposición \ref{v2:topologia-lch} produce $h\in C_c$ que coincide con ella.

Si $f$ no es acotada, los conjuntos $A_n=\{0<|f|\le n\}$ crecen a $A$ y tienen medida finita. Elegimos $n$ con $\mu(A\setminus A_n)<\eps/2$ y aplicamos el caso anterior a $f\ind{A_n}$ con error $\eps/2$. El compacto obtenido está contenido en $A_n$, donde la función truncada coincide con $f$. Esto prueba el enunciado sin suponer normalidad global de $X$.
\end{proof}

\section{Localización estricta: definición y alcance}
\begin{definicion}
Un espacio de medida es \emph{estrictamente localizable} si admite una partición medible $(X_i)_{i\in I}$ con $\mu(X_i)<\infty$ tal que un subconjunto $A\subset X$ es medible si y solo si todas sus secciones $A\cap X_i$ lo son, y, para \emph{todo} medible $A$,
\[
 \mu(A)=\sum_{i\in I}\mu(A\cap X_i),
\]
donde la suma no numerable es el supremo de las sumas finitas.\index{medida!estrictamente localizable}
\end{definicion}
La condición sobre todos los medibles es esencial. La lectura la exige solamente cuando $\mu(A)<\infty$, pero con esa versión no se deduce semifinitud: sobre un conjunto no numerable con todas sus partes medibles, la medida que vale cero en conjuntos numerables e infinito en los demás, particionada en singletons, satisface las condiciones impresas y no es semifinita. Usaremos la definición corregida.
\begin{proposicion}\label{v2:estricta-localizable}
Toda medida $\sigma$-finita es estrictamente localizable; toda medida estrictamente localizable es semifinita y localizable.
\end{proposicion}
\begin{proof}
En el caso $\sigma$-finito, disjuntamos una cubierta numerable por conjuntos de medida finita. La medibilidad se reconstruye por unión numerable y la igualdad de masas es la aditividad numerable. Para una partición estricta general, si $\mu(A)>0$, algún $A\cap X_i$ tiene medida positiva y finita; esto prueba semifinitud cuando $\mu(A)=\infty$.

Dada una familia $\mathcal E$ de medibles, en cada espacio finito $X_i$ elegimos un supremo esencial $H_i$ de las secciones $E\cap X_i$, por la proposición \ref{v2:sigma-localizable}. La unión $H=\bigcup_iH_i$ es medible por la condición de secciones. Para cada $E\in\mathcal E$, la fórmula de sumas da $\mu(E\setminus H)=\sum_i0=0$. Si $G$ es otro mayorante esencial, su sección en cada $X_i$ contiene esencialmente a $H_i$; la misma fórmula da $\mu(H\setminus G)=0$. Así $H$ es un supremo esencial global.
\end{proof}
\begin{teorema}[Descomposición por compactos bajo $\sigma$-compacidad]\label{v2:compactos-localizacion}
Supóngase además que $X$ es $\sigma$-compacto y que la medida regular está completada. Existe una familia a lo sumo numerable $\mathcal F$ de compactos disjuntos, cada uno de medida positiva finita y con la propiedad
\[
 U\text{ abierto},\ U\cap F\ne\varnothing\quad\Longrightarrow\quad\mu(U\cap F)>0,
\]
tal que $D=X\setminus\bigcup\mathcal F$ es nulo. Todo conjunto de medida finita toca a lo sumo numerables miembros de $\mathcal F$, y $B\subset X$ es medible si y solo si lo son sus intersecciones con todos los $F\in\mathcal F$. Añadiendo $D$ se obtiene una descomposición estricta.
\end{teorema}
\begin{proof}
Ordenamos por inclusión las familias de compactos disjuntos de medida positiva que satisfacen la propiedad indicada. La unión de una cadena sigue siendo una familia de ese tipo, por lo que Zorn da una maximal $\mathcal F$. Si $A$ tiene medida finita, la regularidad exterior da un abierto $U\supset A$ de medida finita. Cada compacto de la familia que toca $A$ toca $U$ y aporta una cantidad positiva $\mu(U\cap F)$. Solo puede haber numerables cantidades positivas disjuntas en un espacio de medida finita: para cada $n$, las mayores que $1/n$ son finitas. Esto prueba la afirmación sobre $A$. Como $X$ es unión numerable de compactos de medida finita, la familia $\mathcal F$ es numerable. En consecuencia $D$ es Borel.

Si $\mu(D)>0$, la medida es $\sigma$-finita porque $X$ es $\sigma$-compacto, y la regularidad interior proporciona un compacto $K\subset D$ de medida positiva. Sea $O$ la unión de los abiertos $U$ con $\mu(U\cap K)=0$. Entonces $\mu(O\cap K)=0$: cualquier compacto contenido en $O\cap K$ se cubre con finitos de esos abiertos y es nulo; la regularidad interior en el espacio compacto $K$ da la afirmación para su abierto relativo. Así $F=K\setminus O$ es compacto de medida positiva. Si un abierto $V$ toca $F$, no puede tener $\mu(V\cap K)=0$, pues en tal caso estaría en $O$. Como se retiró un conjunto nulo, $\mu(V\cap F)>0$. El compacto $F$ es disjunto de todos los anteriores y contradice maximalidad. Por tanto $D$ es nulo.

Si todas las secciones $B\cap F$ son medibles, entonces $B=(B\cap D)\cup\bigcup_{F\in\mathcal F}(B\cap F)$ es medible: la unión es numerable y la completitud hace medible cualquier subconjunto de $D$. La implicación inversa es inmediata. Las sumas de medidas se reconstruyen por aditividad numerable. Esto prueba todas las propiedades de la descomposición.
\end{proof}
\begin{observacion}[Por qué la hipótesis añadida importa]
La lectura 32 anuncia una descomposición de este tipo para todo espacio localmente compacto Hausdorff con medida regular, y remite sus detalles a lecturas posteriores no presentes en la carpeta. Con la convención de regularidad exterior usada allí, su conclusión de localizabilidad general es falsa. En $X=\R\times D$, $D$ discreto no numerable, tomemos la medida de Borel descrita en el ejercicio \ref{ej:v2-l31-5}: masa infinita si un conjunto toca no numerables componentes, y suma de medidas de Lebesgue en otro caso. Es finita sobre compactos, exteriormente regular e interiormente regular sobre abiertos. Sin embargo, $H=\{0\}\times D$ tiene medida infinita y cada uno de sus subconjuntos de medida finita es nulo. No es semifinita y, por tanto, no es localizable. Completarla no cambia esos valores ni repara el problema. El teorema anterior es una versión plenamente demostrada que incluye los espacios segundo numerables utilizados habitualmente en análisis; una ampliación a espacios generales requiere otra noción de medida de Radon o una extensión localmente determinada, y no puede identificarse con la completación ordinaria.
\end{observacion}

\section{Ejercicios resueltos: Lusin y soporte de una medida}
\begin{ejercicio}[Lectura 32, ejercicio 1]\label{ej:v2-l32-1}
Para una medida regular y un medible $A$ de medida $\sigma$-finita, demostrar $\mu(A)=\sup\{\mu(K):K\subset A,\ K\text{ compacto}\}$.
\end{ejercicio}
\begin{solucion}
La proposición \ref{v2:regular-interior-sigma} prueba exactamente este resultado. Para medida finita se retira de un compacto aproximante un abierto de pequeña medida que cubre la parte exterior a $A$. Para medida infinita $\sigma$-finita se aplican esas aproximaciones a conjuntos finitos de masa arbitrariamente grande. El ejemplo $H$ anterior muestra por qué no se debe suprimir la condición de $\sigma$-finitud.
\end{solucion}
\begin{ejercicio}[Lectura 32, ejercicio 2]\label{ej:v2-l32-2}
Sea $(x_n)$ una enumeración de $\Q\cap[0,1]$ y $\eps>0$. Para $A_n=(x_n-\eps2^{-n},x_n+\eps2^{-n})$ y $F=[0,1]\setminus\bigcup_nA_n$, demostrar que $\ind{\Q\cap[0,1]}|_F$ es continua y $m(F)\ge1-2\eps$.
\end{ejercicio}
\begin{solucion}
Cada racional de $[0,1]$ pertenece a alguno de los intervalos, por lo que $F$ solo contiene irracionales. La restricción del indicador es la función constante cero, que es continua. Además $m([0,1]\setminus F)\le\sum_n2\eps2^{-n}=2\eps$, de donde la cota. El conjunto $F$ es cerrado en $[0,1]$, luego compacto. Si $2\eps>1$, la cota sigue siendo válida aunque sea menos informativa.
\end{solucion}
\begin{ejercicio}[Lectura 32, ejercicio 3]\label{ej:v2-l32-3}
Demostrar que la función $h$ de Lusin puede elegirse con $\norm h_{\mathrm{sup}}\le\norm f_{\mathrm{sup}}$.
\end{ejercicio}
\begin{solucion}
Si $M=\norm f_{\mathrm{sup}}<\infty$, componemos una extensión con $P_M(z)=z\min\{1,M/|z|\}$, con valor cero en cero. Esta aplicación es continua, no aumenta el módulo y deja fijos todos los valores de módulo a lo sumo $M$. Por ello conserva la igualdad en $K$, mantiene el soporte compacto y da la cota. Si $M=0$, basta $h=0$. Si $M=\infty$, cualquier función obtenida por Lusin satisface la desigualdad extendida.
\end{solucion}
\begin{ejercicio}[Lectura 32, ejercicio 4]\label{ej:v2-l32-4}
Demostrar que la unión $O$ de todos los abiertos de medida cero es abierta y tiene medida cero. Definir el soporte de $\mu$.
\end{ejercicio}
\begin{solucion}
La unión es abierta por definición de topología. Si $K\subset O$ es compacto, una cantidad finita de esos abiertos nulos lo cubre, así que $\mu(K)=0$. La regularidad interior del abierto $O$ da $\mu(O)=\sup_{K\subset O}\mu(K)=0$. Se define $\supp\mu=X\setminus O$, un cerrado. Para una medida real con signo o compleja regular se usa $\supp\nu=\supp|\nu|$. La prueba no usa subaditividad para una unión no numerable.\index{soporte de una medida}
\end{solucion}
\begin{ejercicio}[Lectura 32, ejercicio 5: positividad necesaria]\label{ej:v2-l32-5}
Demostrar que $x\in\supp\mu$ si y solo si toda función \emph{no negativa} $f\in C_c(X)$ con $f(x)>0$ satisface $\int f\,d\mu>0$. La condición $f\ge0$ falta en el enunciado original.
\end{ejercicio}
\begin{solucion}
Si $x$ está en el soporte y $f(x)>0$, la continuidad da un abierto $U\ni x$ y $a>0$ con $f\ge a$ en $U$. Ese abierto tiene medida positiva, pues de lo contrario $x$ no estaría en el soporte. Por tanto $\int f\ge a\mu(U)>0$. Si $x$ no está en el soporte, hay un abierto nulo $U\ni x$ y una función de corte no negativa, de soporte en $U$, con valor uno en $x$. Su integral es cero, lo que contradice la propiedad propuesta.

Sin no negatividad, la afirmación es falsa: en $\R$ con Lebesgue, dos funciones triangulares no negativas de igual integral y soportes disjuntos producen, al restarlas, una función de integral cero que es positiva en el centro del primer triángulo. Ese centro sí pertenece al soporte de Lebesgue.
\end{solucion}
\begin{ejercicio}[Lectura 32, ejercicio 6]\label{ej:v2-l32-6}
Sea $D$ discreto infinito no numerable y $D^*$ su compactificación unipuntual. Demostrar que ninguna medida de Borel regular sobre $D^*$ tiene soporte igual a $D^*$.
\end{ejercicio}
\begin{solucion}
Como $D^*$ es compacto, una medida regular allí es finita. Si su soporte fuera todo el espacio, cada singleton abierto $\{d\}$, $d\in D$, tendría medida positiva. En una medida finita solo hay numerables átomos puntuales de medida positiva: los de medida al menos $1/n$ son finitos para cada $n$, y todo peso positivo supera $1/n$ para algún $n$. Esto contradice que $D$ sea no numerable. El punto añadido no puede reparar la ausencia de masa en una de las vecindades singleton.
\end{solucion}
\section{Regularidad con sus hipótesis visibles}
La aproximación por compactos y funciones continuas funciona sobre conjuntos de medida finita y se extiende a conjuntos $\sigma$-finitos. El paso a todos los medibles requiere hipótesis adicionales; el ejemplo de masa infinita localmente nula muestra qué puede fallar. Mantener estas distinciones permite usar Lusin, densidad y localización sin convertir una propiedad local en una conclusión global injustificada.

\chapter{El teorema de representación de Riesz}\label{v2:cap16}\label{cap:v2-16}
\section{Funciones que se anulan en el infinito}\label{lectura:v2-33}
La dualidad de $L^p$ convierte funcionales en funciones integrables. Al sustituir las funciones medibles por funciones continuas, el objeto que representa una funcional es una medida. La topología decide qué medidas aparecen: son precisamente las medidas de Borel regulares y de variación finita. Desarrollamos la lectura 33 de \cite{valencia2} y completamos la demostración de suprayectividad que allí queda anunciada. A lo largo del capítulo, $X$ es localmente compacto y Hausdorff; no suponemos que sea metrizable ni $\sigma$-compacto.

\begin{definicion}
Una función $f:X\to\F$ \emph{se anula en el infinito} si, para cada $\eps>0$, existe un compacto $K\subset X$ tal que $|f(x)|<\eps$ para $x\notin K$. El espacio $C_0(X)$ consta de las funciones continuas con esta propiedad.\index{funcion@función!que se anula en el infinito}\index{espacio!C0@$C_0(X)$}
\end{definicion}
Para una función continua, la condición equivale a que $\{x:|f(x)|\ge\eps\}$ sea compacto para todo $\eps>0$: este conjunto es cerrado y está contenido en el compacto de la definición. Recíprocamente, dicho conjunto sirve como $K$. En un espacio compacto, $C_0(X)=C(X)$. En $\R$, una función continua pertenece a $C_0(\R)$ exactamente cuando tiende a cero al alejarse hacia ambos extremos. En un espacio discreto, la condición dice que sólo un número finito de valores puede tener módulo mayor o igual que un $\eps$ dado.

\begin{teorema}\label{v2:c0-banach}
Con la norma uniforme, $C_0(X)$ es un espacio de Banach y $C_c(X)$ es denso en él.
\end{teorema}
\begin{proof}
Primero, toda $f\in C_0(X)$ es acotada: fuera de algún compacto tiene módulo menor que uno, y sobre ese compacto es acotada por continuidad. Por tanto $C_0(X)\subset C_b(X)$, donde $C_b(X)$ denota las funciones continuas acotadas. Esta precisión evita usar una norma uniforme finita sobre todas las funciones continuas de un espacio no compacto.

Si $f,g\in C_0(X)$ y $a,b\in\F$, para $\eps>0$ elegimos compactos fuera de los cuales $|f|<\eps/[2(1+|a|)]$ y $|g|<\eps/[2(1+|b|)]$. Fuera de su unión, que es compacta porque toda cubierta abierta tiene subcubiertas finitas de cada uno de los dos compactos, se tiene $|af+bg|<\eps$. Así $C_0(X)$ es un subespacio.

El espacio $C_b(X)$ es completo: una sucesión de Cauchy uniforme tiene límites puntuales, la condición de Cauchy prueba convergencia uniforme hacia el límite, y un límite uniforme de funciones continuas es continuo. Si $f_n\in C_0(X)$ y $\norm{f_n-f}_\infty\to0$, fijamos $n$ con $\norm{f_n-f}_\infty<\eps/2$ y un compacto $K$ fuera del cual $|f_n|<\eps/2$. Entonces $|f|<\eps$ fuera de $K$. Esto muestra que $C_0(X)$ es cerrado en $C_b(X)$ y, por ello, completo.

Finalmente, sea $f\in C_0(X)$. Elegimos $K$ compacto con $|f|<\eps$ fuera de $K$. La proposición \ref{v2:topologia-lch} proporciona $h\in C_c(X)$ con $0\le h\le1$ y $h=1$ en $K$. Entonces $hf\in C_c(X)$, y
\[
|f-hf|=|f|(1-h)=0\quad\text{en }K,\qquad |f-hf|\le |f|<\eps\quad\text{fuera de }K.
\]
Luego $\norm{f-hf}_\infty\le\eps$. Además $C_c(X)\subset C_0(X)$, porque fuera del soporte la función vale cero.
\end{proof}

\section{Las medidas regulares de variación finita}
\begin{definicion}
Una medida real finita o compleja $\nu$ sobre $\Bcal(X)$ es \emph{regular} si para cada $A\in\Bcal(X)$ y cada $\eps>0$ existen un compacto $K$ y un abierto $U$, con $K\subset A\subset U$, tales que
\[
 |\nu(B)|<\eps\quad\text{para todo Borel }B\subset U\setminus K.
\]
Escribimos $M_r(X)$ para el espacio de estas medidas y $\norm{\nu}=|\nu|(X)$.\index{medida!regular compleja}
\end{definicion}
La condición debe controlar todos los subconjuntos de $U\setminus K$: pedir sólo que $|\nu(U\setminus K)|$ sea pequeño permitiría cancelaciones entre masas de signos o fases diferentes.

\begin{proposicion}\label{v2:medidas-regulares-banach}
Una medida finita real o compleja $\nu$ es regular si y sólo si su variación $|\nu|$ es una medida positiva regular. Con la norma de variación, $M_r(X)$ es un espacio de Banach.
\end{proposicion}
\begin{proof}
Si $|\nu|$ es regular, su finitud y la proposición \ref{v2:regular-interior-sigma} permiten elegir $K\subset A\subset U$ con $|\nu|(U\setminus K)<\eps$. Para $B\subset U\setminus K$ tenemos $|\nu(B)|\le |\nu|(B)<\eps$.

Recíprocamente, tomemos los $K,U$ de la definición para $\eps/4$. La estimación de variación del capítulo \ref{v2:cap12}, aplicada a la restricción a $U\setminus K$, da
\[
 |\nu|(U\setminus K)\le 4\sup_{B\subset U\setminus K}|\nu(B)|\le\eps.
\]
Como $|\nu|$ es finita, esto implica a la vez aproximación interior por $K$ y exterior por $U$ para cada Borel $A$. En particular cumple la definición de regularidad positiva del capítulo anterior.

La suma de dos medidas regulares es regular. En efecto, elegimos $K_1,K_2\subset A\subset U_1,U_2$ para cada medida y usamos $K=K_1\cup K_2$, $U=U_1\cap U_2$; el nuevo hueco está contenido en ambos huecos anteriores. La desigualdad triangular controla la suma. La estabilidad por escalares es inmediata, ajustando la tolerancia.

El espacio de todas las medidas finitas es completo por el teorema \ref{v2:medidas-banach}. Si $\nu_n\to\nu$ en variación y cada $\nu_n$ es regular, fijamos $n$ con $\norm{\nu-\nu_n}<\eps/2$. Para $A$ elegimos $K,U$ que hagan $|\nu_n(B)|<\eps/2$ en todo Borel $B\subset U\setminus K$. Entonces
\[
 |\nu(B)|\le |(\nu-\nu_n)(B)|+|\nu_n(B)|<\eps.
\]
Así $M_r(X)$ es un subespacio cerrado del espacio completo de medidas.
\end{proof}

\begin{lema}[Forma polar de una medida]\label{v2:polar-medida}
Para una medida finita real o compleja $\nu$ existe una función medible $h$ tal que $|h|=1$ $|\nu|$-casi en todas partes y
\[
 \nu(A)=\int_A h\,d|\nu|,\qquad
 \int f\,d\nu=\int fh\,d|\nu|\quad(f\in L^1(|\nu|)).
\]
\end{lema}
\begin{proof}
Si $|\nu|(A)=0$, también $\nu(A)=0$. Puesto que $|\nu|$ es finita, Radon--Nikodym proporciona $h\in L^1(|\nu|)$ con $\nu=h|\nu|$. La identidad de variación \ref{v2:densidad-variacion} muestra que $|\nu|=|h|\,|\nu|$. Por unicidad de densidades, $|h|=1$ casi en todas partes. La identidad de integrales vale primero para simples y después para toda función integrable por aproximación en $L^1(|\nu|)$. Si se desea $|h|=1$ en todos los puntos, se redefine $h$ como uno en el conjunto nulo excepcional.
\end{proof}

\begin{teorema}\label{v2:riesz-isometria}
Cada $\nu\in M_r(X)$ determina una funcional continua
\[
 T_\nu:C_0(X)\longrightarrow\F,\qquad T_\nu(f)=\int_X f\,d\nu,
\]
y $\norm{T_\nu}=|\nu|(X)$.
\end{teorema}
\begin{proof}
Toda $f\in C_0(X)$ es acotada y Borel medible, de modo que pertenece a $L^1(|\nu|)$ porque $|\nu|(X)<\infty$. Además,
\[
 |T_\nu(f)|\le\int|f|\,d|\nu|\le\norm f_\infty |\nu|(X).
\]
Falta mostrar que esta cota se alcanza arbitrariamente bien sobre la bola unidad. Escribimos $\nu=h|\nu|$ como en el lema anterior. En el caso complejo debemos aproximar $\overline h$, porque $\overline h\,h=1$; aproximar $h$ no eliminaría su fase.

Por la densidad de $C_c(X)$ en $L^1(|\nu|)$, hay $u\in C_c(X)$ con $\int|u-\overline h|\,d|\nu|$ tan pequeño como se quiera. Sustituimos $u$ por $g=P\circ u$, donde
\[
 P(z)=\frac{z}{\max\{1,|z|\}}.
\]
La aplicación $P$ es continua, fija el disco unidad y vale cero en cero; por eso $g\in C_c(X)$ y $\norm g_\infty\le1$. Para $|w|\le1$,
\[
 |P(z)-w|\le |P(z)-z|+|z-w|\le2|z-w|,
\]
pues $|P(z)-z|=(|z|-1)^+\le |z-w|$. Podemos, por tanto, asegurar $\int|g-\overline h|\,d|\nu|<\eps$. Entonces
\[
 |\nu|(X)=\int\overline h\,d\nu
 \le\left|\int g\,d\nu\right|+\eps
 \le\norm{T_\nu}+\eps.
\]
Al hacer $\eps\downarrow0$ obtenemos la igualdad de normas. En el caso real, la misma proyección es el recorte al intervalo $[-1,1]$.
\end{proof}

\section{El módulo de una funcional}
Para representar una funcional compleja, conviene dominarla por una funcional positiva. La siguiente construcción cumple exactamente ese propósito.
\begin{lema}\label{v2:modulo-funcional}
Sea $\varphi\in C_0(X)^*$. Para $f\in C_0(X)$, $f\ge0$, definimos
\[
 I(f)=\sup\{|\varphi(g)|:g\in C_0(X),\ |g|\le f\}.
\]
Entonces $I$ se extiende de manera única a una funcional lineal positiva y continua de $C_0(X)$, y $\norm I=\norm\varphi$.
\end{lema}
\begin{proof}
La familia del supremo contiene a cero y, si $|g|\le f$, se tiene $|\varphi(g)|\le\norm\varphi\,\norm f_\infty$. Así $I(f)$ es finito. Cambiar $g$ por $ag$, con $a>0$, prueba $I(af)=aI(f)$; para $a=0$ la igualdad es inmediata.

Demostremos la aditividad en el cono positivo. Dados $f_1,f_2\ge0$ y $\eps>0$, elegimos $g_j$ con $|g_j|\le f_j$ y $|\varphi(g_j)|>I(f_j)-\eps/2$. Multiplicamos cada $g_j$ por un escalar de módulo uno que haga $\varphi(g_j)$ real y no negativo. Esto no cambia su módulo. Como $|g_1+g_2|\le f_1+f_2$, resulta
\[
 I(f_1)+I(f_2)-\eps<\varphi(g_1+g_2)\le I(f_1+f_2).
\]
Para la desigualdad contraria, tomemos $g$ con $|g|\le f_1+f_2$ y pongamos
\[
 g_j(x)=\begin{cases}
 g(x)\dfrac{f_j(x)}{f_1(x)+f_2(x)},&f_1(x)+f_2(x)>0,\\[4pt]
 0,&f_1(x)+f_2(x)=0.
 \end{cases}
\]
Estas funciones son continuas donde el denominador es positivo. En sus ceros, la cota $|g_j|\le f_j$ demuestra continuidad porque $f_j$ también se anula allí. Además $g_j\in C_0(X)$, $|g_j|\le f_j$ y $g=g_1+g_2$. Por tanto $|\varphi(g)|\le I(f_1)+I(f_2)$. Tomando supremo y dejando $\eps\downarrow0$, concluimos $I(f_1+f_2)=I(f_1)+I(f_2)$.

Para una función real definimos $I(f)=I(f^+)-I(f^-)$. La definición es compatible con cualquier escritura $f=a-b$, $a,b\ge0$: de $a-b=c-d$ sigue $a+d=b+c$, y la aditividad positiva da $I(a)-I(b)=I(c)-I(d)$. Esta observación demuestra la aditividad y homogeneidad real de la extensión. Para funciones complejas ponemos $I(u+iv)=I(u)+iI(v)$; la linealidad compleja se verifica usando la linealidad real y $i(u+iv)=-v+iu$.

La positividad implica monotonicidad. Si $g$ es complejo, elegimos $\alpha$ de módulo uno tal que $\alpha I(g)=|I(g)|$. Como $\operatorname{Re}(\alpha g)\le|g|$,
\[
 |I(g)|=I(\operatorname{Re}(\alpha g))\le I(|g|)
 \le\norm\varphi\,\norm g_\infty.
\]
En el caso real basta elegir un signo. Así $\norm I\le\norm\varphi$. A la inversa, la propia definición da $|\varphi(g)|\le I(|g|)\le\norm I\,\norm g_\infty$, luego $\norm\varphi\le\norm I$. La extensión es única porque toda función real es diferencia de dos funciones positivas y toda función compleja tiene partes real e imaginaria.
\end{proof}

\section{Construir una medida a partir de una funcional positiva}
La parte central del argumento consiste en obtener una medida a partir de los valores de $I$ sobre funciones que caben dentro de un abierto. Presentamos todos los pasos: definir primero el tamaño de los abiertos, construir una medida exterior, comprobar la medibilidad de Borel y recuperar las integrales.

\begin{teorema}[Representación positiva sobre $C_c$]\label{v2:riesz-positivo-general}\index{Riesz--Markov, teorema de}
Sea $I:C_c(X)\to\F$ una funcional lineal positiva, es decir, $I(f)\ge0$ cuando $f$ es real y no negativa. Existe una única medida positiva $\mu$ sobre $\Bcal(X)$, finita sobre compactos, exteriormente regular en todos los Borel e interiormente regular en los abiertos, tal que
\[
 I(f)=\int_X f\,d\mu\qquad(f\in C_c(X)).
\]
Además, $\mu$ es interiormente regular en cada Borel de medida finita y, más generalmente, en cada Borel de medida $\sigma$-finita. No se exige que $I$ sea acotada para la norma uniforme sobre todo $C_c(X)$ ni que $\mu(X)$ sea finita.
\end{teorema}
\begin{proof}
Para un abierto $U$ definimos
\[
 m(U)=\sup\{I(f):f\in C_c(X),\ 0\le f\le1,\ \supp f\subset U\}.
\]
Se tiene $0\le m(U)\le\infty$, $m(\varnothing)=0$ y monotonía. La positividad proporciona las cotas locales que sustituirán una norma global: si $K$ es compacto y $h\in C_c(X)$ satisface $h\ge0$ y $h=1$ sobre $K$, entonces
\[
 0\le I(f)\le I(h)\qquad(0\le f\le1,\ \supp f\subset K).
\]
Si $V$ es abierto con cerradura compacta, elegimos una función de corte igual a uno sobre $\cl V$, y obtenemos $m(V)<\infty$. Todo compacto tiene una vecindad de esta clase por la proposición \ref{v2:topologia-lch}.

Si $U\subset\bigcup_{n\ge1}U_n$, cada compacto $\supp f\subset U$ queda cubierto por un número finito de los $U_n$. La partición de la unidad de la proposición \ref{v2:particion-compacta}, agrupando los términos correspondientes al mismo abierto si fuera necesario, escribe $f=\sum_j f\theta_j$, donde $0\le f\theta_j\le1$ y cada soporte está contenido en su abierto. De aquí
\[
 I(f)\le\sum_{n\ge1}m(U_n),\qquad m(U)\le\sum_{n\ge1}m(U_n).
\]
Si $U,V$ son abiertos disjuntos, funciones admisibles $f,g$ sobre ellos satisfacen $f+g\le1$. Al aproximar ambos supremos obtenemos $m(U\cup V)\ge m(U)+m(V)$; la desigualdad contraria procede de la subaditividad anterior.

Definamos ahora, para cualquier subconjunto $A\subset X$,
\[
 \mu^*(A)=\inf\{m(U):A\subset U,\ U\text{ abierto}\}.
\]
Ésta es una medida exterior. Sólo requiere explicación la subaditividad numerable: dados $A_n$ y $\eps>0$, si $\sum_n\mu^*(A_n)<\infty$, tomamos abiertos $U_n\supset A_n$ con $m(U_n)<\mu^*(A_n)+\eps2^{-n}$. Entonces
\[
 \mu^*\Bigl(\bigcup_n A_n\Bigr)\le m\Bigl(\bigcup_nU_n\Bigr)
 \le\sum_nm(U_n)\le\sum_n\mu^*(A_n)+\eps.
\]
Si la suma es infinita, la desigualdad es inmediata. Además $\mu^*(U)=m(U)$ para todo abierto $U$, por monotonía de $m$. Para cada compacto $K$, una vecindad $V$ de cerradura compacta da $\mu^*(K)\le m(V)<\infty$.

Necesitamos la identidad
\begin{equation}\label{v2:riesz-interior-abiertos}
 m(U)=\sup\{\mu^*(K):K\subset U,\ K\text{ compacto}\}.
\end{equation}
Una desigualdad es monotonía. Para la otra, si $f$ es admisible y $K=\supp f$, todo abierto $V\supset K$ admite a $f$ entre sus funciones de prueba. Luego $I(f)\le m(V)$ y, tomando ínfimo, $I(f)\le\mu^*(K)$. El supremo sobre $f$ prueba la identidad.

Verifiquemos que todo abierto $U$ satisface el criterio de Carathéodory. Si $\mu^*(E)=\infty$, la desigualdad
$\mu^*(E)\ge\mu^*(E\cap U)+\mu^*(E\setminus U)$ es automática, y la inversa siempre es subaditividad. Basta, pues, considerar $E$ de medida exterior finita. Al tomar el ínfimo que define $\mu^*(E)$, podemos restringirnos a abiertos $V\supset E$ con $m(V)<\infty$. Fijemos uno de ellos y $\eps>0$. Por \eqref{v2:riesz-interior-abiertos} elegimos $K\subset V\cap U$ compacto con
\[
 \mu^*(K)\ge m(V\cap U)-\eps\ge\mu^*(E\cap U)-\eps.
\]
Si $V\cap U$ es vacío, se toma $K=\varnothing$. Existe un abierto $W$ con $K\subset W\subset\cl W\subset V\cap U$ y $\cl W$ compacto. Los abiertos $W$ y $V\setminus\cl W$ son disjuntos y están contenidos en $V$. Por ello,
\begin{align*}
 m(V)&\ge m(W)+m(V\setminus\cl W)\\
 &\ge\mu^*(K)+\mu^*(E\setminus U)\\
 &\ge\mu^*(E\cap U)+\mu^*(E\setminus U)-\eps.
\end{align*}
Hemos usado $E\setminus U\subset V\setminus\cl W$. Tomando ínfimo en $V$ y después $\eps\downarrow0$ obtenemos la desigualdad no trivial del criterio; la otra es subaditividad. Así todos los abiertos, y por consiguiente todos los Borel, son $\mu^*$-medibles.

Sea $\mu$ la restricción a los Borel. Es una medida finita sobre compactos, exteriormente regular por construcción e interiormente regular en abiertos por \eqref{v2:riesz-interior-abiertos}. La proposición \ref{v2:regular-interior-sigma} prueba regularidad interior en los Borel de medida $\sigma$-finita, incluidos los de medida finita. Esta conclusión no afirma regularidad interior de todos los Borel cuando su medida es infinita y no $\sigma$-finita.

Probemos $I(f)=\int f\,d\mu$ para $f\in C_c(X)$ no negativa. Sea $K=\supp f$; su medida es finita y todas las funciones que siguen se anulan fuera de $K$. Fijemos $\delta>0$ y un entero $n$ con $n\delta\ge\norm f_\infty$. Escribamos $U_j=\{f>j\delta\}$. Para $1\le j\le n$, elegimos funciones admisibles $h_j$ en $U_j$ cuyos valores $I(h_j)$ aproximen $m(U_j)=\mu(U_j)$. Punto a punto,
\[
 \delta\sum_{j=1}^n h_j\le f,
\]
porque sólo pueden contribuir los índices con $j\delta<f(x)$ y cada sumando es a lo sumo uno. La positividad y la aproximación arbitraria de los supremos dan
\[
 I(f)\ge\delta\sum_{j=1}^n\mu(U_j).
\]
La función simple $\delta\sum_{j=1}^n\ind{U_j}$ queda entre $f-\delta$ y $f$, por lo que el lado derecho es al menos $\int f\,d\mu-\delta\mu(K)$. Al hacer $\delta\downarrow0$ obtenemos $I(f)\ge\int f\,d\mu$.

Para la desigualdad contraria, consideremos las funciones continuas
\[
 \psi_j=\min\{1,\max\{(f-j\delta)/\delta,0\}\},\qquad 0\le j\le n-1.
\]
Tienen soporte compacto y $f=\delta\sum_{j=0}^{n-1}\psi_j$. Elegimos $h\in C_c(X)$ no negativa e igual a uno sobre $K$. La primera satisface $I(\psi_0)\le I(h)$ por positividad. Para $j\ge1$, su soporte está contenido en el compacto $\{f\ge j\delta\}\subset U_{j-1}$, así que $I(\psi_j)\le\mu(U_{j-1})$. En consecuencia,
\begin{align*}
 I(f)&\le\delta I(h)+\delta\sum_{k=0}^{n-2}\mu(U_k)\\
 &\le\int f\,d\mu+\delta\bigl(I(h)+\mu(K)\bigr).
\end{align*}
La última desigualdad se sigue de $\delta\sum_{k=0}^{n-2}\ind{U_k}\le f+\delta\ind K$; una suma vacía se interpreta como cero. Haciendo otra vez $\delta\downarrow0$ obtenemos la igualdad.

Las partes positiva y negativa, y luego las partes real e imaginaria, extienden la identidad a $C_c(X)$. Estas integrales son finitas porque las funciones son acotadas y tienen soporte de medida finita.

Finalmente, si $\lambda$ es otra medida positiva finita sobre compactos y con las regularidades indicadas que representa a $I$, entonces para todo abierto $U$ se tiene
\[
 \lambda(U)=\sup\left\{\int f\,d\lambda:f\in C_c(X),\ 0\le f\le1,\ \supp f\subset U\right\}.
\]
La desigualdad $\ge$ es inmediata; la otra se obtiene aproximando $U$ por compactos y usando una función de corte que sea uno sobre cada compacto. Por tanto $\lambda(U)=m(U)=\mu(U)$ para todo abierto. La regularidad exterior determina los valores sobre todos los Borel y prueba $\lambda=\mu$.
\end{proof}

\begin{corolario}[Representación positiva sobre $C_0$]\label{v2:riesz-positivo}
Si $I$ es una funcional lineal positiva y continua sobre $C_0(X)$, la medida del teorema anterior es finita y satisface
\[
 I(f)=\int f\,d\mu\qquad(f\in C_0(X)).
\]
Es la única medida positiva finita regular con esta propiedad.
\end{corolario}
\begin{proof}
Restringimos $I$ a $C_c(X)$. En la construcción anterior, cada función admisible para $m(X)$ tiene norma uniforme a lo sumo uno; luego $\mu(X)=m(X)\le\norm I<\infty$. La representación vale sobre $C_c(X)$ por el teorema. Si $f_n\in C_c(X)$ converge uniformemente a $f\in C_0(X)$, entonces $I(f_n)\to I(f)$ por continuidad y
$|\int(f_n-f)\,d\mu|\le\norm{f_n-f}_\infty\mu(X)\to0$. La densidad del teorema \ref{v2:c0-banach} prueba la identidad en todo $C_0(X)$. La unicidad ya vale en $C_c(X)$. Además, como todos los Borel tienen medida finita, la regularidad interior se extiende a todos ellos.
\end{proof}

\section{La representación completa del dual}
\begin{teorema}[Riesz--Markov para $C_0(X)$]\label{v2:riesz-c0}
Para un espacio localmente compacto Hausdorff $X$, la aplicación
\[
 M_r(X)\longrightarrow C_0(X)^*,\qquad
 \nu\longmapsto T_\nu,\quad T_\nu(f)=\int_X f\,d\nu,
\]
es un isomorfismo lineal isométrico. En particular, toda funcional continua admite una única representación mediante una medida de Borel regular de variación finita.\index{teorema!de representacion de Riesz@de representación de Riesz}\index{dual!de C0@$C_0(X)$}
\end{teorema}
\begin{proof}
La linealidad y la isometría ya están demostradas. Sea $\varphi\in C_0(X)^*$. El lema \ref{v2:modulo-funcional} proporciona una funcional positiva $I$ tal que
\[
 |\varphi(g)|\le I(|g|)\qquad(g\in C_0(X)).
\]
Por el teorema anterior, $I$ está representada por una medida positiva finita regular $\mu$. En consecuencia,
\[
 |\varphi(g)|\le\int|g|\,d\mu\qquad(g\in C_c(X)).
\]
Esta desigualdad tiene dos consecuencias. Primero, $\varphi(g)$ sólo depende de la clase de $g$ en $L^1(\mu)$: si $g=0$ casi en todas partes, la cota da $\varphi(g)=0$. Segundo, como $C_c(X)$ es denso en $L^1(\mu)$ por el teorema \ref{v2:cc-denso-lp}, la funcional se extiende de manera única a una funcional continua de $L^1(\mu)$ de norma a lo sumo uno. La extensión se obtiene tomando límites de $\varphi(g_n)$ cuando $g_n\to f$ en $L^1$; la desigualdad anterior prueba tanto la existencia como la independencia de la aproximación.

La medida $\mu$ es finita, por tanto localizable. La representación del dual de $L^1$, teorema \ref{v2:dual-l1-general}, proporciona $h\in L^\infty(\mu)$, $\norm h_\infty\le1$, con
\[
 \varphi(g)=\int gh\,d\mu\qquad(g\in C_c(X)).
\]
Definamos $\nu=h\mu$. Es una medida finita real o compleja y $|\nu|=|h|\mu$. Es regular: para $K\subset A\subset U$ elegidos mediante la regularidad de $\mu$,
\[
 |\nu|(U\setminus K)=\int_{U\setminus K}|h|\,d\mu
 \le\mu(U\setminus K),
\]
y el último término puede hacerse arbitrariamente pequeño. Por densidad de $C_c(X)$ en $C_0(X)$, la identidad $\varphi(g)=T_\nu(g)$ se extiende a todo $C_0(X)$.

Si dos medidas regulares representan la misma funcional, su diferencia representa la funcional cero. La igualdad de normas del teorema \ref{v2:riesz-isometria} implica que la variación de esa diferencia es cero, y por tanto las medidas son iguales.
\end{proof}
\begin{observacion}
El teorema \ref{v2:riesz-positivo-general} permite funcionales positivas sobre $C_c(X)$ sin acotación uniforme global y medidas de masa infinita; el corolario \ref{v2:riesz-positivo} identifica el caso acotado sobre $C_0(X)$ con medidas finitas. En el primero, sobre un espacio localmente compacto Hausdorff general, no se debe afirmar regularidad interior de todos los Borel sin hipótesis adicionales. Las distinciones del capítulo \ref{v2:cap15} siguen siendo necesarias.
\end{observacion}

\section{Ejercicio resuelto: una obstrucción a la reflexividad}
\begin{ejercicio}[Lectura 33, ejercicio 1]\label{ej:v2-l33-1}
Suponga que existe una medida de Borel regular real o compleja $\mu$ tal que $\mu(\{x\})=0$ para todo $x\in X$ y $0<|\mu|(X)<\infty$. Demuestre que $C_0(X)$ no es reflexivo. Considere la funcional sobre medidas dada por $\Phi(\nu)=\sum_{x\in X}\nu(\{x\})$.
\end{ejercicio}
\begin{solucion}
Primero hay que dar sentido a la suma, que puede estar indexada por un conjunto no numerable. Para cualquier conjunto finito $F\subset X$,
\[
 \sum_{x\in F}|\nu(\{x\})|=\sum_{x\in F}|\nu|(\{x\})\le|\nu|(X).
\]
En particular, para cada $n$ sólo hay un número finito de puntos con $|\nu(\{x\})|\ge1/n$; de otro modo las sumas finitas superarían cualquier cota. La unión de estos conjuntos finitos contiene todos los puntos con masa no nula, así que éstos son a lo sumo numerables. La serie de sus masas converge absolutamente y define $\Phi(\nu)$, independientemente del orden. La desigualdad anterior prueba
\[
 |\Phi(\nu)|\le\sum_x|\nu(\{x\})|\le\norm\nu.
\]
La linealidad se obtiene sumando sobre la unión numerable de los soportes atómicos de las medidas involucradas. Así $\Phi\in M_r(X)^*$.

Supongamos que $C_0(X)$ fuera reflexivo. Mediante la identificación isométrica $C_0(X)^*=M_r(X)$, la funcional $\Phi$ pertenecería al bidual y estaría representada por evaluación en alguna $f\in C_0(X)$:
\[
 \Phi(\nu)=\int f\,d\nu\qquad(\nu\in M_r(X)).
\]
Cada medida de Dirac $\delta_x$ es regular, y $\Phi(\delta_x)=1$. Sustituyéndola en la identidad obtenemos $f(x)=1$ para todo $x\in X$.

Ahora usemos la medida positiva $\lambda=|\mu|$, que es regular y no nula. En un singleton, la definición de variación da $\lambda(\{x\})=|\mu(\{x\})|=0$. Por ello $\Phi(\lambda)=0$. Sin embargo, la representación y $f\equiv1$ darían
\[
 0=\Phi(\lambda)=\int f\,d\lambda=\lambda(X)>0,
\]
una contradicción. Usar $|\mu|$ es esencial para esta última contradicción: la masa total de una medida compleja no nula puede ser cero por cancelación.
\end{solucion}

El recorrido del volumen termina con dos representaciones complementarias. En $L^p$, la norma y la medida de partida determinan densidades que representan funcionales. En $C_0(X)$, la continuidad y el comportamiento en el infinito permiten reconstruir la medida misma. Las aproximaciones por simples, los teoremas de convergencia, la variación, Radon--Nikodym y la regularidad intervienen aquí con una función precisa, y las hipótesis que las conectan quedan visibles en la demostración.

